% ════════════════════════════════════════════════════════════════════════════════════════
%  A Hot-Started Nullspace Active-Set Method for Regularized Column Generation
%  with Equality and Inequality Linking Constraints
%
%  Compile with:  pdflatex main  (twice, for the references)
%  Figures in the results chapter are drawn by pgfplots from the tables in data/, which are
%  produced by scripts/generate_data.jl and scripts/make_plot_data.jl.
% ════════════════════════════════════════════════════════════════════════════════════════

% Document properties
\documentclass[11pt, openany]{book}
\usepackage[a4paper, margin=2.6cm]{geometry}
\usepackage[english]{babel}
\usepackage[T1]{fontenc}
\usepackage[utf8]{inputenc}
\usepackage{lmodern}
\title{A Hot-Started Nullspace Active-Set Method for Regularized Column Generation}
\author{Kyrian Bäumer}
\date{\today}


% Mathematics packages
\usepackage{amsmath}
\usepackage{mathtools}
\usepackage{amssymb}
\usepackage{amsthm}
\usepackage{bm}


% Graphics
\usepackage{graphicx}
\usepackage{subcaption}
\usepackage{booktabs}
\usepackage{multirow}
\usepackage{array}
\usepackage{xcolor}


% TikZ plotting
\usepackage{tikz}
\usepackage{pgfplots}
\pgfplotsset{compat=1.18}
\usepackage{pgfplotstable}
\usepgfplotslibrary{groupplots}
\usetikzlibrary{arrows.meta, patterns, positioning, calc, fit, decorations.pathreplacing,
                shapes.geometric, matrix, backgrounds}


% Colours used in figures (colour-blind safe)
\definecolor{cblue}{HTML}{0072B2}
\definecolor{corange}{HTML}{E69F00}
\definecolor{cgreen}{HTML}{009E73}
\definecolor{cred}{HTML}{D55E00}
\definecolor{cpurple}{HTML}{CC79A7}
\definecolor{csky}{HTML}{56B4E9}
\definecolor{cgrey}{HTML}{6E6E6E}


% Hyperreferences
\usepackage{hyperref}
\hypersetup{colorlinks=true, linkcolor=cblue, citecolor=cgreen, urlcolor=cblue}
\usepackage[nameinlink, capitalise]{cleveref}
%\crefname{prop}{Proposition}{Propositions}
%\crefname{cor}{Corollary}{Corollaries}
%\crefname{definition}{Definition}{Definitions}
%\crefname{remark}{Remark}{Remarks}
%\crefname{note}{Note}{Notes}
%\crefname{assumption}{Assumption}{Assumptions}
\crefname{algo}{Algorithm}{Algorithms}


% Theorem environment (shared numbering; aliascnt lets cleveref name each type)
\usepackage{aliascnt}
\theoremstyle{definition}
\newtheorem{theorem}{Theorem}[section]
\newaliascnt{cor}{theorem}\newtheorem{cor}[cor]{Corollary}\aliascntresetthe{cor}
\newaliascnt{prop}{theorem}\newtheorem{prop}[prop]{Proposition}\aliascntresetthe{prop}
\newaliascnt{lemma}{theorem}\newtheorem{lemma}[lemma]{Lemma}\aliascntresetthe{lemma}
\newaliascnt{example}{theorem}\newtheorem{example}[example]{Example}\aliascntresetthe{example}
\newaliascnt{definition}{theorem}\newtheorem{definition}[definition]{Definition}\aliascntresetthe{definition}
\newaliascnt{remark}{theorem}\newtheorem{remark}[remark]{Remark}\aliascntresetthe{remark}
\newaliascnt{note}{theorem}\newtheorem{note}[note]{Note}\aliascntresetthe{note}
\newaliascnt{assumption}{theorem}\newtheorem{assumption}[assumption]{Assumption}\aliascntresetthe{assumption}
\crefname{theorem}{Theorem}{Theorems}
\crefname{cor}{Corollary}{Corollaries}
\crefname{prop}{Proposition}{Propositions}
%\crefname{lemma}{Lemma}{Lemmas}
%\crefname{example}{Example}{Examples}
\crefname{definition}{Definition}{Definitions}
\crefname{remark}{Remark}{Remarks}
\crefname{note}{Note}{Notes}
\crefname{assumption}{Assumption}{Assumptions}


% Enumeration (Roman)
\usepackage{enumitem}
\setlist[enumerate,1]{label=(\roman*)}


% Pseudocode and boxes
\usepackage{tcolorbox}
\usepackage{algorithm}
\usepackage{algpseudocode}
\newcounter{algo}[chapter]
\renewcommand{\thealgo}{\thechapter.\arabic{algo}}
% \begin{algobox}{Title}{label} ... \end{algobox}
\newenvironment{algobox}[2]{%
  \par\medskip\refstepcounter{algo}\label{#2}%
  \begin{tcolorbox}[colback=gray!10!white, colframe=gray!80!black, sharp corners,
                    boxrule=0.8pt, title={Algorithm \thealgo: #1}]
  \small\begin{algorithmic}[1]}{%
  \end{algorithmic}\end{tcolorbox}\medskip}


% Table of Contents
\usepackage{tocloft}
\setlength{\cftbeforesecskip}{6pt}
\setlength{\cftbeforesubsecskip}{3pt}


% Miscellaneous
\usepackage{url}
\setlength{\parindent}{0pt}
\setlength{\parskip}{3pt}
\usepackage{microtype}
\usepackage{needspace}
\usepackage[section]{placeins}   % keep floats inside their section
\raggedbottom
% expandable input for generated table rows (LaTeX's \input breaks \bottomrule in a tabular)
\makeatletter
\newcommand{\tabinput}[1]{\@@input #1 }
\setlength{\@fptop}{0pt}          % float pages start at the top, not vertically centred
\setlength{\@fpsep}{24pt plus 4pt}
\setlength{\@fpbot}{0pt plus 1fil}
\makeatother
\pgfplotsset{legend cell align=left}


% ── notation ──────────────────────────────────────────────────────────────────────────────
\newcommand{\R}{\mathbb{R}}
\newcommand{\one}{\mathbf{1}}
\newcommand{\Sset}{\mathcal{S}}
\newcommand{\Wx}{\mathcal{W}_x}
\newcommand{\Nx}{\mathcal{N}_x}
\newcommand{\Ws}{\mathcal{W}_s}
\newcommand{\Ns}{\mathcal{N}_s}
\newcommand{\Wset}{\mathcal{W}}
\newcommand{\Fset}{\mathcal{F}}
\newcommand{\rep}{\operatorname{rep}}
\newcommand{\Heff}{H_{\mathrm{eff}}}
\newcommand{\Aeff}{A_{\mathrm{eff}}}
\newcommand{\pieff}{\pi_{\mathrm{eff}}}
\newcommand{\Meff}{M_{\mathrm{eff}}}
\newcommand{\Abar}{\bar{A}}
\newcommand{\Cbar}{\bar{C}}
\newcommand{\cbar}{\hat{c}}       % augmented cost [c; 0] (\bar c_j is the reduced cost)
\newcommand{\Hbar}{\bar{H}}
\newcommand{\ctil}{\tilde{c}}
\newcommand{\pihat}{\hat{\pi}}
\newcommand{\nmax}{n_{\max}}
\newcommand{\nnz}{\operatorname{nnz}}
\newcommand{\reach}{\operatorname{reach}}
\newcommand{\diag}{\operatorname{diag}}
\newcommand{\Null}{\operatorname{Null}}
\newcommand{\rank}{\operatorname{rank}}
\newcommand{\spn}{\operatorname{span}}
\newcommand{\code}[1]{\texttt{#1}}



\begin{document}
\let\cleardoublepage\clearpage   % no blank pages before chapters (one-sided reading)

\frontmatter

\begin{titlepage}
\begin{center}

\vspace*{\stretch{2}}

\Large{\textsc{Applied Mathematics}}

\vspace*{\stretch{4}}

\textbf{\huge{\em A Hot-Started Nullspace Active-Set Method}}

\bigskip\medskip

\textbf{\huge{\em for Regularized Column Generation}}

\bigskip\bigskip

\Large{with Equality and Inequality Linking Constraints}

\vspace*{\stretch{5}}

\large{Author: \emph{Kyrian Bäumer}}

\vspace*{\stretch{3}}

\textsc{\today}

\end{center}
\end{titlepage}


\chapter*{Abstract}
\addcontentsline{toc}{chapter}{Abstract}

Column generation solves linear programs with a very large number of variables by working
on a restricted master problem that is enlarged one column at a time. Its classical form is
known for unstable dual iterates and slow convergence. A quadratic proximal term in the dual
space removes most of this instability, but it turns every master problem into a convex
quadratic program, which is considerably more expensive to solve than a linear program.

\medskip

This work develops a \textbf{hot-started nullspace active-set method} for these quadratic
master problems, for multi-block problems whose linking constraints may be equalities,
``$\ge$'' inequalities or ``$\le$'' inequalities. Surplus variables bring all linking rows
into one equality form. The convexity constraints of the blocks then admit a
reordering of the variables into five sets, which yields an explicit and sparse nullspace
basis. The reduced KKT system decouples: the surplus part of the search direction is given
by an explicit formula, and only an \emph{effective reduced Hessian} on the column weights
has to be factorized. This matrix is the reduced Hessian of the equality rows together with
the inequality rows whose surplus is fixed at zero. Every change of the working set becomes
a low-rank modification of an $LDL^\top$ factorization of this matrix, which is updated in
place and carried over from one column generation iteration to the next. The Lagrange
multipliers of the method coincide with the reduced costs of column generation and with the
dual sign conditions of the inequality rows, and singular reduced Hessians are resolved by
an exact zero-curvature step.

\medskip

Two factorizations are developed. The first is a sparse factorization of fixed dimension,
updated with the sparse direct solver CHOLMOD, in which periodic refactorization controls the
structural zeros left by the updates. The second exploits that the effective reduced Hessian
is small and dense: it is factorized densely in the order in which its indices became active,
so that a new index is appended and a removed index is deleted by a rank-one update. This
factor never contains structural zeros, needs no periodic refactorization, and its pivots
detect singular reduced Hessians exactly.

\medskip

On generated airline crew scheduling instances with up to 960 linking rows the method
reproduces the optimum of the full linear program, and it solves the regularized master
problems about three orders of magnitude faster than a general-purpose quadratic programming
solver. The dense factorization is two to four times faster than the sparse one, with a
speedup that grows with the number of linking rows, and needs an order of magnitude less
memory. The complete column generation run is much faster than classical column generation
with a linear master problem, and, on all but the smallest instance, faster than solving the
full linear program directly.


\cleardoublepage
\tableofcontents


\mainmatter

\chapter{Introduction}
\label{chap:introduction}

\vspace*{\fill}

Many large-scale planning problems, such as crew scheduling, vehicle routing and cutting stock,
lead to linear programs with a block structure: a number of independent subproblems, one per
crew member, vehicle or machine, are coupled by a comparatively small set of
\textbf{linking constraints}. In crew scheduling, for instance, every pilot follows a
schedule of their own, while the linking constraints require that each flight is covered by
the pilots together.

\medskip

\textbf{Column generation} is the standard method for such problems. The problem is first
rewritten in terms of complete solutions of the subproblems (a schedule, a route, a cutting
pattern), so that every variable, or \emph{column}, represents one such solution. There are
far too many columns to write down, but only a few of them are needed in an optimal
solution. Column generation therefore works with a \textbf{restricted master problem} over a
small set of columns and asks the subproblems, in the \textbf{pricing problem}, for a new
column that can improve the current solution.

\medskip

In its classical form the restricted master problem is a linear program. Its dual solutions,
which drive the pricing problem, are known to jump from one extreme point to another, which
makes the method unstable and slow in its final phase \cite{lubbecke2005selected}. A
well-established remedy is \textbf{proximal stabilization}
\cite{dumerle1999stabilized, briant2008comparison}: the dual problem is augmented by a
quadratic term that keeps the dual iterate close to a \emph{stability center}. The price of this regularization is
that every restricted master problem becomes a \textbf{convex quadratic program}.

\medskip

This work is about solving this sequence of quadratic programs efficiently, for linking
constraints that may be equalities as well as inequalities.

\vspace*{\fill}

\newpage


\section{The Computational Problem}

Two consecutive restricted master problems in column generation are almost identical. The
new problem contains one more column, and with stabilization the stability center may have
moved. A general-purpose quadratic programming solver does not exploit this: it rebuilds and
refactorizes its internal linear algebra for every master problem, and its cost per master
problem grows with the number of generated columns.

\bigskip

\textbf{Active-set methods} are well suited for sequences of closely related quadratic
programs. They move between faces of the feasible region by adding or removing one
constraint at a time, and the linear system they solve at each step changes only by a
low-rank term. In the \textbf{nullspace} variant, this system is the reduced Hessian on the
nullspace of the active constraints. Its size equals the number of degrees of freedom of the
current face, not the number of variables.

\bigskip

The question addressed in this work is how far this structure can be exploited for the
regularized master problem of column generation, in particular:
\begin{enumerate}
    \item What is the structure of the active constraints of the regularized master problem,
          and does it allow an explicit nullspace basis?
    \item How small is the linear system that has to be factorized, and how does it change
          when the working set changes?
    \item Can the factorization be carried over from one column generation iteration to the
          next, so that a new column or a new stability center costs (almost) nothing?
    \item What changes when the linking constraints are inequalities, which add sign
          conditions on the dual variables and destroy part of the structure?
    \item How does the resulting method compare in practice with general-purpose solvers and
          with classical column generation?
\end{enumerate}


\section{Contributions}

The main results of this work are the following.

\begin{itemize}
    \item \textbf{A unified formulation for mixed linking constraints.} Equality rows and
          inequality rows of both directions are written as equalities with surplus
          variables. The regularized master problem then has a single form for all cases
          (\cref{chap:formulation,chap:regularization}).

    \item \textbf{A five-set column ordering and an explicit nullspace basis.} The convexity
          constraints of the blocks, together with the surplus variables, admit an ordering
          of the variables into representatives, working-set columns, nullspace columns,
          active surpluses and free surpluses. The nullspace basis consists of vectors
          $e_j - e_{\rep(j)}$ and unit vectors, and it is never stored
          (\cref{chap:qpas,chap:nullspace}).

    \item \textbf{Decoupling and the effective reduced Hessian.} The reduced KKT system splits
          into an explicit formula for the surplus directions and a system with the
          \emph{effective reduced Hessian} $\Heff = (\Aeff Z_x)^\top(\Aeff Z_x)$. Here $\Aeff$
          holds the equality rows together with the inequality rows whose surplus is fixed at
          zero. Its dimension never exceeds the number of linking constraints
          (\cref{chap:nullspace}).

    \item \textbf{Multipliers are reduced costs.} The Lagrange multipliers of the active-set
          method are exactly the reduced costs of column generation and the dual sign
          conditions of the inequality rows. The optimality test of the master problem and
          the pricing test are therefore the same test, applied inside and outside the
          restricted master problem (\cref{chap:nullspace}).

    \item \textbf{Updates and hot start.} Every change of the working set is a row
          addition, a row deletion or a rank-one modification of a fixed-size $LDL^\top$
          factorization. A new column or a new stability center leaves the factorization
          unchanged (\cref{chap:updates}).

    \item \textbf{Exact treatment of singular reduced Hessians.} A dependent nullspace
          direction is resolved by a zero-curvature step that keeps every working-set
          variable except the released one exactly at its bound (\cref{chap:singular}).

    \item \textbf{A computational study} on generated crew scheduling instances with equality
          and inequality coverage rows: correctness against the full linear program, timing,
          the two phases of convergence, and the effect of refactorization on the fill of the
          factor (\cref{chap:results}).
\end{itemize}


\section{Outline}

\Cref{chap:formulation} introduces the multi-block linear program with mixed linking
constraints, its Dantzig--Wolfe reformulation, column generation and the surplus
formulation. \Cref{chap:regularization} adds the proximal regularization, eliminates the
auxiliary variable and derives the relation between the gradient of the quadratic master
problem and the reduced costs. \Cref{chap:qpas} sets up the active-set method for the
regularized master problem and establishes the five-set reordering of the constraint matrix.
\Cref{chap:nullspace} constructs the nullspace basis, solves the reduced KKT system with the
effective reduced Hessian and recovers the Lagrange multipliers. \Cref{chap:updates}
describes the updates of the factorization when the working set changes, the fixed-length
representation, periodic refactorization and the hot start between column generation
iterations. \Cref{chap:singular} treats singular reduced Hessians. \Cref{chap:results}
presents the computational results, and \cref{chap:conclusion} concludes.

\bigskip

Throughout, the reader is assumed to be familiar with linear programming duality, the basic
idea of column generation and the basic idea of a primal active-set method for quadratic
programming, as found in standard texts such as \cite{nocedal2006numerical, gill1981practical,
desaulniers2005column}. Everything beyond this is derived in the text.

\chapter{Problem Formulation}
\label{chap:formulation}

\vspace*{\fill}

This chapter introduces the class of problems considered throughout this work: linear
programs with a multi-block structure, whose blocks are coupled by linking constraints that
may be equalities or inequalities.

\medskip

We first state the problem in its original (compact) form and then apply the
\textbf{Dantzig--Wolfe reformulation}, which replaces the variables of each block by convex
combinations of the extreme points of the block. This leads to the \textbf{master problem},
a linear program with a very large number of columns, which is solved by
\textbf{column generation}. We define the restricted master problem, the reduced cost, the
pricing problem and the Lagrangian bound, and collect them in the classical column
generation algorithm.

\medskip

Finally, we introduce \textbf{surplus variables} for the inequality rows. They bring every
linking constraint into equality form and lead to a single, \emph{unified} formulation of the
master problem in which equality and inequality rows are treated alike. This formulation is
the starting point for everything that follows. It also reveals the one structural property
that the inequality rows destroy: the surplus variables belong to no block.

\vspace*{\fill}

\newpage


\section{Multi-Block Linear Programs with Mixed Linking Constraints}
\label{sec:multiblock}

We consider a linear program whose variables are partitioned into $K$ blocks
$x_1, \dots, x_K$, with $x_k \in \R^{n_k}$. Each block has its own constraints, and the blocks
interact only through a set of linking constraints.

\begin{definition}[Multi-block linear program] \label{def:multiblock}
A \textbf{multi-block linear program with mixed linking constraints} is a problem of the
form
\begin{subequations} \label{eq:compact}
\begin{align}
\min_{x_1,\dots,x_K} \quad & \sum_{k=1}^{K} c_k^\top x_k \\
\text{s.t.} \quad
& \sum_{k=1}^{K} A_{E,k}\, x_k = b_E, \label{eq:compact-eq} \\
& \sum_{k=1}^{K} A_{I,k}\, x_k \ge b_I, \label{eq:compact-ineq} \\
& x_k \in P_k, \qquad k = 1,\dots,K, \label{eq:compact-block}
\end{align}
\end{subequations}
where $A_{E,k} \in \R^{m_E\times n_k}$ and $A_{I,k}\in\R^{m_I\times n_k}$ are the
coefficients of block $k$ in the $m_E$ \textbf{equality linking rows} and the $m_I$
\textbf{inequality linking rows}, and
\[
P_k = \{\, x_k \in \R^{n_k} \;:\; F_k x_k = d_k,\; l_k \le x_k \le u_k \,\}
\]
is the \textbf{block polytope} of block $k$. We write $m = m_E + m_I$ for the total number of
linking rows. Constraints~\eqref{eq:compact-eq}--\eqref{eq:compact-ineq} are the
\textbf{linking constraints}, and constraints~\eqref{eq:compact-block} are the
\textbf{block constraints}.
\end{definition}

\begin{assumption} \label{ass:bounded}
Every block polytope $P_k$ is nonempty and bounded, and problem~\eqref{eq:compact} is
feasible.
\end{assumption}

Under \cref{ass:bounded} the feasible region of~\eqref{eq:compact} is a nonempty polytope,
so the problem has an optimal solution. Its optimal value is denoted by $z^\star$.

\bigskip

The constraint matrix of~\eqref{eq:compact} has the \textbf{block-angular} structure shown
in \cref{fig:block-angular}. The block constraints form a block-diagonal matrix, and only
the linking rows couple different blocks. If the linking rows were absent, the problem would
decompose into $K$ independent problems over $P_1, \dots, P_K$. Column generation exploits
exactly this observation.

\begin{figure}[ht]
\centering
\begin{tikzpicture}[x=0.55cm, y=0.55cm]
    % linking rows
    \fill[cblue!25] (0,10) rectangle (16,12);
    \fill[corange!35] (0,8.6) rectangle (16,10);
    \foreach \k in {0,1,2,3} {
        \draw[cblue!70!black] (4*\k,10) rectangle (4*\k+4,12);
        \draw[corange!80!black] (4*\k,8.6) rectangle (4*\k+4,10);
        \node at (4*\k+2,11) {\small $A_{E,\the\numexpr\k+1\relax}$};
        \node at (4*\k+2,9.3) {\small $A_{I,\the\numexpr\k+1\relax}$};
    }
    % block-diagonal part
    \foreach \k in {0,1,2,3} {
        \fill[cgreen!25] (4*\k, 8.6-2*\k-2) rectangle (4*\k+4, 8.6-2*\k);
        \draw[cgreen!70!black] (4*\k, 8.6-2*\k-2) rectangle (4*\k+4, 8.6-2*\k);
        \node at (4*\k+2, 8.6-2*\k-1) {\small $F_{\the\numexpr\k+1\relax}$};
    }
    \draw[thick] (0,0.6) rectangle (16,12);
    % labels
    \draw[decorate, decoration={brace, amplitude=4pt}] (16.3,12) -- (16.3,10)
        node[midway, right=5pt] {\small equality rows, $=b_E$};
    \draw[decorate, decoration={brace, amplitude=4pt}] (16.3,10) -- (16.3,8.6)
        node[midway, right=5pt] {\small inequality rows, $\ge b_I$};
    \draw[decorate, decoration={brace, amplitude=4pt}] (16.3,8.6) -- (16.3,0.6)
        node[midway, right=5pt] {\small block constraints, $x_k\in P_k$};
    \foreach \k in {0,1,2,3}
        \node at (4*\k+2, 12.6) {\small $x_{\the\numexpr\k+1\relax}$};
\end{tikzpicture}
\caption{Block-angular structure of the multi-block linear program~\eqref{eq:compact} for
$K = 4$ blocks. The linking rows (top) couple all blocks; the block constraints (bottom) are
block-diagonal.}
\label{fig:block-angular}
\end{figure}

\begin{remark}[Inequalities of the other direction] \label{rem:le-rows}
A linking row of the form $\sum_k a_{i,k}^\top x_k \le b_i$ is the row
$\sum_k (-a_{i,k})^\top x_k \ge -b_i$, so the form~\eqref{eq:compact-ineq} covers both
directions. For the presentation we therefore only consider ``$\ge$'' rows. In
\cref{sec:unified} we show that a ``$\le$'' row can equivalently be kept as it is, with a
surplus variable of the opposite sign. This is how such rows are handled in the
implementation.
\end{remark}

\begin{example}[Crew scheduling] \label{ex:crew}
In the crew scheduling problem studied in \cref{chap:results}, block $k$ is the schedule of
pilot $k$: a path from a source to a sink in a time--space network, with $F_k$ the
node--edge incidence matrix of the network. The linking rows are the \emph{coverage rows}:
row $i$ counts how often flight $i$ is operated, summed over all pilots. An equality row
``$=1$'' requires that the flight is operated exactly once, and an inequality row ``$\ge 1$''
requires it at least once, which allows a pilot to travel on a flight as a passenger.
\end{example}


\section{Dantzig--Wolfe Reformulation}
\label{sec:dantzig-wolfe}

Since $P_k$ is a nonempty polytope, it is the convex hull of its finitely many extreme points
(Minkowski's theorem). Let $\mathcal{J}_k$ be an index set of the extreme points of $P_k$,
and let $x^{(j)}$, $j\in\mathcal{J}_k$, denote these extreme points. Every $x_k\in P_k$ can be
written as
\begin{equation} \label{eq:convex-combination}
x_k = \sum_{j\in\mathcal{J}_k} \lambda_j\, x^{(j)}, \qquad
\sum_{j\in\mathcal{J}_k}\lambda_j = 1, \qquad \lambda_j \ge 0 .
\end{equation}
Substituting~\eqref{eq:convex-combination} into~\eqref{eq:compact} replaces the variables
$x_k$ by the weights $\lambda_j$. Each extreme point turns into one column of a new problem,
which motivates the following definitions.

\begin{definition}[Columns] \label{def:columns}
Let $\mathcal{J} = \mathcal{J}_1\cup\dots\cup\mathcal{J}_K$ (disjoint union). For
$j\in\mathcal{J}_k$ we define
\begin{itemize}
    \item the \textbf{block} of column $j$: $\psi(j) = k$;
    \item the \textbf{cost} of column $j$: $c_j = c_k^\top x^{(j)}$;
    \item the \textbf{linking coefficients} of column $j$:
          $a_j^E = A_{E,k}\,x^{(j)}\in\R^{m_E}$, $a_j^I = A_{I,k}\,x^{(j)}\in\R^{m_I}$, and
          $a_j = \bigl[\,a_j^E;\, a_j^I\,\bigr]\in\R^{m}$.
\end{itemize}
We collect them in the cost vector $c = (c_j)_{j\in\mathcal{J}}$ and the matrices
\[
A_E = \bigl[\, a_j^E \,\bigr]_{j\in\mathcal{J}}\in\R^{m_E\times|\mathcal{J}|}, \qquad
A_I = \bigl[\, a_j^I \,\bigr]_{j\in\mathcal{J}}\in\R^{m_I\times|\mathcal{J}|}, \qquad
A = \begin{bmatrix} A_E \\ A_I\end{bmatrix}\in\R^{m\times|\mathcal{J}|}.
\]
\end{definition}

\begin{definition}[Convexity matrix] \label{def:convexity}
The \textbf{convexity matrix} $C\in\{0,1\}^{K\times|\mathcal{J}|}$ has the entries
$C_{kj} = 1$ if $\psi(j) = k$ and $C_{kj} = 0$ otherwise. The constraint $C\lambda = \one$
states that the weights of the columns of every block sum to one.
\end{definition}

\begin{definition}[Master problem] \label{def:master}
The \textbf{master problem} is the linear program
\begin{subequations} \label{eq:master}
\begin{align}
\min_{\lambda} \quad & c^\top \lambda \\
\text{s.t.} \quad & A_E\,\lambda = b_E, && (\pi_E) \label{eq:master-eq}\\
& A_I\,\lambda \ge b_I, && (\pi_I) \label{eq:master-ineq}\\
& C\lambda = \one, && (\mu) \label{eq:master-conv}\\
& \lambda \ge 0, \label{eq:master-nonneg}
\end{align}
\end{subequations}
where the symbols in parentheses denote the dual variables of the constraints.
\end{definition}

\begin{theorem}[Dantzig--Wolfe \cite{dantzig1960decomposition}] \label{thm:dw}
Under \cref{ass:bounded}, problem~\eqref{eq:master} has the optimal value $z^\star$. If
$\lambda$ is optimal for~\eqref{eq:master}, then $x_k = \sum_{j\in\mathcal{J}_k}\lambda_j\,
x^{(j)}$, $k=1,\dots,K$, is optimal for~\eqref{eq:compact}.
\end{theorem}

\begin{proof}
By~\eqref{eq:convex-combination}, $x_k\in P_k$ if and only if $x_k$ is a convex combination
of the $x^{(j)}$, $j\in\mathcal{J}_k$. Under this substitution the objective and the
linking rows are linear in $\lambda$, with exactly the coefficients of \cref{def:columns}.
The feasible points of~\eqref{eq:compact} are therefore the images of the feasible points
of~\eqref{eq:master}, with equal objective values. 
\end{proof}

The convexity matrix has a very particular structure, which will be essential for the
nullspace method of \cref{chap:nullspace}.

\begin{definition}[Column generation property] \label{def:cg-property}
A matrix $C\in\R^{K\times n}$ has the \textbf{column generation property} if
\begin{enumerate}
    \item every entry of $C$ is $0$ or $1$;
    \item every column of $C$ contains exactly one entry equal to $1$;
    \item every row of $C$ contains at least one entry equal to $1$.
\end{enumerate}
\end{definition}

\begin{lemma} \label{lem:cg-property}
The convexity matrix $C$ of \cref{def:convexity} has the column generation property, and so
does every submatrix of $C$ formed by a subset of its columns that contains at least one
column of every block.
\end{lemma}

\begin{proof}
Property~(i) holds by definition. Every column belongs to exactly one block, which gives
property~(ii) for every subset of columns. Property~(iii) holds for a subset of columns if
and only if every block is represented in it. 
\end{proof}

\begin{lemma} \label{lem:C-rank}
If $C\in\R^{K\times n}$ has the column generation property, then $\rank(C) = K$.
\end{lemma}

\begin{proof}
By property~(iii), for every row $k$ there is a column $j_k$ with $C_{kj_k} = 1$. By
property~(ii), this column has no other nonzero entry, i.e.\ it equals the unit vector
$e_k$. The columns $j_1,\dots,j_K$ therefore form the identity matrix $I_K$. 
\end{proof}

\subsection*{The dual of the master problem}

The dual of~\eqref{eq:master} has one constraint per column and the sign condition of the
inequality rows:
\begin{subequations} \label{eq:master-dual}
\begin{align}
\max_{\pi_E,\,\pi_I,\,\mu} \quad & b_E^\top \pi_E + b_I^\top\pi_I + \one^\top\mu \\
\text{s.t.} \quad & c_j - a_j^\top \pi - \mu_{\psi(j)} \ge 0, \qquad j\in\mathcal{J},
\label{eq:dual-columns}\\
& \pi_I \ge 0, \label{eq:dual-sign}
\end{align}
\end{subequations}
where $\pi = [\pi_E;\pi_I]\in\R^m$. Compared with a problem that has only equality linking
rows, the inequality rows add the \textbf{sign condition}~\eqref{eq:dual-sign}. This is the
only change in the dual, but much of the work in the following chapters is devoted to it.


\section{Column Generation}
\label{sec:column-generation}

The index set $\mathcal{J}$ is typically far too large to enumerate. In crew scheduling, for
example, $\mathcal{J}_k$ contains every feasible schedule of pilot $k$. Column generation
solves~\eqref{eq:master} by working with a small subset of the columns at a time.

\begin{definition}[Restricted master problem] \label{def:rmp}
For a subset $\mathcal{J}'\subseteq\mathcal{J}$ that contains at least one column of every
block, the \textbf{restricted master problem} is problem~\eqref{eq:master} with all columns
outside $\mathcal{J}'$ removed (equivalently, with $\lambda_j = 0$ fixed for
$j\notin\mathcal{J}'$). We write $n = |\mathcal{J}'|$ for the number of columns in the
restricted master problem.
\end{definition}

\begin{definition}[Reduced cost] \label{def:reduced-cost}
For dual values $(\pi,\mu)$ the \textbf{reduced cost} of a column $j\in\mathcal{J}$ is
\[
\bar c_j(\pi,\mu) \;=\; c_j - a_j^\top\pi - \mu_{\psi(j)} .
\]
\end{definition}

The reduced cost is the slack of the dual constraint~\eqref{eq:dual-columns}. A column with
a negative reduced cost violates the dual constraints and can improve the restricted master
problem. Finding the column of most negative reduced cost in block $k$ is the
\textbf{pricing problem}
\begin{equation} \label{eq:pricing}
\min_{j\in\mathcal{J}_k}\bar c_j(\pi,\mu)
\;=\; \min_{x_k\in P_k}\bigl(c_k - A_k^\top\pi\bigr)^\top x_k \;-\;\mu_k,
\qquad A_k = \begin{bmatrix}A_{E,k}\\ A_{I,k}\end{bmatrix},
\end{equation}
where the equality holds because a linear function attains its minimum over a polytope at an
extreme point. The pricing problem is a problem over a single block. In crew scheduling it
is a shortest path problem with edge weights $c_{e,k} - (A_k^\top\pi)_e$.

\begin{prop}[Optimality test] \label{prop:cg-optimality}
Let $\lambda'$ be optimal for the restricted master problem, with optimal duals
$(\pi,\mu)$. If $\pi_I\ge0$ and the pricing problem~\eqref{eq:pricing} has a nonnegative
optimal value for every block, then $\lambda'$, extended by zeros, is optimal for the master
problem~\eqref{eq:master}.
\end{prop}

\begin{proof}
The extended $\lambda'$ is feasible for~\eqref{eq:master}. The pair $(\pi,\mu)$ satisfies
the sign condition~\eqref{eq:dual-sign} by assumption and every dual
constraint~\eqref{eq:dual-columns} by the pricing result, so it is dual feasible. Complementary
slackness holds for the columns in $\mathcal{J}'$ by optimality of the restricted problem,
and trivially for the zero weights outside $\mathcal{J}'$. 
\end{proof}

For a linear restricted master problem the duals satisfy $\pi_I\ge0$ automatically, so the
pricing test alone decides optimality. With the regularized master problems of the next
chapter this will no longer be the case. Column generation also provides a lower bound on
$z^\star$ at every iteration in which all blocks are priced.

\begin{prop}[Lagrangian bound] \label{prop:lagrangian-bound}
For every $\pi = [\pi_E;\pi_I]$ with $\pi_I\ge0$, the \textbf{Lagrangian bound}
\[
\theta(\pi) \;=\; b^\top\pi + \sum_{k=1}^{K}\,\min_{x_k\in P_k}\bigl(c_k - A_k^\top\pi\bigr)^\top x_k,
\qquad b = \begin{bmatrix}b_E\\ b_I\end{bmatrix},
\]
satisfies $\theta(\pi)\le z^\star$.
\end{prop}

\begin{proof}
Let $x^\star$ be optimal for~\eqref{eq:compact}. Because $\sum_k A_{E,k}x^\star_k = b_E$,
$\sum_k A_{I,k}x^\star_k\ge b_I$ and $\pi_I\ge0$, we have
$\pi^\top\bigl(b - \sum_k A_k x_k^\star\bigr)\le 0$. Hence
\[
z^\star = \sum_k c_k^\top x^\star_k
\;\ge\; \sum_k c_k^\top x^\star_k + \pi^\top\Bigl(b - \sum_k A_k x_k^\star\Bigr)
= b^\top\pi + \sum_k\bigl(c_k - A_k^\top\pi\bigr)^\top x^\star_k
\;\ge\; \theta(\pi),
\]
since $x_k^\star\in P_k$. 
\end{proof}

\begin{remark} \label{rem:bound-sign}
The sign condition $\pi_I\ge0$ in \cref{prop:lagrangian-bound} is essential: for $\pi_i<0$
in an inequality row, the term $\pi_i\,(b_i - \sum_k a_{i,k}^\top x_k)$ can become arbitrarily
large. In \cref{sec:center-update} the bound is therefore evaluated at the projection of
the dual vector onto $\{\pi_I\ge0\}$.
\end{remark}

The resulting method is summarized in \cref{alg:classical-cg}. The initial columns are
usually obtained by solving the pricing problem once with zero duals, which gives one
column per block.

\begin{algobox}{Classical Column Generation}{alg:classical-cg}
\State \textbf{Input:} initial columns $\mathcal{J}'$ with at least one column per block
\Repeat
    \State Solve the (linear) restricted master problem over $\mathcal{J}'$; obtain
           $\lambda'$ and duals $(\pi,\mu)$
    \For{each block $k = 1,\dots,K$}
        \State Solve the pricing problem~\eqref{eq:pricing}; let $j$ be the best column
        \If{$\bar c_j(\pi,\mu) < 0$}
            \State $\mathcal{J}'\gets\mathcal{J}'\cup\{j\}$
        \EndIf
    \EndFor
\Until{no column with negative reduced cost was found}
\State \textbf{Output:} $\lambda'$, optimal for the master problem by \cref{prop:cg-optimality}
\end{algobox}

\begin{remark}[Instability] \label{rem:instability}
The duals of a linear restricted master problem are extreme points of its dual feasible
region. When a column is added, the dual solution can jump to a very different extreme point,
even if the primal solution hardly changes. As a result, the duals oscillate strongly from
one iteration to the next, the pricing problem generates many columns that are never used,
and the method converges slowly in its final phase (the \emph{tailing-off effect}). In the
dual space, column generation is Kelley's cutting-plane method
\cite{kelley1960cutting} applied to the Lagrangian dual function, a method known for exactly
this behaviour. This is the motivation for the regularization of \cref{chap:regularization}.
\end{remark}


\section{Surplus Variables and the Unified Formulation}
\label{sec:unified}

The two kinds of linking rows in~\eqref{eq:master} behave differently: the equality rows
fix $A_E\lambda$, while the inequality rows only bound $A_I\lambda$ from below. To treat both
in the same way, we introduce for every inequality row $i$ a \textbf{surplus variable}
$s_i\ge0$ with
\[
(A_I\lambda)_i - s_i = (b_I)_i, \qquad\text{i.e.}\qquad A_I\lambda - s = b_I .
\]
The surplus measures how much the row is over-satisfied.

\begin{definition}[Unified formulation] \label{def:unified}
Let $s\in\R^{m_I}$ be the vector of surplus variables. The \textbf{unified variable} and the
\textbf{augmented data} are
\[
y = \begin{bmatrix}\lambda\\ s\end{bmatrix}, \qquad
\Abar = \begin{bmatrix} A_E & 0\\ A_I & -I\end{bmatrix}, \qquad
\Cbar = \begin{bmatrix} C & 0\end{bmatrix}, \qquad
\cbar = \begin{bmatrix} c\\ 0\end{bmatrix}, \qquad
b = \begin{bmatrix}b_E\\ b_I\end{bmatrix}.
\]
The \textbf{unified master problem} is
\begin{equation} \label{eq:unified}
\min_{y}\;\cbar^\top y \qquad\text{s.t.}\qquad \Abar y = b,\quad \Cbar y = \one,\quad y\ge0 .
\end{equation}
\end{definition}

\begin{lemma} \label{lem:unified}
Problems~\eqref{eq:master} and~\eqref{eq:unified} are equivalent: $\lambda$ is feasible
for~\eqref{eq:master} if and only if $y = [\lambda;\, A_I\lambda - b_I]$ is feasible
for~\eqref{eq:unified}, with the same objective value. Moreover, the dual
of~\eqref{eq:unified} is the dual~\eqref{eq:master-dual} of the master problem.
\end{lemma}

\begin{proof}
The first statement follows from $s = A_I\lambda - b_I\ge0$. For the second, the dual
of~\eqref{eq:unified} has multipliers $\pi$ for $\Abar y = b$ and $\mu$ for $\Cbar y = \one$,
and one constraint per column of $[\Abar;\Cbar]$. For the columns of $\lambda$ these are the
constraints~\eqref{eq:dual-columns}. The column of surplus $s_i$ has cost $0$, coefficient
$-e_i$ in the inequality rows and $0$ in the convexity rows, so its dual constraint reads
$0 - (-e_i)^\top\pi_I - 0 \ge 0$, i.e.\ $(\pi_I)_i\ge0$. These are exactly the sign
conditions~\eqref{eq:dual-sign}. 
\end{proof}

\Cref{lem:unified} shows that the surplus columns \emph{encode the sign conditions as
ordinary dual constraints}. Problem~\eqref{eq:unified} has the same form as a master problem
with equality linking rows only, with $(A, C, c)$ replaced by $(\Abar,\Cbar,\cbar)$. All the
following chapters work with this form.

\begin{remark}[The column generation property is lost] \label{rem:cg-property-lost}
The augmented convexity matrix $\Cbar = [\,C\;\;0\,]$ has $m_I$ zero columns, one per
surplus variable, so it violates property~(ii) of \cref{def:cg-property}. The surpluses do
not belong to any block. \Cref{chap:qpas} shows that this causes no difficulty: zero
columns are even simpler to handle than unit columns, but they have to be treated as a
separate group of variables.
\end{remark}

\begin{remark}[``$\le$'' rows and the slack sign] \label{rem:kappa}
A row $a_i^\top\lambda\le b_i$ can be written with a \emph{slack} $s_i\ge0$ as
$a_i^\top\lambda + s_i = b_i$. In general, let $\kappa_i\in\{-1,+1\}$ be the coefficient of
the slack of inequality row $i$ ($\kappa_i = -1$ for a ``$\ge$'' row, $\kappa_i = +1$ for a
``$\le$'' row) and $K_s = \diag(\kappa)$. The unified formulation then uses
\[
\Abar = \begin{bmatrix} A_E & 0\\ A_I & K_s\end{bmatrix}
\]
instead of \cref{def:unified}, and the dual constraint of slack $i$ becomes
$-\kappa_i(\pi_I)_i\ge0$, i.e.\ $(\pi_I)_i\le0$ for a ``$\le$'' row. Since
$\kappa_i^2 = 1$, every result of the following chapters holds for both kinds of rows; only a
few formulas carry a factor $\kappa_i$, and these are collected in \cref{rem:kappa-summary}.
In the text we write all formulas for ``$\ge$'' rows ($K_s = -I$).
\end{remark}

\begin{remark}[Range rows] \label{rem:range}
A row with two different finite bounds $\ell_i\le a_i^\top\lambda\le u_i$ needs a slack with
an upper bound. Such rows are not considered in this work.
\end{remark}


\section{Summary of Notation}

\begin{table}[ht]
\centering
\renewcommand{\arraystretch}{1.2}
\begin{tabular}{ll}
\toprule
symbol & meaning \\
\midrule
$K$, $k$ & number of blocks, block index \\
$m_E$, $m_I$, $m$ & number of equality, inequality and all linking rows \\
$\mathcal{J}$, $\mathcal{J}_k$, $\mathcal{J}'$ & all columns, columns of block $k$, columns of the restricted problem \\
$n$ & number of columns in the restricted master problem \\
$\psi(j)$ & block of column $j$ \\
$c_j$, $a_j = [a_j^E; a_j^I]$ & cost and linking coefficients of column $j$ \\
$A_E$, $A_I$, $A$ & linking coefficient matrices of the columns \\
$C$ & convexity matrix \\
$\lambda$, $s$, $y = [\lambda; s]$ & column weights, surpluses, unified variable \\
$\Abar$, $\Cbar$, $\cbar$ & augmented data of the unified formulation \\
$\pi = [\pi_E;\pi_I]$, $\mu$ & duals of the linking rows and of the convexity rows \\
$\bar c_j(\pi,\mu)$ & reduced cost of column $j$ \\
$\theta(\pi)$ & Lagrangian bound \\
\bottomrule
\end{tabular}
\caption{Notation of \cref{chap:formulation}.}
\end{table}

\chapter{Regularization of the Master Problem}
\label{chap:regularization}

\vspace*{\fill}

This chapter replaces the linear restricted master problem by a regularized one. The
regularization is formulated in the dual space, where it adds a quadratic penalty on the
distance between the dual iterate and a \textbf{stability center}. Translated back to the
primal space, it yields a convex quadratic program with an auxiliary variable, the
\textbf{explicit regularized master problem}.

\medskip

Because all linking rows of the unified formulation are equalities, the auxiliary variable
can be eliminated for equality and inequality rows alike. The result, the
\textbf{eliminated regularized master problem}, is a convex quadratic program whose only
constraints are the convexity constraints and the nonnegativity of the variables. This is the
problem that the active-set method of the following chapters solves.

\medskip

Two identities derived in this chapter are used throughout the rest of the work. The first
expresses the duals of the linking rows as an explicit function of the primal iterate. The
second shows that the gradient of the quadratic objective consists of the reduced costs of
the columns and the duals of the inequality rows. The chapter ends with the regularized
column generation algorithm, its exact stopping criterion and the accuracy that this
criterion guarantees.

\vspace*{\fill}

\newpage


\section{Proximal Stabilization of the Dual}
\label{sec:proximal}

Consider the restricted master problem in the unified form~\eqref{eq:unified}, over the
columns $\mathcal{J}'$ and all $m_I$ surpluses. Its dual is
\begin{equation} \label{eq:restricted-dual}
\max_{\pi,\,\mu}\; b^\top\pi + \one^\top\mu
\qquad\text{s.t.}\qquad
c_j - a_j^\top\pi - \mu_{\psi(j)}\ge0 \;\;(j\in\mathcal{J}'),\qquad \pi_I\ge0 .
\end{equation}
Its optimal dual solutions are extreme points of the feasible region, which causes the
instability described in \cref{rem:instability}. Proximal stabilization adds a penalty on the
distance to a given point $\pihat\in\R^m$.

\begin{definition}[Stabilized dual problem] \label{def:stabilized-dual}
Let $\pihat\in\R^m$ be the \textbf{stability center} and $\delta>0$ the
\textbf{stabilization parameter}. The \textbf{stabilized dual problem} is
\begin{equation} \label{eq:stabilized-dual}
\max_{\pi,\,\mu}\; b^\top\pi + \one^\top\mu - \frac{\delta}{2}\,\lVert\pi-\pihat\rVert^2
\qquad\text{s.t.}\qquad
c_j - a_j^\top\pi - \mu_{\psi(j)}\ge0 \;\;(j\in\mathcal{J}'),\qquad \pi_I\ge0 .
\end{equation}
\end{definition}

The objective of~\eqref{eq:stabilized-dual} is strictly concave in $\pi$, so the optimal
$\pi$ is unique. A large $\delta$ keeps $\pi$ close to $\pihat$ (strong stabilization), and
$\delta\to0$ recovers~\eqref{eq:restricted-dual}. Replacing the center $\pihat$ by the
solution $\pi$ and solving again is one step of the \textbf{proximal point method}
\cite{rockafellar1976monotone} applied to the restricted dual problem. If the center is only
moved when it improves the Lagrangian bound, we obtain a \textbf{proximal bundle method}
\cite{kiwiel1990proximity, briant2008comparison}. Both rules are discussed in
\cref{sec:center-update}.


\section{The Explicit Regularized Master Problem}
\label{sec:explicit}

The primal counterpart of~\eqref{eq:stabilized-dual} has an auxiliary vector
$w\in\R^m$, one entry per linking row.

\begin{definition}[Explicit regularized master problem] \label{def:explicit}
The \textbf{explicit regularized master problem} is
\begin{subequations} \label{eq:explicit}
\begin{align}
\min_{y,\,w}\quad & \cbar^\top y + \frac{\delta}{2}\lVert w + \pihat\rVert^2
                     - \frac{\delta}{2}\lVert\pihat\rVert^2 \\
\text{s.t.}\quad & \Abar y + \delta w = b, && (\pi) \label{eq:explicit-link}\\
& \Cbar y = \one, && (\mu)\\
& y\ge0 .
\end{align}
\end{subequations}
\end{definition}

\begin{theorem} \label{thm:explicit-duality}
The dual of~\eqref{eq:explicit} is the stabilized dual problem~\eqref{eq:stabilized-dual}.
Strong duality holds, and at a pair of optimal solutions
\begin{equation} \label{eq:pi-w}
\pi = \pihat + w .
\end{equation}
\end{theorem}

\begin{proof}
The Lagrangian of~\eqref{eq:explicit}, with multipliers $\pi$, $\mu$ and $\nu\ge0$ for
$y\ge0$, is
\[
\mathcal{L} = \cbar^\top y + \tfrac{\delta}{2}\lVert w+\pihat\rVert^2
- \tfrac{\delta}{2}\lVert\pihat\rVert^2
+ \pi^\top(b - \Abar y - \delta w) + \mu^\top(\one - \Cbar y) - \nu^\top y .
\]
It is a strictly convex quadratic function of $w$, minimized where
$\delta(w + \pihat) - \delta\pi = 0$, i.e.\ at $w = \pi - \pihat$. Inserting this,
\[
\tfrac{\delta}{2}\lVert\pi\rVert^2 - \tfrac{\delta}{2}\lVert\pihat\rVert^2
- \delta\pi^\top(\pi-\pihat) = -\tfrac{\delta}{2}\lVert\pi-\pihat\rVert^2 .
\]
The Lagrangian is linear in $y$, and its infimum over $y$ is finite only if
$\cbar - \Abar^\top\pi - \Cbar^\top\mu - \nu = 0$ for some $\nu\ge0$, i.e.\ only if
$\cbar - \Abar^\top\pi - \Cbar^\top\mu\ge0$. The dual function is therefore
$b^\top\pi + \one^\top\mu - \frac{\delta}{2}\lVert\pi-\pihat\rVert^2$, subject to this
condition. Written per column, the condition for the columns of $\lambda$
is~\eqref{eq:dual-columns}, and the condition for the surplus columns is $\pi_I\ge0$
(\cref{lem:unified}). This is~\eqref{eq:stabilized-dual}. Problem~\eqref{eq:explicit} is a
feasible convex quadratic program with linear constraints (it is feasible because $w$ is
free), so strong duality holds, and the minimizing $w$ satisfies~\eqref{eq:pi-w}. 
\end{proof}

\begin{remark} \label{rem:sign-in-dual}
\Cref{thm:explicit-duality} shows that the sign condition $\pi_I\ge0$ is part of the
stabilized dual problem automatically, through the surplus columns. The stability center
$\pihat$ itself may violate it. The proximal step then moves $\pi$ back into the region
$\pi_I\ge0$.
\end{remark}

\begin{remark} \label{rem:why-surplus}
Without surplus variables, an inequality row $(A_I\lambda)_i\ge(b_I)_i$ would enter the
explicit problem as $(A_I\lambda)_i + \delta w_i\ge(b_I)_i$, and $w_i$ would no longer be
determined by $\lambda$. The surpluses make every row of~\eqref{eq:explicit-link} an
equality, so that $w$ can be eliminated for all rows in the next section.
\end{remark}


\section{The Eliminated Regularized Master Problem}
\label{sec:eliminated}

From~\eqref{eq:explicit-link} we get $w = (b - \Abar y)/\delta$ for every feasible point.
Substituting this into the objective of~\eqref{eq:explicit} removes $w$ completely.

\begin{theorem}[Eliminated regularized master problem] \label{thm:eliminated}
Problem~\eqref{eq:explicit} is equivalent to the convex quadratic program
\begin{subequations} \label{eq:eliminated}
\begin{align}
\min_{y}\quad & f(y) = \frac{1}{2\delta}\, y^\top\Abar^\top\Abar\, y + \ctil^\top y \\
\text{s.t.}\quad & \Cbar y = \one, \label{eq:eliminated-conv}\\
& y\ge0, \label{eq:eliminated-bounds}
\end{align}
\end{subequations}
with the linear cost
\begin{equation} \label{eq:ctilde}
\ctil = \cbar - \Abar^\top\Bigl(\pihat + \frac{b}{\delta}\Bigr).
\end{equation}
More precisely, $(y, w)$ is optimal for~\eqref{eq:explicit} if and only if $y$ is optimal
for~\eqref{eq:eliminated} and $w = (b - \Abar y)/\delta$. The objective values are related by
\begin{equation} \label{eq:objective-relation}
\Phi(y) := \cbar^\top y + \frac{\delta}{2}\lVert w+\pihat\rVert^2 - \frac{\delta}{2}\lVert\pihat\rVert^2
= f(y) + \frac{1}{2\delta}\lVert b\rVert^2 + \pihat^\top b .
\end{equation}
\end{theorem}

\begin{proof}
Let $r = b - \Abar y$, so $w = r/\delta$. Then
\[
\tfrac{\delta}{2}\lVert r/\delta + \pihat\rVert^2 - \tfrac{\delta}{2}\lVert\pihat\rVert^2
= \tfrac{1}{2\delta}\lVert r\rVert^2 + \pihat^\top r
= \tfrac{1}{2\delta}\lVert\Abar y\rVert^2 - \tfrac{1}{\delta}b^\top\Abar y
  + \tfrac{1}{2\delta}\lVert b\rVert^2 + \pihat^\top b - \pihat^\top\Abar y .
\]
Adding $\cbar^\top y$ and collecting the terms that are linear in $y$ gives $f(y)$ plus the
constant in~\eqref{eq:objective-relation}. The constraint~\eqref{eq:explicit-link} only
defines $w$, so the remaining constraints are~\eqref{eq:eliminated-conv}
and~\eqref{eq:eliminated-bounds}. 
\end{proof}

The Hessian of $f$ is $\frac{1}{\delta}\Abar^\top\Abar$, with the block structure
\begin{equation} \label{eq:hessian-blocks}
\frac{1}{\delta}\Abar^\top\Abar
= \frac{1}{\delta}\begin{bmatrix} A_E^\top A_E + A_I^\top A_I & -A_I^\top\\ -A_I & I\end{bmatrix}.
\end{equation}
It is positive semidefinite but in general singular: it has rank at most $m$, while the
number of variables $n + m_I$ grows with every generated column. Problem~\eqref{eq:eliminated}
nevertheless has an optimal solution, because its feasible region is bounded in the $\lambda$
part and the objective grows quadratically in $s$.

\subsection{The dual vector as a function of the primal iterate}

Relation~\eqref{eq:pi-w} with $w = (b-\Abar y)/\delta$ suggests the following definition,
which applies to \emph{every} $y$, not only to the optimal one.

\begin{definition}[Dual estimate] \label{def:dual-estimate}
For $y = [\lambda; s]$ the \textbf{dual estimate} is
\begin{equation} \label{eq:pi-of-y}
\pi(y) = \pihat + \frac{1}{\delta}\bigl(b - \Abar y\bigr),
\qquad\text{i.e.}\qquad
\begin{aligned}
\pi_E(y) &= \pihat_E + \tfrac{1}{\delta}\bigl(b_E - A_E\lambda\bigr),\\
\pi_I(y) &= \pihat_I + \tfrac{1}{\delta}\bigl(b_I - A_I\lambda + s\bigr).
\end{aligned}
\end{equation}
\end{definition}

A surplus raises the dual of its row: an over-covered inequality row with a positive
surplus has a larger $\pi_i$ than it would have without the surplus. The next lemma is the
central identity of this work.

\begin{lemma}[Gradient--dual relation] \label{lem:gradient-dual}
For every $y$,
\begin{equation} \label{eq:gradient-dual}
\nabla f(y) = \cbar - \Abar^\top\pi(y),
\qquad\text{i.e.}\qquad
g_j := \frac{\partial f}{\partial\lambda_j} = c_j - a_j^\top\pi(y), \qquad
\frac{\partial f}{\partial s_i} = \pi_i(y)\quad (i\in\{1,\dots,m_I\}).
\end{equation}
\end{lemma}

\begin{proof}
$\nabla f(y) = \frac{1}{\delta}\Abar^\top\Abar y + \cbar - \Abar^\top\pihat - \frac{1}{\delta}\Abar^\top b
= \cbar - \Abar^\top\bigl(\pihat + \frac{1}{\delta}(b - \Abar y)\bigr)$. The rows of
$\Abar^\top$ belonging to $\lambda_j$ are $a_j^\top$, and the rows belonging to $s_i$ are
$-e_i^\top$ (in the inequality part). 
\end{proof}

\Cref{lem:gradient-dual} has two consequences that are used throughout this work.
\begin{itemize}
    \item The gradient component of a column, $g_j = c_j - a_j^\top\pi$, is the
          \textbf{reduced cost of column $j$ without the convexity dual}. The gradient
          component of a surplus is the dual of its row.
    \item The gradient, an $(n+m_I)$-vector whose length grows with every column, is
          determined by the $m$-vector $\pi(y)$. An implementation can store $\pi$ and
          evaluate $g_j$ only where it is needed.
\end{itemize}

\begin{prop}[Objective value] \label{prop:objective-value}
For every $y$, with $\pi = \pi(y)$,
\[
\Phi(y) = c^\top\lambda + \frac{\delta}{2}\lVert\pi\rVert^2 - \frac{\delta}{2}\lVert\pihat\rVert^2 .
\]
\end{prop}

\begin{proof}
In~\eqref{eq:objective-relation}, $\Phi(y) = \cbar^\top y + \frac{\delta}{2}\lVert w+\pihat\rVert^2
- \frac{\delta}{2}\lVert\pihat\rVert^2$ with $w+\pihat = \pi(y)$, and $\cbar^\top y = c^\top\lambda$
because the surpluses have zero cost. 
\end{proof}

\begin{prop}[Uniqueness of the dual] \label{prop:pi-unique}
The optimal value of $\Abar y$, and hence the optimal dual vector $\pi$, is the same for all
optimal solutions of~\eqref{eq:eliminated}, even if the optimal $y$ is not unique.
\end{prop}

\begin{proof}
Let $y_1, y_2$ be optimal with $\Abar y_1\ne\Abar y_2$. Their midpoint $\bar y$ is feasible.
The linear part of $f$ is affine, while $\lVert\Abar y\rVert^2$ is strictly convex as a
function of $\Abar y$, so $f(\bar y) < \frac12\bigl(f(y_1)+f(y_2)\bigr)$, a contradiction.

\end{proof}

\begin{prop}[Residual and dual distance] \label{prop:residual}
For every $y$,
\[
b - \Abar y = \delta\bigl(\pi(y) - \pihat\bigr).
\]
In particular, $y$ is feasible for the unified master problem~\eqref{eq:unified} if and only
if $\pi(y) = \pihat$.
\end{prop}

\begin{proof}
This is~\eqref{eq:pi-of-y} rearranged. 
\end{proof}


\section{Regularized Column Generation}
\label{sec:center-update}

Column generation with the regularized master problem follows the classical scheme of
\cref{alg:classical-cg}, with two differences: the duals passed to the pricing problem are
the stabilized duals $\pi(y)$, and the stability center has to be updated.

\subsection{Center updates}

Two rules for the center update are common.
\begin{itemize}
    \item \textbf{Proximal point rule:} after every master problem, $\pihat\gets\pi$. Every
          iteration is a step of the proximal point method on the current restricted dual
          problem.
    \item \textbf{Bundle rule:} $\pihat\gets\pi$ only if the Lagrangian bound improves,
          $\theta(\pi) > \theta(\pihat)$ (a \emph{serious step}). Otherwise the center stays
          (a \emph{null step}). This requires the Lagrangian bound at every iteration, i.e.\
          all blocks must be priced every iteration.
\end{itemize}

\subsection{Exact termination}

In classical column generation, the pricing test alone decides optimality
(\cref{prop:cg-optimality}). With regularization this is no longer true: when pricing finds
no improving column, the regularized master problem may still be infeasible for the
unregularized one, since $\Abar y\ne b$ as long as $\pi\ne\pihat$ (\cref{prop:residual}).

\begin{theorem}[Exact termination] \label{thm:exact-termination}
Let $y = [\lambda; s]$ be optimal for the eliminated regularized master
problem~\eqref{eq:eliminated} over the columns $\mathcal{J}'$, with dual estimate
$\pi = \pi(y)$ and convexity duals $\mu$. If
\begin{enumerate}
    \item $\pi = \pihat$, and
    \item the pricing problem~\eqref{eq:pricing} has a nonnegative optimal value for every
          block,
\end{enumerate}
then $\lambda$, extended by zeros, is optimal for the master problem~\eqref{eq:master}, and
$(\pi,\mu)$ is optimal for its dual~\eqref{eq:master-dual}.
\end{theorem}

\begin{proof}
By \cref{prop:residual}, $\pi = \pihat$ gives $\Abar y = b$, i.e.\ $A_E\lambda = b_E$ and
$A_I\lambda = b_I + s\ge b_I$, so $\lambda$ is primal feasible. The optimality conditions
of~\eqref{eq:eliminated}, which are derived in \cref{thm:kkt-master} below, state that every
column of $\mathcal{J}'$ has a nonnegative reduced cost (with equality where $\lambda_j>0$),
that $\pi_I\ge0$, and that $\pi_i = 0$ whenever $s_i>0$. Condition~(ii) extends the
nonnegativity of the reduced costs to all columns. So $(\pi,\mu)$ is dual feasible, and
complementary slackness holds for the columns, the zero weights outside $\mathcal{J}'$ and the
inequality rows. 
\end{proof}

If pricing finds no improving column while $\pi\ne\pihat$, the restricted master problem
contains the right columns for the moment, but the proximal term still keeps $\pi$ away from
the restricted dual optimum. The algorithm then moves the center to $\pi$ and solves the
master problem again \emph{without adding a column}. We call this a
\textbf{center-only iteration}. It is one proximal point step on the current restricted
model. Along the way, the duals move, and pricing may find improving columns again. A run of
regularized column generation therefore consists of two kinds of iterations, which
\cref{sec:results-phases} studies in detail: iterations that add columns, and center-only
iterations that drive $\pi$ towards $\pihat$.

\begin{prop}[Accuracy of the stopping test] \label{prop:stop-accuracy}
In practice condition~(i) of \cref{thm:exact-termination} is tested as
\begin{equation} \label{eq:center-test}
\lVert\pi - \pihat\rVert_\infty \;\le\; \tau\,\max\bigl(1,\lVert\pi\rVert_\infty\bigr)
\end{equation}
with a tolerance $\tau>0$. When this test holds, the violation of the linking rows by
$\lambda$ is bounded by
\[
\lVert b - \Abar y\rVert_\infty \;\le\; \delta\,\tau\,\max\bigl(1,\lVert\pi\rVert_\infty\bigr).
\]
\end{prop}

\begin{proof}
This follows from \cref{prop:residual}. 
\end{proof}

The accuracy of the final solution is thus relative to the size of the duals, and it scales
with $\delta$. The objective value $\Phi$ and the Lagrangian bound are much less affected by
this residual than the cost $c^\top\lambda$ of a slightly infeasible $\lambda$, as the
results in \cref{sec:results-correctness} show.

\subsection{The Lagrangian bound with inequality rows}

The Lagrangian bound of \cref{prop:lagrangian-bound} requires $\pi_I\ge0$. The duals of the
regularized master problem satisfy this condition only up to the tolerance of the master
solver, and a tiny negative entry would make the bound useless (\cref{rem:bound-sign}). The
bound is therefore evaluated at the projection
\[
\pi^+ = \bigl[\,\pi_E;\;\max(\pi_I, 0)\,\bigr],
\]
which always satisfies the sign condition, so that $\theta(\pi^+)\le z^\star$ holds by
\cref{prop:lagrangian-bound}. The same projected vector is passed to the pricing problem.
For problems with equality rows only, $\pi^+ = \pi$.

\subsection{The algorithm}

\begin{algobox}{Regularized Column Generation}{alg:rcg}
\State \textbf{Input:} initial columns $\mathcal{J}'$ (one per block), center $\pihat$,
       parameter $\delta>0$, tolerances $\tau$, $\varepsilon_{\mathrm{rc}}$
\Repeat
    \State Solve the eliminated regularized master problem~\eqref{eq:eliminated}; obtain $y$,
           $\pi = \pi(y)$ and $\mu$
    \State $\pi^+\gets[\pi_E;\max(\pi_I,0)]$
    \State Price the blocks in random order with $(\pi^+,\mu)$; add the first column with
           $\bar c_j < -\varepsilon_{\mathrm{rc}}$ to $\mathcal{J}'$
    \If{no column was added (all blocks priced)}
        \State $\theta\gets\theta(\pi^+)$ \Comment{valid lower bound on $z^\star$}
        \If{$\lVert\pi-\pihat\rVert_\infty\le\tau\max(1,\lVert\pi\rVert_\infty)$}
            \State \textbf{stop}: $\lambda$ is optimal (\cref{thm:exact-termination})
        \EndIf
        \State $\pihat\gets\pi$ \Comment{center-only iteration}
    \Else
        \State $\pihat\gets\pi$ \Comment{proximal point rule (bundle rule: only if $\theta$ improves)}
    \EndIf
\Until{stopped}
\end{algobox}

In the experiments of \cref{chap:results} we use the proximal point rule,
$\delta = 1$, $\tau = 10^{-6}$ and $\varepsilon_{\mathrm{rc}} = 10^{-6}$. The remaining
chapters develop the method used to solve the master problem in the first line of the loop.

\chapter{An Active-Set Method for the Regularized Master Problem}
\label{chap:qpas}
\chaptermark{Active-Set Method for the Master Problem}
\chaptermark{Active-Set Method for the Master Problem}

\vspace*{\fill}

This chapter sets up the \textbf{quadratic programming active-set method (QPAS)} for the
eliminated regularized master problem. We first derive its optimality conditions and show
that the Lagrange multipliers of the nonnegativity constraints are the reduced costs of the
columns and the duals of the inequality rows. The optimality test of the master problem and
the pricing test of column generation are therefore the same test.

\medskip

We then recall the primal active-set method: it fixes a \textbf{working set} of variables at
zero, minimizes over the remaining free variables, and adds or removes one constraint at a
time. Every iteration requires the solution of a KKT system whose constraint matrix consists
of the convexity rows and the rows of the working set.

\medskip

The main result of this chapter is a \textbf{reordering theorem} for this constraint matrix.
The columns of the regularized master problem are partitioned into five sets:
representatives, working-set columns, nullspace columns, active surpluses and free surpluses.
In this ordering the constraint matrix has a block structure with identity blocks, which
is the basis of the explicit nullspace method of the next chapter.

\vspace*{\fill}

\newpage


\section{The Master Problem in Simplified Notation}
\label{sec:simplified}

From now on we consider a single restricted master problem with $n$ columns. To keep the
notation light, we drop the explicit reference to the restricted index set and write
\begin{itemize}
    \item $x\in\R^n$ for the column weights (the vector $\lambda$ restricted to $\mathcal{J}'$),
    \item $A_E\in\R^{m_E\times n}$, $A_I\in\R^{m_I\times n}$, $A = [A_E;A_I]$ and
          $C\in\{0,1\}^{K\times n}$ for the corresponding columns of the linking and convexity
          matrices,
    \item $c\in\R^n$ for the costs of the columns, and $a_j\in\R^m$ for column $j$ of $A$,
    \item $y = [x; s]\in\R^{n+m_I}$ for the unified variable.
\end{itemize}
The vector $c$ always denotes the original column costs; the linear term of the quadratic
objective is $\ctil$ from~\eqref{eq:ctilde}. With this notation the eliminated regularized
master problem~\eqref{eq:eliminated} reads
\begin{subequations} \label{eq:qp}
\begin{align}
\min_{y\in\R^{n+m_I}}\quad & f(y) = \frac{1}{2\delta}\,y^\top\Abar^\top\Abar\,y + \ctil^\top y\\
\text{s.t.}\quad & \Cbar y = \one, \label{eq:qp-conv}\\
& y\ge0, \label{eq:qp-bounds}
\end{align}
\end{subequations}
with
\[
\Abar = \begin{bmatrix} A_E & 0\\ A_I & -I\end{bmatrix}\in\R^{m\times(n+m_I)},\qquad
\Cbar = \begin{bmatrix} C & 0\end{bmatrix}\in\R^{K\times(n+m_I)} .
\]
We write $g = \nabla f(y)$ for the gradient at the current iterate and, by
\cref{lem:gradient-dual},
\begin{equation} \label{eq:gradient-split}
g = \begin{bmatrix} g_x\\ g_s\end{bmatrix} = \begin{bmatrix} c - A^\top\pi\\ \pi_I\end{bmatrix},
\qquad \pi = \pi(y) = \pihat + \tfrac{1}{\delta}(b - \Abar y).
\end{equation}


\section{Optimality Conditions}
\label{sec:kkt}

Problem~\eqref{eq:qp} is a convex quadratic program with linear constraints, so the KKT
conditions are necessary and sufficient for optimality.

\begin{theorem}[KKT conditions of the regularized master problem] \label{thm:kkt-master}
A feasible point $y = [x;s]$ of~\eqref{eq:qp} is optimal if and only if there exist
$\mu\in\R^K$ and $\nu\in\R^{n+m_I}$ with
\[
g = \Cbar^\top\mu + \nu,\qquad \nu\ge0,\qquad \nu_\ell\,y_\ell = 0\quad\text{for all }\ell .
\]
Written per variable, with $\pi = \pi(y)$, these conditions are equivalent to:
\begin{enumerate}
    \item \textbf{reduced costs:} $\bar c_j(\pi,\mu) = c_j - a_j^\top\pi - \mu_{\psi(j)}\ge0$
          for every column $j$, with equality if $x_j>0$;
    \item \textbf{sign condition:} $\pi_i\ge0$ for every inequality row $i$;
    \item \textbf{slackness:} $\pi_i = 0$ for every inequality row $i$ with $s_i>0$.
\end{enumerate}
The multipliers are $\nu_j = \bar c_j(\pi,\mu)$ for the columns and $\nu_{n+i} = \pi_i$ for
the surpluses.
\end{theorem}

\begin{proof}
The first part is the KKT theorem for convex quadratic programs. For a column $j$, the
stationarity condition reads $g_j = \mu_{\psi(j)} + \nu_j$, so
$\nu_j = c_j - a_j^\top\pi - \mu_{\psi(j)} = \bar c_j(\pi,\mu)$ by~\eqref{eq:gradient-split}.
For surplus $i$ the column of $\Cbar$ is zero, so $\nu_{n+i} = g_{n+i} = \pi_i$. The
conditions $\nu\ge0$ and $\nu_\ell y_\ell = 0$ then become (i)--(iii). 
\end{proof}

\begin{remark}[The optimality test is the pricing test] \label{rem:same-test}
By \cref{thm:kkt-master}, the multipliers of the nonnegativity constraints of the columns
are their reduced costs, the same quantity that the pricing problem minimizes over the
columns that have not been generated yet. Solving the regularized master problem to
optimality establishes ``reduced cost $\ge0$'' for the columns inside the restricted problem,
and pricing establishes it for the columns outside. The multipliers of the surpluses carry
the sign conditions~\eqref{eq:dual-sign} of the inequality rows. Together with $\pi=\pihat$,
this is exactly the optimality certificate of \cref{thm:exact-termination}.
\end{remark}

\begin{remark}[Matching tolerances] \label{rem:tolerances}
The active-set method accepts a multiplier as nonnegative if $\nu_\ell\ge-\varepsilon$, and
pricing accepts a column only if $\bar c_j<-\varepsilon_{\mathrm{rc}}$. Choosing
$\varepsilon\ll\varepsilon_{\mathrm{rc}}$ (we use $\varepsilon = 10^{-10}\max(1,\max_j|c_j|)$
and $\varepsilon_{\mathrm{rc}} = 10^{-6}$) guarantees that pricing never returns a column that
is already in the restricted master problem.
\end{remark}


\section{The Primal Active-Set Method}
\label{sec:active-set}

\begin{definition}[Active set and working set] \label{def:working-set}
The \textbf{active set} of a feasible point $y$ is
$\mathcal{A}(y) = \{\ell : y_\ell = 0\}$. A \textbf{working set} is a set
$\Wset\subseteq\{1,\dots,n+m_I\}$ of indices that are treated as active, i.e.\ $y_\ell = 0$
for all $\ell\in\Wset$. We split it into the \textbf{column working set} $\Wx$ and the
\textbf{surplus working set} $\Ws$:
\[
\Wset = \Wx\cup\{\,n+i : i\in\Ws\,\},\qquad \Wx\subseteq\{1,\dots,n\},\quad
\Ws\subseteq\{1,\dots,m_I\} .
\]
A surplus $i\in\Ws$ is called \textbf{active}; a surplus $i\notin\Ws$ is called
\textbf{free}.
\end{definition}

Given a working set and a feasible iterate $y_k$ with $y_\ell = 0$ on $\Wset$, the
active-set method computes a search direction from the \textbf{equality-constrained
subproblem}
\begin{equation} \label{eq:eqp}
\min_{p}\;\frac{1}{2\delta}\,p^\top\Abar^\top\Abar\,p + g^\top p
\qquad\text{s.t.}\qquad \Cbar p = 0,\qquad p_\ell = 0\;\;(\ell\in\Wset),
\end{equation}
which minimizes $f(y_k+p)$ over the face of the feasible region defined by the working set.
To write its constraints compactly, let $W_x\in\R^{|\Wx|\times n}$ and
$W_s\in\R^{|\Ws|\times m_I}$ be the selection matrices whose rows are the unit vectors
$e_j^\top$, $j\in\Wx$, and $e_i^\top$, $i\in\Ws$.

\begin{definition}[Constraint matrix of a working set] \label{def:DW}
The \textbf{constraint matrix associated with $\Wset$} is
\begin{equation} \label{eq:DW}
D_\Wset = \begin{bmatrix} C & 0\\ W_x & 0\\ 0 & W_s\end{bmatrix}
\in\R^{(K+|\Wx|+|\Ws|)\times(n+m_I)} .
\end{equation}
The constraints of~\eqref{eq:eqp} are $D_\Wset\,p = 0$.
\end{definition}

The KKT conditions of~\eqref{eq:eqp} form the linear \textbf{KKT system}
\begin{equation} \label{eq:kkt-system}
\begin{bmatrix} \frac{1}{\delta}\Abar^\top\Abar & D_\Wset^\top\\ D_\Wset & 0\end{bmatrix}
\begin{bmatrix} p\\ -\mu\\ -\nu_x\\ -\nu_s\end{bmatrix}
= \begin{bmatrix} -g\\ 0\\ 0\\ 0\end{bmatrix},
\end{equation}
where $\mu\in\R^K$, $\nu_x\in\R^{|\Wx|}$ and $\nu_s\in\R^{|\Ws|}$ are the Lagrange
multipliers of the convexity constraints and of the working-set constraints of the columns
and of the surpluses. If $p = 0$ solves~\eqref{eq:kkt-system}, the iterate minimizes $f$
over its face, and we call it a \textbf{subspace minimizer}. The generic method is
\cref{alg:qpas}.

\begin{algobox}{Primal Active-Set Method for the Regularized Master Problem}{alg:qpas}
\State \textbf{Input:} feasible $y$, working set $\Wset$ with $y_\ell = 0$ on $\Wset$
\Repeat
    \State Solve the KKT system~\eqref{eq:kkt-system} for $p$, $\mu$, $\nu = (\nu_x,\nu_s)$
    \If{$p = 0$}
        \If{$\nu_\ell\ge-\varepsilon$ for all $\ell\in\Wset$}
            \State \textbf{stop}: $y$ is optimal
        \Else
            \State $r\gets\arg\min_{\ell\in\Wset}\nu_\ell$;\quad $\Wset\gets\Wset\setminus\{r\}$
                   \Comment{reduce the working set}
        \EndIf
    \Else
        \State $\alpha\gets\min\Bigl\{1,\;\min\limits_{\ell\notin\Wset,\;p_\ell<0}\;
               -\dfrac{y_\ell}{p_\ell}\Bigr\}$ \Comment{ratio test}
        \State $y\gets y+\alpha p$
        \If{$\alpha<1$}
            \State $q\gets$ the index attaining the minimum;\quad $\Wset\gets\Wset\cup\{q\}$
                   \Comment{expand the working set}
        \EndIf
    \EndIf
\Until{stopped}
\end{algobox}

Under the usual nondegeneracy assumptions, every iteration either decreases $f$ strictly or
changes the working set in a way that cannot cycle, so the method terminates after finitely
many iterations \cite{nocedal2006numerical, fletcher1971quadratic}. Its cost is dominated by
the solution of the KKT system~\eqref{eq:kkt-system}, which has dimension
$n + m_I + K + |\Wset|$ and grows with every generated column. The rest of this work shows
how to avoid ever forming it. The nullspace method of \cref{chap:nullspace} needs two
properties:

\begin{assumption} \label{ass:nullspace}
\begin{enumerate}
    \item $D_\Wset$ has full row rank;
    \item the reduced Hessian $Z^\top\Abar^\top\Abar Z$ is positive definite, where the columns
          of $Z$ form a basis of $\Null(D_\Wset)$.
\end{enumerate}
\end{assumption}

Property~(i) holds for every working set that keeps one column per block free
(\cref{thm:rank} below). Property~(ii) can fail, and \cref{chap:singular} shows how the
method guarantees it.


\section{The Constraint Matrix and the Reordering Theorem}
\label{sec:reordering}

The structure of $D_\Wset$ is determined by the convexity matrix $C$, which has the column
generation property, and by the zero columns of the surpluses, which do not. We first
arrange the columns of the weights $x$, and then add the surpluses.

\begin{definition}[Free columns] \label{def:free-columns}
The \textbf{free columns} are $\Fset_x = \{1,\dots,n\}\setminus\Wx$.
\end{definition}

\begin{lemma} \label{lem:free-block}
If $y$ is feasible and $x_j = 0$ for all $j\in\Wx$, then every block has at least one free
column, and the submatrix $C_{\Fset_x}$ of $C$ has the column generation property.
\end{lemma}

\begin{proof}
For every block $k$, $\sum_{\psi(j) = k}x_j = 1$ and $x_j = 0$ on $\Wx$, so some column of
block $k$ lies in $\Fset_x$. The claim follows from \cref{lem:cg-property}. 
\end{proof}

\begin{definition}[Representatives] \label{def:representatives}
A \textbf{list of representatives} is an ordered list $\Sset = [S_1,\dots,S_K]$ of free
columns with $\psi(S_k) = k$, one per block. The \textbf{representative of column} $j$ is
\[
\rep(j) = S_{\psi(j)} .
\]
The remaining free columns form the \textbf{nullspace set} $\Nx = \Fset_x\setminus\Sset$.
\end{definition}

The representatives exist by \cref{lem:free-block}. Their columns in $C$ are the unit vectors
$e_1,\dots,e_K$, so $C_\Sset = I_K$.

\begin{definition}[Extended ordering] \label{def:extended-ordering}
The \textbf{extended ordering} of the variables of~\eqref{eq:qp} is the partition into the
five sets
\begin{center}
\renewcommand{\arraystretch}{1.25}
\begin{tabular}{lll}
\toprule
set & contents & status \\
\midrule
$\Sset$ & one representative column per block & free \\
$\Wx$   & columns in the working set & $x_j = 0$ fixed \\
$\Nx$   & the remaining free columns & free \\
$\Ws$   & active surpluses & $s_i = 0$ fixed \\
$\Ns$   & free surpluses & free \\
\bottomrule
\end{tabular}
\end{center}
We write $[\,\Sset,\Wx,\Nx\mid\Ws,\Ns\,]$ for this ordering.
\end{definition}

\begin{theorem}[Reordering theorem] \label{thm:reordering}
Let $\Wset = \Wx\cup\Ws$ be a working set such that every block has a free column, and let
$[\,\Sset,\Wx,\Nx\mid\Ws,\Ns\,]$ be an extended ordering. Let $Q$ be the permutation matrix
that orders the variables in this way. Then
\begin{equation} \label{eq:reordered}
D_\Wset\,Q =
\begin{bmatrix}
I_K & C_{\Wx} & C_{\Nx} & 0 & 0\\
0 & I_{|\Wx|} & 0 & 0 & 0\\
0 & 0 & 0 & I_{|\Ws|} & 0
\end{bmatrix},
\end{equation}
where $C_{\Wx}$ and $C_{\Nx}$ are the columns of $C$ indexed by $\Wx$ and $\Nx$.
\end{theorem}

\begin{proof}
The permutation is block-diagonal, $Q = \diag(Q_x, Q_s)$: $Q_x$ orders the column weights as
$[\Sset,\Wx,\Nx]$ and $Q_s$ orders the surpluses as $[\Ws,\Ns]$. Apply it to the three block
rows of~\eqref{eq:DW}.
\begin{itemize}
    \item The convexity rows $[\,C\;\;0\,]$ become $[\,C_\Sset\;\;C_{\Wx}\;\;C_{\Nx}\;\;0\;\;0\,]$,
          and $C_\Sset = I_K$ by \cref{def:representatives}.
    \item The rows $[\,W_x\;\;0\,]$ select the columns in $\Wx$, which appear in the second
          position: $[\,0\;\;I_{|\Wx|}\;\;0\;\;0\;\;0\,]$.
    \item The rows $[\,0\;\;W_s\,]$ select the surpluses in $\Ws$, which appear in the fourth
          position: $[\,0\;\;0\;\;0\;\;I_{|\Ws|}\;\;0\,]$. 
\end{itemize}
\end{proof}

\begin{figure}[ht]
\centering
\begin{tikzpicture}[x=0.6cm, y=0.6cm, font=\small]
    % column group widths: S=2, Wx=6, Nx=3, Ws=2, Ns=2 ; rows: K=2, Wx=6, Ws=2
    \def\Sx{0}\def\Wxx{2}\def\Nxx{8}\def\Wsx{11}\def\Nsx{13}\def\Ed{15}
    \def\Tp{10}\def\rK{8}\def\rW{2}\def\rB{0}
    % shaded nonzero blocks
    \fill[cblue!35]   (\Sx,\rK)  rectangle (\Wxx,\Tp);
    \fill[cblue!12]   (\Wxx,\rK) rectangle (\Nxx,\Tp);
    \fill[cblue!12]   (\Nxx,\rK) rectangle (\Wsx,\Tp);
    \fill[corange!35] (\Wxx,\rW) rectangle (\Nxx,\rK);
    \fill[cgreen!35]  (\Wsx,\rB) rectangle (\Nsx,\rW);
    % identity diagonals
    \draw[very thick, cblue!80!black]   (\Sx,\Tp) -- (\Wxx,\rK);
    \draw[very thick, corange!80!black] (\Wxx,\rK) -- (\Nxx,\rW);
    \draw[very thick, cgreen!60!black]  (\Wsx,\rW) -- (\Nsx,\rB);
    % grid
    \draw[thick] (\Sx,\rB) rectangle (\Ed,\Tp);
    \foreach \c in {\Wxx,\Nxx,\Wsx,\Nsx} \draw (\c,\rB) -- (\c,\Tp);
    \draw[thick, dashed] (\Wsx,\rB) -- (\Wsx,\Tp);
    \foreach \r in {\rK,\rW} \draw (\Sx,\r) -- (\Ed,\r);
    % labels inside
    \node at (1.4,9.4) {\scriptsize $I_K$};
    \node at (5,9) {$C_{\Wx}$};
    \node at (9.5,9) {$C_{\Nx}$};
    \node at (12,9) {$0$}; \node at (14,9) {$0$};
    \node[fill=white, inner sep=1pt] at (5.6,4.6) {$I_{|\Wx|}$};
    \node at (12,1) [right=-2pt] {};
    \node[fill=white, inner sep=1pt] at (12.3,0.6) {\scriptsize $I_{|\Ws|}$};
    % column headers
    \node at (1,10.6) {$\Sset$};
    \node at (5,10.6) {$\Wx$};
    \node at (9.5,10.6) {$\Nx$};
    \node at (12,10.6) {$\Ws$};
    \node at (14,10.6) {$\Ns$};
    \draw[decorate, decoration={brace, amplitude=4pt}] (0,11.2) -- (11,11.2)
        node[midway, above=4pt] {column weights $x$};
    \draw[decorate, decoration={brace, amplitude=4pt}] (11,11.2) -- (15,11.2)
        node[midway, above=4pt] {surpluses $s$};
    % row labels
    \node[left] at (0,9) {$C$};
    \node[left] at (0,5) {$W_x$};
    \node[left] at (0,1) {$W_s$};
\end{tikzpicture}
\caption{Block structure of the reordered constraint matrix $D_\Wset Q$ of
\cref{thm:reordering}. The diagonal lines are identity blocks. The surplus part (right of the
dashed line) only contributes the identity block of the active surpluses; the free
surpluses $\Ns$ have empty columns.}
\label{fig:reordered}
\end{figure}

\begin{theorem}[Full row rank] \label{thm:rank}
Under the assumptions of \cref{thm:reordering}, $D_\Wset$ has full row rank,
\[
\rank(D_\Wset) = K + |\Wx| + |\Ws|,
\]
and $\dim\Null(D_\Wset) = |\Nx| + |\Ns|$.
\end{theorem}

\begin{proof}
In~\eqref{eq:reordered}, the identity blocks $I_K$, $I_{|\Wx|}$ and $I_{|\Ws|}$ lie in
disjoint rows and disjoint columns, so the rows are linearly independent. Permutations do not
change the rank. The dimension of the nullspace follows from the rank--nullity theorem:
$(n+m_I) - (K+|\Wx|+|\Ws|) = |\Nx| + |\Ns|$, because $n = K + |\Wx| + |\Nx|$ and
$m_I = |\Ws| + |\Ns|$. 
\end{proof}

\Cref{thm:rank} establishes property~(i) of \cref{ass:nullspace} for every working set the
method can produce: the ratio test never removes the last free column of a block, since the
representative of the block can only reach zero while some other column of the block grows
(\cref{lem:replacement} below).

\begin{remark}[Comparison with problems without inequality rows] \label{rem:three-sets}
Without inequality rows ($m_I = 0$), the surplus sets are empty, and the extended ordering
reduces to the three sets $[\Sset,\Wx,\Nx]$. The inequality rows thus only add two groups of
columns: active surpluses, which behave like working-set columns, and free surpluses, which
behave like nullspace columns without a representative. The loss of the column generation
property (\cref{rem:cg-property-lost}) does not affect the reordering. It only means that the
surpluses form groups of their own.
\end{remark}

\Cref{fig:spy-DW} shows the effect of the reordering on an actual constraint matrix, taken
from the final working set of a column generation run on one of the test instances of
\cref{chap:results}.

\begin{figure}[ht]
\centering
\IfFileExists{data/spy_DW_orig.csv}{%
\begin{tikzpicture}
\begin{groupplot}[group style={group size=2 by 1, horizontal sep=1.2cm},
    width=0.5\textwidth, height=0.42\textwidth, y dir=reverse,
    xmin=0, ymin=0, enlargelimits=false, tick label style={font=\scriptsize},
    title style={font=\small}, xlabel style={font=\small}, ylabel style={font=\small}]
\nextgroupplot[title={original order}, xlabel={variable}, ylabel={row of $D_\Wset$}]
    \addplot[only marks, mark=*, mark size=0.25pt, cblue]
        table[x=col, y=row, col sep=comma] {data/spy_DW_orig.csv};
\nextgroupplot[title={order $[\Sset,\Wx,\Nx\mid\Ws,\Ns]$}, xlabel={variable}]
    \addplot[only marks, mark=*, mark size=0.25pt, cblue]
        table[x=col, y=row, col sep=comma] {data/spy_DW_perm.csv};
\end{groupplot}
\end{tikzpicture}}{\fbox{spy plot data missing: run scripts/make\_plot\_data.jl}}
\caption{Sparsity pattern of the constraint matrix $D_\Wset$ at the end of a column
generation run (instance I2 with 30\% inequality rows), before and after the reordering of
\cref{thm:reordering}. After reordering, the identity blocks of the representatives (top
left), the working-set columns (diagonal) and the active surpluses (bottom right) are
visible.}
\label{fig:spy-DW}
\end{figure}

\chapter{The Nullspace Method}
\label{chap:nullspace}

\vspace*{\fill}

The nullspace method solves the KKT system~\eqref{eq:kkt-system} in two smaller steps. The
search direction is written as $p = Z\eta$, where the columns of $Z$ span the nullspace of
the constraint matrix. This turns the KKT system into a \textbf{reduced system} with the
reduced Hessian, whose dimension equals the number of degrees of freedom of the current face.
The Lagrange multipliers are then recovered from the remaining equations.

\medskip

For the regularized master problem both steps can be carried out explicitly. The reordering
theorem of the previous chapter gives a nullspace basis that consists of the vectors
$e_j - e_{\rep(j)}$ and of unit vectors for the free surpluses. This basis is
\textbf{never stored}: it is determined by the five sets alone.

\medskip

The central result of this chapter is a \textbf{decoupling} of the reduced system. The
surplus part of the direction is given by an explicit formula, and only an
$|\Nx|\times|\Nx|$ system has to be solved. Its matrix is the \textbf{effective reduced
Hessian} $\Heff = (\Aeff Z_x)^\top(\Aeff Z_x)$, the reduced Hessian of the equality rows
together with the inequality rows whose surplus is active. Finally, the multipliers turn out
to be the reduced costs of the columns and the duals of the inequality rows, available at
almost no cost.

\vspace*{\fill}

\newpage


\section{The Nullspace Basis}
\label{sec:basis}

By \cref{thm:rank}, the nullspace of $D_\Wset$ has dimension $|\Nx|+|\Ns|$. We label its basis
vectors by the indices of these two sets.

\begin{theorem}[Nullspace basis] \label{thm:basis}
Under the assumptions of \cref{thm:reordering}, the matrix
\begin{equation} \label{eq:Z-reordered}
Z = Q\tilde Z,\qquad
\tilde Z = \begin{bmatrix}
-C_{\Nx} & 0\\
0 & 0\\
I_{|\Nx|} & 0\\
0 & 0\\
0 & I_{|\Ns|}
\end{bmatrix}
\begin{matrix} \leftarrow\Sset\\ \leftarrow\Wx\\ \leftarrow\Nx\\ \leftarrow\Ws\\ \leftarrow\Ns\end{matrix}
\end{equation}
is a basis of $\Null(D_\Wset)$. In the original order of the variables it is block-diagonal,
\begin{equation} \label{eq:Z-blocks}
Z = \begin{bmatrix} Z_x & 0\\ 0 & Z_s\end{bmatrix},
\qquad Z_x\in\R^{n\times|\Nx|},\quad Z_s\in\R^{m_I\times|\Ns|},
\end{equation}
where $Z_x$ is a nullspace basis for the column weights and $Z_s$ is the selection matrix of
the free surpluses.
\end{theorem}

\begin{proof}
Write $\tilde Z = Q^\top Z$ in the five row blocks of the extended ordering,
\[
\tilde Z = \bigl[\,Z_\Sset;\; Z_{\Wx};\; Z_{\Nx};\; Z_{\Ws};\; Z_{\Ns}\,\bigr].
\]
By~\eqref{eq:reordered},
$D_\Wset Z = (D_\Wset Q)\tilde Z = 0$ is equivalent to
\[
Z_\Sset + C_{\Wx}Z_{\Wx} + C_{\Nx}Z_{\Nx} = 0,\qquad Z_{\Wx} = 0,\qquad Z_{\Ws} = 0 .
\]
The blocks $Z_{\Nx}$ and $Z_{\Ns}$ are free, so the nullspace has the stated dimension, and the
choice $Z_{\Nx} = I$, $Z_{\Ns} = I$ (with the corresponding zero columns) gives a basis with
$Z_\Sset = -C_{\Nx}$ in the first group of columns and $Z_\Sset = 0$ in the second. The column
weights and the surpluses do not mix, which gives~\eqref{eq:Z-blocks}. 
\end{proof}

\begin{cor}[Explicit nullspace vectors] \label{cor:z-explicit}
The columns of $Z$ are
\begin{equation} \label{eq:z-explicit}
z_j = e_j - e_{\rep(j)}\quad (j\in\Nx),\qquad\qquad z_i = e_{n+i}\quad (i\in\Ns),
\end{equation}
where $e_\ell$ denotes the $\ell$-th unit vector in $\R^{n+m_I}$.
\end{cor}

\begin{proof}
Column $j\in\Nx$ of $\tilde Z$ has a $1$ in the position of $j$ and the entries of
$-C_{\Nx}e_j = -e_{\psi(j)}$ in the representative block. By the column generation property,
this is a single $-1$ in the row of block $\psi(j)$, i.e.\ at the position of the
representative $S_{\psi(j)} = \rep(j)$. Column $i\in\Ns$ has a single $1$ at the position of
surplus $i$. 
\end{proof}

\begin{figure}[ht]
\centering
\begin{tikzpicture}[x=0.75cm, y=1cm, font=\small]
    % one block with 5 columns: rep = 2, N = {1,4}, W = {3,5}
    \foreach \j/\h/\lab/\col in {1/0.25/j_1/cblue, 2/0.5/{\rep}/cgreen, 3/0/w/cgrey, 4/0.25/j_2/cblue, 5/0/w/cgrey} {
        \fill[\col!45] (\j-0.35,0) rectangle (\j+0.35,{\h*3});
        \draw[\col!80!black] (\j-0.35,0) rectangle (\j+0.35,{\h*3});
    }
    \node[below] at (1,0) {$x_{j_1}$};
    \node[below] at (2,0) {$x_{\rep}$};
    \node[below] at (3,0) {$x_{w}$};
    \node[below] at (4,0) {$x_{j_2}$};
    \node[below] at (5,0) {$x_{w'}$};
    \draw (0.4,0) -- (5.6,0);
    % arrows for z_{j1}: +1 at j1, -1 at rep
    \draw[-{Stealth[length=2.5mm]}, thick, cred] (1,0.85) -- (1,1.45) node[above] {$+\eta_{j_1}$};
    \draw[-{Stealth[length=2.5mm]}, thick, cred] (2,1.6) -- (2,1.0) node[pos=0, above] {$-\eta_{j_1}$};
    \node[align=left, right] at (6.2,1.1) {$x + \eta_{j_1} z_{j_1}$:\\ weight moves from the\\ representative to $j_1$;\\ the block sum stays $1$,\\ and $x_w = x_{w'} = 0$ stay.};
    % surplus
    \begin{scope}[xshift=11.5cm]
        \fill[corange!45] (-0.35,0) rectangle (0.35,0.9);
        \draw[corange!80!black] (-0.35,0) rectangle (0.35,0.9);
        \node[below] at (0,0) {$s_i$};
        \draw (-0.6,0) -- (0.6,0);
        \draw[-{Stealth[length=2.5mm]}, thick, cred] (0,1.0) -- (0,1.6) node[above] {$+\eta_i$};
        \node[align=left, right] at (0.8,0.8) {$z_i = e_{n+i}$:\\ a free surplus\\ moves on its own};
    \end{scope}
\end{tikzpicture}
\caption{Interpretation of the nullspace vectors~\eqref{eq:z-explicit}. A vector
$z_j = e_j - e_{\rep(j)}$ transfers weight between a nullspace column and the
representative of its block, which keeps the convexity constraint satisfied and the
working-set columns at zero. A free surplus belongs to no block and moves independently.}
\label{fig:nullspace-vectors}
\end{figure}

\Cref{cor:z-explicit} shows that $Z$ is completely determined by the extended ordering: by
the sets $\Nx$, $\Ns$ and the representative mapping. It therefore never has to be formed.
The same holds for the product of the linking matrix with the nullspace vectors,
\begin{equation} \label{eq:AZ-columns}
\Abar z_j = \begin{bmatrix} a_j^E - a_{\rep(j)}^E\\ a_j^I - a_{\rep(j)}^I\end{bmatrix}
= a_j - a_{\rep(j)} =: u_j\quad(j\in\Nx),\qquad
\Abar z_i = -e_i^{(I)}\quad(i\in\Ns),
\end{equation}
where $e_i^{(I)}\in\R^m$ is the unit vector of inequality row $i$.


\section{The Reduced KKT System}
\label{sec:reduced}

\subsection{The reduced system}

Every direction that satisfies the constraints of~\eqref{eq:eqp} has the form
$p = Z\eta$ with $\eta = [\eta_x;\eta_s]\in\R^{|\Nx|+|\Ns|}$. Substituting this into the first
block row of~\eqref{eq:kkt-system} and multiplying by $Z^\top$, which annihilates
$D_\Wset^\top$, gives the \textbf{reduced KKT system}
\begin{equation} \label{eq:reduced-system}
\frac{1}{\delta}\,Z^\top\Abar^\top\Abar Z\,\eta = -Z^\top g .
\end{equation}
We call $\Hbar = Z^\top\Abar^\top\Abar Z$ the (unscaled) \textbf{reduced Hessian}. The vector
$\eta$ minimizes $f(y_k + Z\eta) = f(y_k) + (Z^\top g)^\top\eta + \frac{1}{2\delta}\eta^\top\Hbar\eta$,
so~\eqref{eq:reduced-system} is the subproblem~\eqref{eq:eqp} expressed in the coordinates
of the nullspace. The parameter $\delta$ only appears on the right-hand side if we write the
system as $\Hbar\eta = -\delta Z^\top g$.

\subsection{Block structure}

To exhibit the structure of $\Hbar$, we introduce the products of the linking matrices with
the column part of the basis,
\[
B_E := A_E Z_x\in\R^{m_E\times|\Nx|},\qquad B_I := A_I Z_x\in\R^{m_I\times|\Nx|},
\]
and split the rows of $B_I$ and of $\pi_I$ according to the status of the surplus of each row:
\[
B_\Wset := W_s B_I,\quad \pi_\Wset := W_s\pi_I \qquad\text{(rows with an active surplus)},
\]
\[
B_\mathcal{N} := Z_s^\top B_I,\quad \pi_\mathcal{N} := Z_s^\top\pi_I \qquad\text{(rows with a free surplus)}.
\]
Since every inequality row belongs to exactly one of the two groups, $W_s^\top W_s + Z_sZ_s^\top = I$
and therefore
\begin{equation} \label{eq:BI-split}
B_I^\top B_I = B_\Wset^\top B_\Wset + B_\mathcal{N}^\top B_\mathcal{N} .
\end{equation}

\begin{lemma}[Blocks of the reduced system] \label{lem:reduced-blocks}
With~\eqref{eq:Z-blocks},
\[
\Abar Z = \begin{bmatrix} B_E & 0\\ B_I & -Z_s\end{bmatrix},\qquad
\Hbar = \begin{bmatrix} H & -B_\mathcal{N}^\top\\ -B_\mathcal{N} & I\end{bmatrix},\qquad
Z^\top g = \begin{bmatrix} Z_x^\top g_x\\ \pi_\mathcal{N}\end{bmatrix},
\]
where $H = B_E^\top B_E + B_I^\top B_I = (AZ_x)^\top(AZ_x)$ is the reduced Hessian of the column
weights with all linking rows.
\end{lemma}

\begin{proof}
The first identity follows from the block forms of $\Abar$ and $Z$. Then
$\Hbar = (\Abar Z)^\top(\Abar Z)$ has the blocks $B_E^\top B_E + B_I^\top B_I$,
$-B_I^\top Z_s = -B_\mathcal{N}^\top$ and $Z_s^\top Z_s = I$. Finally
$Z^\top g = [Z_x^\top g_x;\, Z_s^\top g_s]$ and $g_s = \pi_I$ by~\eqref{eq:gradient-split}.

\end{proof}

\subsection{Decoupling and the effective reduced Hessian}

The identity block in the lower right corner of $\Hbar$ allows us to eliminate the surplus
part of the direction.

\begin{theorem}[Decoupling of the reduced system] \label{thm:decoupling}
The reduced system~\eqref{eq:reduced-system} is equivalent to
\begin{align}
\Heff\,\eta_x &= -\delta\,Z_x^\top\bigl(c - A_E^\top\pi_E - A_I^\top W_s^\top\pi_\Wset\bigr),
\label{eq:eff-system}\\[3pt]
\eta_s &= B_\mathcal{N}\,\eta_x - \delta\,\pi_\mathcal{N},
\label{eq:surplus-direction}
\end{align}
with the \textbf{effective reduced Hessian}
\begin{equation} \label{eq:Heff}
\Heff := H - B_\mathcal{N}^\top B_\mathcal{N} = B_E^\top B_E + B_\Wset^\top B_\Wset .
\end{equation}
\end{theorem}

\begin{proof}
Multiply~\eqref{eq:reduced-system} by $\delta$ and write it in the blocks of
\cref{lem:reduced-blocks}:
\[
H\eta_x - B_\mathcal{N}^\top\eta_s = -\delta Z_x^\top g_x,\qquad
-B_\mathcal{N}\eta_x + \eta_s = -\delta\pi_\mathcal{N}.
\]
The second equation is~\eqref{eq:surplus-direction}. Substituting it into the first gives
$(H - B_\mathcal{N}^\top B_\mathcal{N})\eta_x = -\delta\bigl(Z_x^\top g_x + B_\mathcal{N}^\top\pi_\mathcal{N}\bigr)$.
The second expression for $\Heff$ follows from~\eqref{eq:BI-split}. For the right-hand side,
insert $g_x = c - A_E^\top\pi_E - A_I^\top\pi_I$ and split
$Z_x^\top A_I^\top\pi_I = B_I^\top\pi_I = B_\Wset^\top\pi_\Wset + B_\mathcal{N}^\top\pi_\mathcal{N}$.
The terms with $\pi_\mathcal{N}$ cancel, and $B_\Wset^\top\pi_\Wset = Z_x^\top A_I^\top W_s^\top\pi_\Wset$.

\end{proof}

\begin{figure}[ht]
\centering
\begin{tikzpicture}[x=0.5cm, y=0.5cm, font=\small]
    % H bar
    \fill[cblue!25] (0,3) rectangle (6,9);
    \fill[corange!30] (6,3) rectangle (9,9);
    \fill[corange!30] (0,0) rectangle (6,3);
    \fill[cgreen!25] (6,0) rectangle (9,3);
    \draw[thick] (0,0) rectangle (9,9);
    \draw (6,0) -- (6,9); \draw (0,3) -- (9,3);
    \draw[cgreen!60!black, very thick] (6,3) -- (9,0);
    \node at (3,6) {$H$};
    \node at (7.5,6) {$-B_\mathcal{N}^\top$};
    \node at (3,1.5) {$-B_\mathcal{N}$};
    \node[fill=cgreen!25, inner sep=1pt] at (7.9,2.1) {$I$};
    \node at (4.5,10) {$\Hbar$};
    \node[left] at (0,6) {$\Nx$}; \node[left] at (0,1.5) {$\Ns$};
    % arrow
    \draw[-{Stealth[length=3mm]}, thick] (10,4.5) -- (14,4.5)
        node[midway, above] {eliminate $\eta_s$};
    % H eff
    \begin{scope}[xshift=7.5cm]
    \fill[cblue!25] (0,1.5) rectangle (6,7.5);
    \draw[thick] (0,1.5) rectangle (6,7.5);
    \node[align=center] at (3,4.5) {$\Heff =$\\ $H - B_\mathcal{N}^\top B_\mathcal{N}$};
    \node at (3,8.5) {$|\Nx|\times|\Nx|$};
    \end{scope}
\end{tikzpicture}
\caption{The reduced Hessian $\Hbar$ of \cref{lem:reduced-blocks} and its Schur complement.
Eliminating the identity block of the free surpluses leaves the effective reduced Hessian
$\Heff$ on the nullspace columns only.}
\label{fig:schur}
\end{figure}

$\Heff$ is the Schur complement of the identity block in $\Hbar$
(\cref{fig:schur}). Only an $|\Nx|\times|\Nx|$ system has to be solved; the surplus
part~\eqref{eq:surplus-direction} costs one sparse matrix--vector product. The following
corollary gives the interpretation that the rest of this work builds on.

\begin{cor}[Effective linking matrix] \label{cor:Aeff}
Define the \textbf{effective linking matrix} and the corresponding duals
\[
\Aeff = \begin{bmatrix} A_E\\ W_sA_I\end{bmatrix},\qquad
\pieff = \begin{bmatrix}\pi_E\\ \pi_\Wset\end{bmatrix},
\]
i.e.\ the equality rows together with the inequality rows whose surplus is active. Then
\begin{equation} \label{eq:Heff-Aeff}
\Heff = (\Aeff Z_x)^\top(\Aeff Z_x),\qquad
\Heff\,\eta_x = -\delta\,Z_x^\top\bigl(c - \Aeff^\top\pieff\bigr).
\end{equation}
This is the reduced system of the master problem with equality linking rows only, whose
linking rows are the rows of $\Aeff$.
\end{cor}

\begin{proof}
$\Aeff Z_x = [B_E;\, B_\Wset]$, which gives the first identity by~\eqref{eq:Heff}. The second is
the right-hand side of~\eqref{eq:eff-system}, since
$\Aeff^\top\pieff = A_E^\top\pi_E + A_I^\top W_s^\top\pi_\Wset$. 
\end{proof}

\begin{remark}[Interpretation] \label{rem:Aeff-interpretation}
An inequality row with a \emph{free} surplus does not restrict the column weights: any
change of its residual is absorbed by the surplus, so the row contributes neither curvature
nor dual pressure to the column part of the problem. An inequality row with an \emph{active}
surplus behaves like an equality row. The column part of the method is therefore the
method for equality rows, applied to a linking matrix $\Aeff$ whose rows change with the
surplus working set. Without inequality rows, or with all surpluses active, $\Aeff = A$ and
$\Heff = H$.
\end{remark}

\begin{cor}[Positive definiteness and dimension] \label{cor:dimension}
$\Hbar$ is positive definite if and only if $\Heff$ is, i.e.\ if and only if $\Aeff Z_x$ has
full column rank. In this case
\[
|\Nx|\;\le\;\rank(\Aeff)\;\le\;m_E + |\Ws|,\qquad\text{and hence}\qquad |\Nx| + |\Ns|\le m .
\]
\end{cor}

\begin{proof}
A symmetric block matrix with a positive definite diagonal block is positive definite if
and only if the Schur complement of that block is (Haynsworth inertia additivity). The
rank bound holds because $\Heff = (\Aeff Z_x)^\top(\Aeff Z_x)$ is nonsingular only if
$|\Nx|\le\rank(\Aeff Z_x)\le\rank(\Aeff)$, and $\Aeff$ has $m_E+|\Ws|$ rows. Adding
$|\Ns| = m_I - |\Ws|$ gives the last bound. 
\end{proof}

\Cref{cor:dimension} is the reason the method scales: the system that has to be factorized
never exceeds the number of linking rows, \emph{however many columns are generated}.

\subsection{The search direction}

With $\eta$ from \cref{thm:decoupling}, the direction $p = Z\eta$ has the components
\begin{equation} \label{eq:direction-components}
p_j = \eta_j\;(j\in\Nx),\qquad
p_{S_k} = -\!\!\sum_{j\in\Nx,\,\psi(j) = k}\!\!\eta_j\;(k = 1,\dots,K),\qquad
(p_s)_i = (A_Ip_x)_i - \delta\,\pi_i\;(i\in\Ns),
\end{equation}
and zero on $\Wx$ and $\Ws$. The surplus component uses $B_\mathcal{N}\eta_x = Z_s^\top A_I Z_x\eta_x = Z_s^\top A_I p_x$.
The direction is supported on $\Sset\cup\Nx\cup\Ns$, and its column part and the vector
$A_Ip_x$ only involve the columns in $\Sset\cup\Nx$.

\begin{prop}[Full step] \label{prop:full-step}
If the ratio test returns $\alpha = 1$, the new iterate $y_k+p$ is a subspace minimizer, i.e.\
the next reduced gradient $Z^\top\nabla f(y_k+p)$ vanishes.
\end{prop}

\begin{proof}
On the affine set $\{y_k + Z\eta\}$, $f$ is the convex quadratic
$\phi(\eta) = f(y_k) + (Z^\top g)^\top\eta + \frac{1}{2\delta}\eta^\top\Hbar\eta$, whose
minimizer is the solution of~\eqref{eq:reduced-system}. 
\end{proof}

After a full step the method can therefore skip the next direction computation and test the
multipliers directly. The next proposition shows how the free surpluses act on the duals.

\begin{prop}[Duals of free-surplus rows] \label{prop:free-duals}
A step $y\gets y + \alpha p$ changes the dual of every row with a free surplus as
\[
\pi_i\;\gets\;(1-\alpha)\,\pi_i\qquad(i\in\Ns).
\]
In particular $\pi_i = 0$ on these rows at every subspace minimizer.
\end{prop}

\begin{proof}
By~\eqref{eq:pi-of-y}, $\pi$ changes by $-\frac{\alpha}{\delta}\Abar p$. For row $i\in\Ns$,
$(\Abar p)_i = (A_Ip_x)_i - (p_s)_i = \delta\pi_i$ by~\eqref{eq:direction-components}. 
\end{proof}

\begin{remark}[Independence of $\delta$ and $\pihat$] \label{rem:factor-independent}
$\Heff$ depends only on the columns and the working set, not on $\delta$, $\pihat$ or the
iterate. A change of the stability center or of $\delta$ only changes the right-hand side
of~\eqref{eq:eff-system}. This is what makes the hot start of \cref{sec:hot-start} possible.
\end{remark}

\subsection{Implicit evaluation}

Neither $Z$, nor $\Aeff Z_x$, nor the gradient has to be stored. With $u_j = a_j - a_{\rep(j)}$
from~\eqref{eq:AZ-columns} and the diagonal \textbf{row mask} $D_\rho\in\R^{m\times m}$ with
$\rho_\ell = 0$ for the rows of free surpluses and $\rho_\ell = 1$ otherwise:
\begin{equation} \label{eq:implicit}
(\Heff)_{kj} = u_k^\top D_\rho\,u_j,\qquad
\bigl(Z_x^\top(c - \Aeff^\top\pieff)\bigr)_j = \bigl(c_j - a_j^\top D_\rho\pi\bigr)
- \bigl(c_{\rep(j)} - a_{\rep(j)}^\top D_\rho\pi\bigr).
\end{equation}
Both are computed from the sparse columns $a_j$ and the $m$-vector $\pi$. The cost of one entry
of the right-hand side is proportional to the number of nonzeros of two columns, independent
of the total number of columns $n$.


\section{The Lagrange Multipliers}
\label{sec:multipliers}

The multipliers follow from the first block row of the KKT system~\eqref{eq:kkt-system}:
\begin{equation} \label{eq:multiplier-system}
D_\Wset^\top\begin{bmatrix}\mu\\ \nu_x\\ \nu_s\end{bmatrix} = h,\qquad
h := \frac{1}{\delta}\Abar^\top\Abar\,p + g .
\end{equation}

\begin{theorem}[Lagrange multipliers] \label{thm:multipliers}
With the extended ordering of \cref{thm:reordering} and $h = [h_\Sset; h_{\Wx}; h_{\Nx}; h_{\Ws}; h_{\Ns}]$,
the solution of~\eqref{eq:multiplier-system} is
\begin{equation} \label{eq:multipliers}
\mu = h_\Sset,\qquad \nu_x = h_{\Wx} - C_{\Wx}^\top\mu,\qquad \nu_s = h_{\Ws},
\end{equation}
and moreover $h_{\Nx} = C_{\Nx}^\top\mu$ and $h_{\Ns} = 0$.
\end{theorem}

\begin{proof}
Multiply~\eqref{eq:multiplier-system} by $Q^\top$ and use the transpose of~\eqref{eq:reordered}:
\[
\begin{bmatrix}
I_K & 0 & 0\\
C_{\Wx}^\top & I & 0\\
C_{\Nx}^\top & 0 & 0\\
0 & 0 & I\\
0 & 0 & 0
\end{bmatrix}
\begin{bmatrix}\mu\\ \nu_x\\ \nu_s\end{bmatrix}
=
\begin{bmatrix} h_\Sset\\ h_{\Wx}\\ h_{\Nx}\\ h_{\Ws}\\ h_{\Ns}\end{bmatrix}.
\]
The first, second and fourth block rows give~\eqref{eq:multipliers}; the third and fifth
are the remaining identities. 
\end{proof}

The multipliers are only needed at a subspace minimizer, where $p = 0$ and therefore $h = g$.
In this case \cref{thm:multipliers} takes a particularly simple form.

\begin{cor}[Multipliers at a subspace minimizer] \label{cor:multipliers-minimizer}
At a subspace minimizer, with $\pi = \pi(y)$,
\begin{align*}
\mu_k &= c_{S_k} - a_{S_k}^\top\pi && (k = 1,\dots,K),\\
(\nu_x)_j &= c_j - a_j^\top\pi - \mu_{\psi(j)} = \bar c_j(\pi,\mu) && (j\in\Wx),\\
(\nu_s)_i &= \pi_i && (i\in\Ws),
\end{align*}
and $\bar c_j(\pi,\mu) = 0$ for $j\in\Sset\cup\Nx$, $\pi_i = 0$ for $i\in\Ns$.
\end{cor}

\begin{proof}
Insert $h = g$ and~\eqref{eq:gradient-split} into \cref{thm:multipliers}. For $j\in\Nx$,
$h_{\Nx} = C_{\Nx}^\top\mu$ reads $g_j = \mu_{\psi(j)}$, i.e.\ $\bar c_j = 0$, and $h_{\Ns} = 0$
reads $\pi_i = 0$. 
\end{proof}

\Cref{cor:multipliers-minimizer} gives a concrete form of \cref{rem:same-test}: the optimality
test of the active-set method checks the reduced costs of the working-set columns and the
signs of the duals of the inequality rows with active surpluses. All of them are computed
from the single vector $\pi$. The test touches every column of $\Wx$ once and costs
$O(\nnz(A))$; it is the only operation of the method whose cost grows with the number of
generated columns.

\begin{remark}[``$\le$'' rows] \label{rem:kappa-summary}
For a ``$\le$'' row with slack coefficient $\kappa_i = +1$ (\cref{rem:kappa}), the formulas
above change as follows; everything else, in particular $\Heff$ and $\Aeff$, is unchanged
because $\kappa_i^2 = 1$:
\[
\pi_I = \pihat_I + \tfrac{1}{\delta}\bigl(b_I - A_Ix - K_s s\bigr),\qquad
g_{n+i} = -\kappa_i\pi_i,\qquad
(\nu_s)_i = -\kappa_i\pi_i,\qquad
(p_s)_i = \kappa_i\bigl(\delta\pi_i - (A_Ip_x)_i\bigr).
\]
For $\kappa_i = -1$ these are the formulas of the text.
\end{remark}


\section{Summary of One Iteration}
\label{sec:iteration-summary}

\Cref{alg:nullspace-iteration} collects the computations of one iteration of the active-set
method in nullspace form, and \cref{tab:costs} lists their costs. Here $\bar\nu$ denotes the
average number of nonzeros of a column $a_j$.

\begin{algobox}{One Iteration of the Active-Set Method in Nullspace Form}{alg:nullspace-iteration}
\State \textbf{Given:} iterate $y = [x;s]$, the extended ordering, $\pi = \pi(y)$, the factor $\Heff = LDL^\top$
\State Compute $\mu^{\rho}_k = c_{S_k} - a_{S_k}^\top D_\rho\pi$ and the right-hand side
       $r_j = -\delta\bigl(c_j - a_j^\top D_\rho\pi - \mu^{\rho}_{\psi(j)}\bigr)$, $j\in\Nx$
\State Solve $LDL^\top\eta_x = r$
\State Expand $p_x = Z_x\eta_x$ and compute $(p_s)_i = (A_Ip_x)_i - \delta\pi_i$, $i\in\Ns$
\If{$p\approx0$}
    \State $\mu_k = c_{S_k} - a_{S_k}^\top\pi$; \;$\nu_x = $ reduced costs on $\Wx$; \;$\nu_s = \pi$ on $\Ws$
    \State Stop if all multipliers are $\ge-\varepsilon$; otherwise reduce the working set
           (\cref{sec:reducing})
\Else
    \State Ratio test over $\Sset\cup\Nx\cup\Ns$; \;$y\gets y+\alpha p$; \;$\pi\gets\pi - \frac{\alpha}{\delta}\Abar p$
    \State If blocked, expand the working set (\cref{sec:expanding})
\EndIf
\end{algobox}

\begin{table}[ht]
\centering
\renewcommand{\arraystretch}{1.25}
\begin{tabular}{lll}
\toprule
operation & cost & grows with $n$? \\
\midrule
right-hand side of~\eqref{eq:eff-system} & $O\bigl((|\Nx|+K)\,\bar\nu\bigr)$ & no \\
solve with the factor of $\Heff$ & $O(\nnz(L))$, at most $O(|\Nx|^2)$ & no \\
surplus direction~\eqref{eq:surplus-direction} & $O\bigl((|\Nx|+K)\,\bar\nu + m_I\bigr)$ & no \\
ratio test, step and update of $\pi$ & $O\bigl((|\Nx|+K)\,\bar\nu + m\bigr)$ & no \\
multipliers (\cref{cor:multipliers-minimizer}) & $O(\nnz(A) + m_I)$ & yes \\
\bottomrule
\end{tabular}
\caption{Cost of the operations of one iteration. Only the multiplier test touches all
columns.}
\label{tab:costs}
\end{table}

\chapter{Efficient Updates of the KKT System}
\label{chap:updates}

\vspace*{\fill}

Every iteration of the active-set method changes the working set by exactly one index: a
blocking variable is added, or a variable with a negative multiplier is released. The
previous chapter showed that everything the method needs is determined by the extended
ordering and by a factorization of the effective reduced Hessian $\Heff$. This chapter shows
that a change of the working set changes $\Heff$ only by a \textbf{low-rank modification}, so
that its $LDL^\top$ factorization can be updated in place instead of being recomputed.

\medskip

We first introduce a \textbf{fixed-length representation}, in which the factorization has a
fixed dimension and inactive indices are represented by identity rows. We then derive the
updates for the three ways of \textbf{expanding} the working set and the two ways of
\textbf{reducing} it, and translate them into row additions, row deletions and rank-one
modifications of the factorization. Row deletions leave structural zeros in the factor; we
analyse this effect and the \textbf{periodic refactorization} that controls it.

\medskip

Finally, we show that the state of the method can be carried over from one column generation
iteration to the next. A new column and a new stability center leave the factorization
unchanged, which gives a \textbf{hot start} of every master problem. We also analyse the
center-only iterations of \cref{sec:center-update} in this light.

\vspace*{\fill}

\newpage


\section{Fixed-Length Representation}
\label{sec:fixed-length}

The number of nullspace columns $|\Nx|$ changes in almost every iteration, and the number of
columns $n$ grows during column generation. Resizing a sparse factorization at every change
would be expensive and complicated. Instead, the factorized matrix is given a fixed dimension.

\begin{definition}[Fixed-length representation] \label{def:fixed-length}
Let $\nmax\ge n$ be a fixed \textbf{column capacity}, the maximal number of columns that column
generation may generate. The \textbf{embedded effective reduced Hessian} is the
$\nmax\times\nmax$ matrix
\[
\bigl(\Heff^{\mathrm{emb}}\bigr)_{kj} =
\begin{cases}
(\Heff)_{kj} = u_k^\top M_\rho u_j, & k,j\in\Nx,\\
1, & k = j\notin\Nx,\\
0, & \text{otherwise,}
\end{cases}
\]
with $u_j = a_j - a_{\rep(j)}$ and the row mask $M_\rho$ of~\eqref{eq:implicit}. Index $j$ of
the embedded matrix is column $j$ of the master problem. Indices of columns that are not in
$\Nx$, including columns that have not been generated yet, carry an identity row and column.
\end{definition}

In this representation, a column entering $\Nx$ replaces an identity row by a row of $\Heff$,
and a column leaving $\Nx$ is replaced by an identity row (\cref{fig:fixed-length}). The
dimension never changes. The surpluses do not appear at all: by \cref{thm:decoupling} they
have no index in the factorized matrix, so the inequality rows do not increase the size of
the factorization.

\begin{figure}[ht]
\centering
\begin{tikzpicture}[x=0.32cm, y=0.32cm]
    % active indices (row b of the matrix is drawn at height 16-b .. 17-b)
    \foreach \a in {2,5,6,10,13} {
        \foreach \b in {2,5,6,10,13} {
            \fill[cblue!60] ({\a-1},{16-\b}) rectangle ({\a},{17-\b});
        }
    }
    % identity on inactive
    \foreach \k in {1,3,4,7,8,9,11,12,14,15,16} {
        \fill[cgrey!45] ({\k-1},{16-\k}) rectangle ({\k},{17-\k});
    }
    \draw[thick] (0,0) rectangle (16,16);
    \node[below] at (8,0) {\small $\Heff^{\mathrm{emb}}$, $\nmax = 16$};
    \begin{scope}[xshift=7.5cm]
        \fill[cblue!60] (0,10) rectangle (0.6,10.6); \node[right] at (0.8,10.3) {\small entries of $\Heff$ ($\Nx = \{2,5,6,10,13\}$)};
        \fill[cgrey!45] (0,8.6) rectangle (0.6,9.2); \node[right] at (0.8,8.9) {\small identity on inactive indices};
        \node[align=left, right] at (0,5.5) {\small A column entering $\Nx$ fills its row and\\ \small column (row addition); a column leaving\\ \small $\Nx$ empties them (row deletion).};
    \end{scope}
\end{tikzpicture}
\caption{Fixed-length representation of the effective reduced Hessian: the active indices
$\Nx$ are scattered over the fixed index range $1,\dots,\nmax$, and every other index carries
an identity row.}
\label{fig:fixed-length}
\end{figure}

The identity rows do not interfere with the computation of the search direction:

\begin{lemma}[Decoupling of inactive indices] \label{lem:embedding}
Let $\Heff^{\mathrm{emb}} = LDL^\top$ be the $LDL^\top$ factorization without pivoting. Then
for every $k\notin\Nx$, row and column $k$ of $L$ are the unit vectors $e_k^\top$ and $e_k$,
and $D_{kk} = 1$. Consequently, for a right-hand side $r$ that vanishes outside $\Nx$, the
solution of $\Heff^{\mathrm{emb}}\eta = r$ vanishes outside $\Nx$, and its restriction to $\Nx$
solves $\Heff\eta_x = r_{\Nx}$.
\end{lemma}

\begin{proof}
Let $P$ be the permutation that moves the indices of $\Nx$ to the front. Then
$P^\top\Heff^{\mathrm{emb}}P = \diag(\Heff, I)$, because an inactive index has a $1$ on the
diagonal and zeros elsewhere in its row and column. With $\Heff = L_1D_1L_1^\top$, this
block-diagonal matrix has the block-diagonal factorization
\[
\diag(\Heff, I) = \diag(L_1,I)\,\diag(D_1,I)\,\diag(L_1,I)^\top,
\]
and the $LDL^\top$ factorization without pivoting is unique. Applying the
permutation backwards only moves rows and columns of these factors, so in the natural order
every inactive index $k$ still has row and column $e_k$ in $L$ and $D_{kk} = 1$. For the second
statement, $\diag(\Heff, I)\,[\eta_x;\eta_0] = [r_{\Nx}; 0]$ gives $\Heff\eta_x = r_{\Nx}$ and
$\eta_0 = 0$.
\end{proof}


\section{Expanding the Working Set}
\label{sec:expanding}

A step is blocked when a free variable reaches zero. The blocking index $q$ is then added to
the working set. There are three cases, depending on the set that $q$ belongs to. For a block
$k$ we write
\[
\mathcal{N}_x^{(k)} = \{\, j\in\Nx : \psi(j) = k \,\}
\]
for the nullspace columns of block $k$; these are the columns whose nullspace vectors
$z_j = e_j - e_{S_k}$ refer to the representative $S_k$.

\begin{lemma}[Ordering after an expansion] \label{lem:expand-ordering}
Let $q$ be the blocking index. The extended ordering is updated as follows.
\begin{enumerate}
    \item $q\in\Nx$: $\;\Nx\gets\Nx\setminus\{q\}$, $\;\Wx\gets\Wx\cup\{q\}$.
    \item $q = S_k\in\Sset$: choose a \textbf{replacement} $r\in\mathcal{N}_x^{(k)}$;
          then $S_k\gets r$, $\;\Nx\gets\Nx\setminus\{r\}$, $\;\Wx\gets\Wx\cup\{q\}$.
    \item $q$ is the free surplus $i\in\Ns$: $\;\Ns\gets\Ns\setminus\{i\}$, $\;\Ws\gets\Ws\cup\{i\}$.
\end{enumerate}
\end{lemma}

In case~(ii) the representative itself reaches zero, and another free column of its block
has to take its place in the list $\Sset$. The next lemma shows that one is always
available.

\begin{lemma}[Existence of a replacement] \label{lem:replacement}
In case~(ii), the set $\mathcal{N}_x^{(k)}$ is not empty.
\end{lemma}

\begin{proof}
By~\eqref{eq:direction-components}, the direction changes the representative by
$p_{S_k} = -\sum_{j\in\mathcal{N}_x^{(k)}}\eta_j$. The representative can only reach zero if it
decreases, $p_{S_k}<0$, so the sum is not empty.
\end{proof}

Any $r\in\mathcal{N}_x^{(k)}$ is a valid replacement. A good choice is the column with the largest
weight, $r = \arg\max_{j\in\mathcal{N}_x^{(k)}}x_j$, since it is the least likely to block in the
next step.

\begin{theorem}[Update of the nullspace basis] \label{thm:Z-update}
After an expansion, the nullspace basis changes as follows.
\begin{enumerate}
    \item $q\in\Nx$: the vector $z_q$ is removed.
    \item $q = S_k$ with replacement $r$: the vector $z_r$ is removed, and for every
          $j\in\mathcal{N}_x^{(k)}\setminus\{r\}$ the vector $z_j$ is replaced by
          \[
          z_j' = e_j - e_r = z_j - z_r .
          \]
          All other vectors are unchanged.
    \item $q$ the surplus $i$: the vector $e_{n+i}$ is removed.
\end{enumerate}
\end{theorem}

\begin{proof}
By \cref{cor:z-explicit} there is one basis vector per index of $\Nx$ and of $\Ns$, given by
$z_j = e_j - e_{\rep(j)}$ and $z_i = e_{n+i}$. In cases~(i) and~(iii) one index leaves $\Nx$ or
$\Ns$, and the representatives do not change, so only its vector disappears. In case~(ii), $r$
leaves $\Nx$ and the representative of block $k$ changes from $q$ to $r$. For
$j\in\mathcal{N}_x^{(k)}\setminus\{r\}$ the new vector is
$z_j' = e_j - e_r = (e_j - e_q) - (e_r - e_q) = z_j - z_r$. The vectors of other blocks refer
to other representatives and do not change.
\end{proof}

\begin{theorem}[Update of the effective reduced Hessian, expansion] \label{thm:H-expand}
After an expansion, $\Heff$ changes as follows.
\begin{enumerate}
    \item $q\in\Nx$: row and column $q$ are removed.
    \item $q = S_k$ with replacement $r$: row and column $r$ are removed, and for
          $i,j\in\mathcal{N}_x^{(k)}\setminus\{r\}$ and $\ell\in\Nx\setminus\mathcal{N}_x^{(k)}$
          \[
          (\Heff')_{\ell j} = (\Heff)_{\ell j} - (\Heff)_{\ell r},\qquad
          (\Heff')_{ij} = (\Heff)_{ij} - (\Heff)_{ir} - (\Heff)_{rj} + (\Heff)_{rr}.
          \]
    \item $q$ the surplus of inequality row $i$: row $i$ of $A_I$ joins $\Aeff$, and
          \begin{equation} \label{eq:rank-one-up}
          \Heff' = \Heff + v_iv_i^\top,\qquad
          v_i = Z_x^\top (A_I)_{i,:}^\top,\qquad (v_i)_j = (A_I)_{ij} - (A_I)_{i,\rep(j)}\;\;(j\in\Nx).
          \end{equation}
\end{enumerate}
\end{theorem}

\begin{proof}
We use $(\Heff)_{\ell j} = (\Aeff z_\ell)^\top(\Aeff z_j)$ from \cref{cor:Aeff}.

(i) The vector $z_q$ disappears, so the column $\Aeff z_q$ of $\Aeff Z_x$ is removed, and with
it row and column $q$ of the Gram matrix $\Heff$.

(ii) The vector $z_r$ disappears, which removes row and column $r$. For the remaining columns
of block $k$, $\Aeff z_j' = \Aeff z_j - \Aeff z_r$ by \cref{thm:Z-update}. For
$\ell\notin\mathcal{N}_x^{(k)}$ the vector $z_\ell$ is unchanged, so
$(\Heff')_{\ell j} = (\Aeff z_\ell)^\top(\Aeff z_j - \Aeff z_r) = (\Heff)_{\ell j} - (\Heff)_{\ell r}$.
For two columns of block $k$, expanding the product
$(\Aeff z_i - \Aeff z_r)^\top(\Aeff z_j - \Aeff z_r)$ gives the second formula.

(iii) The basis $Z_x$ does not change, but $\Aeff$ gains the row $(A_I)_{i,:}$. Hence
\[
\Aeff' Z_x = \begin{bmatrix}\Aeff Z_x\\ (A_I)_{i,:}Z_x\end{bmatrix},\qquad
\Heff' = (\Aeff Z_x)^\top(\Aeff Z_x) + \bigl((A_I)_{i,:}Z_x\bigr)^\top\bigl((A_I)_{i,:}Z_x\bigr)
= \Heff + v_iv_i^\top .
\]
Entry $j$ of $v_i = Z_x^\top(A_I)_{i,:}^\top$ is the product of column $j$ of $Z_x$, which is
$e_j - e_{\rep(j)}$, with row $i$ of $A_I$, i.e.\ $(v_i)_j = (A_I)_{ij} - (A_I)_{i,\rep(j)}$.
\end{proof}

The vector $v_i$ is row $i$ of $B_I = A_IZ_x$. Its entry $(v_i)_j$ is nonzero only if exactly
one of the columns $j$ and $\rep(j)$ has a nonzero coefficient in row $i$; in crew scheduling,
if exactly one of the two schedules covers flight $i$.

\begin{prop} \label{prop:expand-pd}
If $\Heff$ is positive definite before an expansion, it is positive definite after it.
\end{prop}

\begin{proof}
By \cref{cor:dimension}, $\Heff$ is positive definite if and only if the vectors $\Aeff z_j$,
$j\in\Nx$, are linearly independent.
(i) Removing one vector from a linearly independent set leaves it independent.
(ii) The new set consists of the unchanged vectors $\Aeff z_\ell$ and the vectors
$\Aeff z_j - \Aeff z_r$ for $j\in\mathcal{N}_x^{(k)}\setminus\{r\}$. A vanishing linear combination
$\sum_\ell\alpha_\ell\Aeff z_\ell + \sum_j\beta_j(\Aeff z_j - \Aeff z_r) = 0$ is a vanishing
combination of the old, independent vectors, in which the coefficient of $\Aeff z_j$ is
$\beta_j$ and that of $\Aeff z_\ell$ is $\alpha_\ell$; hence all $\alpha_\ell$ and $\beta_j$ are zero.
(iii) For every $\eta_x\ne0$, $\eta_x^\top(\Heff + v_iv_i^\top)\eta_x = \eta_x^\top\Heff\eta_x +
(v_i^\top\eta_x)^2 > 0$.
\end{proof}


\section{Reducing the Working Set}
\label{sec:reducing}

At a subspace minimizer with a negative multiplier, the index with the most negative
multiplier is released from the working set. By \cref{cor:multipliers-minimizer}, the
multipliers of columns (reduced costs) and of surpluses (duals) are both measured in units
of cost, so they can be compared directly. Releasing an index adds a basis vector to $Z$:

\begin{theorem}[Update of the effective reduced Hessian, reduction] \label{thm:H-reduce}
The two kinds of reduction change $\Heff$ as follows.
\begin{enumerate}
    \item[(iv)] Column $r\in\Wx$ is released ($\Wx\gets\Wx\setminus\{r\}$, $\Nx\gets\Nx\cup\{r\}$).
          The basis gains $z_r = e_r - e_{\rep(r)}$, and with $u = \Aeff z_r$,
          \begin{equation} \label{eq:bordered}
          \Heff' = \begin{bmatrix}\Heff & h\\ h^\top & h_{rr}\end{bmatrix},\qquad
          h = (\Aeff Z_x)^\top u,\qquad h_{rr} = u^\top u .
          \end{equation}
    \item[(v)] The surplus of inequality row $i$ is released ($\Ws\gets\Ws\setminus\{i\}$,
          $\Ns\gets\Ns\cup\{i\}$). The basis gains $e_{n+i}$, row $i$ leaves $\Aeff$, and
          \begin{equation} \label{eq:rank-one-down}
          \Heff' = \Heff - v_iv_i^\top
          \end{equation}
          with $v_i = Z_x^\top(A_I)_{i,:}^\top$ as in~\eqref{eq:rank-one-up}.
\end{enumerate}
\end{theorem}

\begin{proof}
(iv) The new column part of the basis is $Z_x' = [\,Z_x\;\;z_r\,]$, so
$\Aeff Z_x' = [\,\Aeff Z_x\;\;u\,]$ and
\[
\Heff' = \begin{bmatrix}(\Aeff Z_x)^\top\\ u^\top\end{bmatrix}\begin{bmatrix}\Aeff Z_x & u\end{bmatrix}
= \begin{bmatrix}\Heff & (\Aeff Z_x)^\top u\\ u^\top\Aeff Z_x & u^\top u\end{bmatrix}.
\]
(v) $Z_x$ does not change, and $\Aeff$ loses the row $(A_I)_{i,:}$. This is the reverse of
case~(iii) of \cref{thm:H-expand}, so $\Heff = \Heff' + v_iv_i^\top$.
\end{proof}

\begin{lemma}[Descent after a reduction] \label{lem:descent}
Suppose the reduced Hessian is positive definite after the reduction. Then the next search
direction moves the released variable away from its bound: $p_r>0$ in case~(iv), and
$(p_s)_i>0$ in case~(v).
\end{lemma}

\begin{proof}
Before the reduction the iterate is a subspace minimizer, so every entry of the reduced
gradient $Z^\top g$ is zero. The reduction adds one basis vector $z$ ($z_r$ or $e_{n+i}$) and
leaves the others unchanged, so the new reduced gradient has the single nonzero entry
$z^\top g$. For $z = z_r$ it is $g_r - g_{\rep(r)}$, and for $z = e_{n+i}$ it is $\pi_i$; by
\cref{cor:multipliers-minimizer} both equal the multiplier $\nu<0$ of the released index. Let
$e$ denote the unit vector of the new entry. The new reduced system
$\Hbar'\eta = -\delta Z'^\top g = -\delta\nu\,e$ gives $\eta = -\delta\nu\,\Hbar'^{-1}e$, whose own
component is
\[
\eta_e = -\delta\nu\,(\Hbar'^{-1})_{ee} > 0,
\]
since $\Hbar'^{-1}$ is positive definite and $\nu<0$. By~\eqref{eq:direction-components} this
component is the change of the released variable ($p_r = \eta_r$, $(p_s)_i = \eta_i$).
\end{proof}

Unlike an expansion, a reduction can make the reduced Hessian singular: the new column $u$ in
case~(iv) can lie in the span of the old ones, and the downdate in case~(v) can remove the
last direction of curvature. This is the subject of \cref{chap:singular}.


\section{Updating the $LDL^\top$ Factorization}
\label{sec:ldl-updates}

The updates of \cref{sec:expanding,sec:reducing} are applied to the factorization
$\Heff^{\mathrm{emb}} = LDL^\top$, with $L$ unit lower triangular and $D$ diagonal, in the
natural order of the indices. We use the sparse direct solver
CHOLMOD~\cite{chen2008cholmod}, which provides the three required modifications
\cite{davis1999modifying, davis2005row}:

\begin{center}
\renewcommand{\arraystretch}{1.25}
\begin{tabular}{lll}
\toprule
change of the working set & change of $\Heff^{\mathrm{emb}}$ & CHOLMOD routine \\
\midrule
(i) column leaves $\Nx$ & row/column $q$ becomes $e_q$ & \code{rowdel} \\
(ii) representative swap & rows/columns of $\mathcal{N}_x^{(k)}$ rewritten & \code{rowdel} + \code{rowadd} \\
(iii) surplus becomes active & $+\,v_iv_i^\top$ & \code{updown} (update) \\
(iv) column enters $\Nx$ & row/column $r$ set to $[h; h_{rr}]$ & \code{rowadd} \\
(v) surplus becomes free & $-\,v_iv_i^\top$ & \code{updown} (downdate) \\
\bottomrule
\end{tabular}
\end{center}

\subsection{Row additions and deletions}

The following proposition shows what a row addition computes. It explains the cost of the
operation and the meaning of the new pivot, which is used to detect singularity in
\cref{chap:singular}.

\begin{prop}[Row addition] \label{prop:rowadd}
Let row and column $k$ of $\Heff^{\mathrm{emb}}$ currently be $e_k$, and partition the indices
into $1 = \{1,\dots,k-1\}$, $k$ and $3 = \{k+1,\dots,\nmax\}$:
\[
L = \begin{bmatrix} L_{11} & & \\ 0 & 1 & \\ L_{31} & 0 & L_{33}\end{bmatrix},\qquad
D = \diag(D_1, 1, D_3).
\]
If row and column $k$ are set to $(h_1, h_{kk}, h_3)$, the new factor has
\begin{align*}
\ell_{k1} &= D_1^{-1}L_{11}^{-1}h_1 &&\text{(triangular solve with the leading block)},\\
d_k &= h_{kk} - \ell_{k1}^\top D_1\ell_{k1} &&\text{(new pivot)},\\
\ell_{3k} &= \bigl(h_3 - L_{31}D_1\ell_{k1}\bigr)/d_k &&\text{(new column below $k$)},\\
L_{33}'D_3'L_{33}'^\top &= L_{33}D_3L_{33}^\top - d_k\,\ell_{3k}\ell_{3k}^\top &&\text{(rank-one downdate)},
\end{align*}
while $L_{11}$, $D_1$ and $L_{31}$ are unchanged.
\end{prop}

\begin{proof}
The new factor has the form
$L' = \left[\begin{smallmatrix} L_{11} & & \\ \ell_{k1}^\top & 1 & \\ L_{31} & \ell_{3k} & L_{33}'\end{smallmatrix}\right]$,
$D' = \diag(D_1, d_k, D_3')$. Multiplying out $L'D'L'^\top$ block by block and comparing with
the new matrix:
\begin{itemize}
    \item block $(k,1)$: $\ell_{k1}^\top D_1L_{11}^\top = h_1^\top$, which gives $\ell_{k1}$;
    \item entry $(k,k)$: $\ell_{k1}^\top D_1\ell_{k1} + d_k = h_{kk}$, which gives $d_k$;
    \item block $(3,k)$: $L_{31}D_1\ell_{k1} + \ell_{3k}d_k = h_3$, which gives $\ell_{3k}$;
    \item block $(3,3)$: $L_{31}D_1L_{31}^\top + d_k\ell_{3k}\ell_{3k}^\top + L_{33}'D_3'L_{33}'^\top$ must
          equal the unchanged block $L_{31}D_1L_{31}^\top + L_{33}D_3L_{33}^\top$, which gives the
          downdate.
\end{itemize}
The blocks $(1,1)$ and $(3,1)$ do not involve row $k$ and stay as they are.
\end{proof}

A row deletion is the reverse operation: row and column $k$ are reset to $e_k$, which
requires the rank-one \emph{update} $L_{33}D_3L_{33}^\top + d_k\ell_{3k}\ell_{3k}^\top$ of the trailing
block. Two consequences of \cref{prop:rowadd} are important.

\begin{itemize}
    \item If $k$ is larger than every active index, the trailing block contains no active
          index, the downdate is trivial, and a row addition costs a single triangular solve
          with the active part, $O(|\Nx|^2)$ for a dense factor. This is used by the
          refactorization in \cref{sec:refactorization}.
    \item The determinant of the new active part is $\det\Heff' = \det\Heff\cdot d_k$, since
          the determinant of an $LDL^\top$ factorization is the product of the pivots and the
          other pivots of the active part do not change. A vanishing pivot $d_k$ therefore
          signals a singular $\Heff'$ (\cref{chap:singular}). If $k$ is the last active index,
          $d_k = h_{kk} - h^\top\Heff^{-1}h$ in the notation of~\eqref{eq:bordered}.
\end{itemize}

\subsection{Rank-one modifications}

The surplus cases (iii) and (v) modify the whole matrix by $\pm v_iv_i^\top$. The factorization
$LDL^\top\pm vv^\top = L'D'L'^\top$ is computed by the classical rank-one update and downdate
algorithms \cite{gill1974methods, davis1999modifying}, which process the columns of $L$ along
the path in the elimination tree that starts at the first nonzero of $v$. Their cost is
proportional to the number of entries of $L$ in these columns. An update keeps the matrix
positive definite (\cref{prop:expand-pd}). A downdate can make it singular; by the matrix
determinant lemma,
\[
\det\bigl(\Heff - v_iv_i^\top\bigr) = \det\Heff\cdot\bigl(1 - v_i^\top\Heff^{-1}v_i\bigr),
\]
so this happens exactly when $\omega_i := 1 - v_i^\top\Heff^{-1}v_i$ vanishes.
\Cref{chap:singular} shows that $\omega_i\ge0$ always holds for the downdates of the
active-set method, and treats the case $\omega_i = 0$.

\subsection{Sparse solves}

The right-hand side of~\eqref{eq:eff-system} vanishes outside $\Nx$. By \cref{lem:embedding},
the solve only has to touch the rows of $L$ that can be reached from $\Nx$ in the graph of
$L$, the set $\reach_L(\Nx)$. CHOLMOD supports such solves directly. Their cost is independent
of $\nmax$, which is what makes a large capacity $\nmax$ affordable.

\subsection{Assembly of a new row}

A row addition in case~(iv) needs the vector $h$ of~\eqref{eq:bordered}. By~\eqref{eq:implicit},
$h_k = u_k^\top M_\rho u$ with $u = a_r - a_{\rep(r)}$. In an implementation, $M_\rho u$ is
scattered into a dense $m$-vector once, with the entries of the rows of free surpluses set to
zero. Then $h_k = a_k^\top(M_\rho u) - a_{\rep(k)}^\top(M_\rho u)$ is computed for every
$k\in\Nx$, where the products with the $K$ representatives are computed once and reused. The
cost is $O\bigl((|\Nx| + K)\,\bar\ell\bigr)$, independent of $n$. A representative swap
(case~(ii)) rewrites the rows of $\mathcal{N}_x^{(k)}$ by a row deletion followed by a row addition
of the freshly assembled row; this costs the same as applying the formulas of
\cref{thm:H-expand}, and it needs no stored copy of $\Heff$.


\section{Structural Zeros and Periodic Refactorization}
\label{sec:refactorization}

A row deletion sets row and column $k$ of $L$ to $e_k$ \emph{numerically}, but a sparse
factorization keeps the positions of the removed entries in its sparsity pattern: the
pattern of a column of a simplicial factor can only grow. The same holds for rank-one
modifications. Over a long sequence of updates, the stored pattern of $L$ therefore becomes
the union of the patterns of all factors that ever existed. To quantify this effect we
distinguish the following quantities, which are measured in \cref{sec:results-fill}.

\begin{definition}[Stored and effective entries] \label{def:stored-effective}
The \textbf{stored entries} of $L$ are the off-diagonal positions in its sparsity pattern.
The \textbf{effective entries} are the stored entries whose value is nonzero. The difference
are the \textbf{structural zeros}. The \textbf{dense equivalent} is $|\Nx|(|\Nx|-1)/2$, the
number of off-diagonal entries of a dense factor of $\Heff$.
\end{definition}

Every solve and every update iterates over the stored entries of the columns it touches,
including the structural zeros. Moreover, structural zeros can connect active indices to
indices that are no longer active, which enlarges $\reach_L(\Nx)$. The stored pattern only
returns to the effective one when the factorization is rebuilt.

\begin{definition}[Refactorization] \label{def:refactorization}
A \textbf{refactorization} resets the factor to the identity and re-activates all indices of
$\Nx$ by row additions \emph{in increasing index order}, with the rows assembled as in
\cref{sec:ldl-updates}.
\end{definition}

\begin{prop}[Cost of a refactorization] \label{prop:refactor-cost}
A refactorization costs $O(|\Nx|^3)$ arithmetic operations for the factor (about
$|\Nx|^3/3$ for a dense $\Heff$) and $O\bigl(|\Nx|(|\Nx|+K)\,\bar\ell\bigr)$ for the assembly of
the rows. Its stored pattern contains no structural zeros.
\end{prop}

\begin{proof}
When the indices are added in increasing order, every new index is the largest active index,
so by \cref{prop:rowadd} each row addition is a single triangular solve with the current
active part, which costs $O(t^2)$ for the $t$-th index. Summing over $t = 1,\dots,|\Nx|$ gives
$O(|\Nx|^3)$, about $|\Nx|^3/3$ as for a dense Cholesky factorization. Each row costs
$O\bigl((|\Nx|+K)\,\bar\ell\bigr)$ to assemble (\cref{sec:ldl-updates}). Since no row is ever
deleted during a refactorization, every stored entry is computed from the current matrix.
\end{proof}

The method refactorizes after every $\tau$ updates of the factorization (row additions, row
deletions and rank-one modifications), and whenever a pivot has become too small
(\cref{chap:singular}). The \textbf{refactorization interval} $\tau$ balances the cost of a
rebuild against the growth of the structural zeros, which is roughly proportional to the number
of updates. Its effect is studied in \cref{sec:results-fill}.

\begin{remark}[Natural ordering] \label{rem:natural-ordering}
Sparse direct solvers usually reorder a matrix to reduce fill-in. For $\Heff$ this would gain
nothing: the effective factor is almost exactly as dense as a dense factor of $\Heff$
(\cref{sec:results-fill}), because columns of the same block share their representative and
are therefore strongly coupled. The natural ordering has two advantages: index $j$ of the
factor is column $j$ of the master problem, and the refactorization of
\cref{def:refactorization} only appends rows.
\end{remark}


\section{Hot Start Between Column Generation Iterations}
\label{sec:hot-start}

Between two master problems, regularized column generation adds columns and moves the
stability center (\cref{alg:rcg}). Both changes can be absorbed without touching the
factorization.

\begin{theorem}[Hot start] \label{thm:hot-start}
Let $y = [x;s]$ be the final iterate of a master problem, with its extended ordering and the
factorization of $\Heff$.
\begin{enumerate}
    \item \textbf{New column.} If a column $n+1\le\nmax$ is added with weight $x_{n+1} = 0$
          and placed in the working set $\Wx$, then $y$ remains feasible, $\pi(y)$ is
          unchanged, and the nullspace basis, $\Aeff$, $\Heff$ and its factorization are
          unchanged. The multiplier of the new column is its reduced cost, which is negative
          because pricing selected it.
    \item \textbf{New center.} If the center changes from $\pihat$ to $\pihat + \Delta$,
          then $y$ remains feasible, the ordering and the factorization are unchanged, and the
          dual estimate becomes $\pi(y) + \Delta$.
\end{enumerate}
\end{theorem}

\begin{proof}
(i) The new column has weight zero, so the convexity constraint of its block and the product
$\Abar y$ are unchanged; hence $y$ stays feasible and $\pi(y) = \pihat + \frac{1}{\delta}(b - \Abar y)$
is unchanged. The new index is placed in $\Wx$, so the sets $\Sset$, $\Nx$, $\Ns$, and with them
$Z$, $\Aeff$ and $\Heff$, are unchanged. In the fixed-length representation, index $n+1$
already carries an identity row, so the factorization is unchanged too. By
\cref{cor:multipliers-minimizer}, the multiplier of a working-set column is its reduced cost
$c_{n+1} - a_{n+1}^\top\pi - \mu_{\psi(n+1)}$, which pricing found to be negative.
(ii) The constraints $\Cbar y = \one$, $y\ge0$ do not involve $\pihat$, and
$\Heff = (\Aeff Z_x)^\top(\Aeff Z_x)$ does not depend on it either. Finally
$\pi(y) = (\pihat + \Delta) + \frac{1}{\delta}(b - \Abar y)$ is the old dual estimate plus $\Delta$.
\end{proof}

After a new column, the previous optimum is a subspace minimizer whose only negative
multipliers are those of the new columns. The next master problem therefore typically starts
by releasing the new column (case~(iv)). After a new center, the gradient changes and the
iterate is in general no longer a subspace minimizer; the next master problem starts with a
direction computation on the inherited factorization. The surpluses are never added or
removed, and their status is carried over like that of the columns.

\subsection{Center-only iterations}

The hot start also explains the behaviour of the center-only iterations of
\cref{sec:center-update}, in which the center moves to the current dual and no column is
added. Let $\Pi = \Abar Z\,\Hbar^{-1}(\Abar Z)^\top$ denote the orthogonal projector onto the range
of $\Abar Z$.

\begin{prop}[Center-only iterations] \label{prop:center-only}
Let $y$ be optimal for the current master problem, with residual $r = b - \Abar y$ and dual
$\pi = \pihat + r/\delta$, and let the center move to $\pihat' = \pi$. Suppose that the next
master problem takes one full step and keeps the working set. Then
\begin{enumerate}
    \item the new residual is $r' = (I-\Pi)\,r$ and the new dual is $\pi' = \pi + (I-\Pi)\,r/\delta$;
    \item if the working set stays the same in the following center-only iterations, the
          primal iterate no longer changes, and every further iteration moves the dual by the
          fixed vector $r'/\delta$:
          \[
          \pi^{(t+1)} = \pi^{(t)} + \frac{r'}{\delta},\qquad
          (\nu_x)_j^{(t+1)} = (\nu_x)_j^{(t)} - \frac{u_j^\top r'}{\delta},\qquad
          (\nu_s)_i^{(t+1)} = (\nu_s)_i^{(t)} + \frac{r'_i}{\delta}.
          \]
\end{enumerate}
\end{prop}

\begin{proof}
(i) The center moves by $\Delta = \pihat' - \pihat = \pi - \pihat = r/\delta$. By
\cref{thm:hot-start}, the dual at the old iterate becomes $\pi + \Delta$, and since the gradient
is $g = \cbar - \Abar^\top\pi$ (\cref{lem:gradient-dual}), it changes by $-\Abar^\top\Delta$. The
old iterate was optimal, so its reduced gradient was zero, and the new reduced gradient is
\[
Z^\top g' = Z^\top g - Z^\top\Abar^\top\Delta = -(\Abar Z)^\top\Delta .
\]
The reduced system $\Hbar\eta = -\delta Z^\top g'$ gives $\eta = \delta\,\Hbar^{-1}(\Abar Z)^\top\Delta$,
hence $\Abar p = \Abar Z\eta = \delta\,\Pi\Delta$. After the full step, the dual is
\[
\pi' = (\pi + \Delta) - \tfrac{1}{\delta}\Abar p = \pi + \Delta - \Pi\Delta = \pi + (I - \Pi)\Delta,
\]
and by \cref{prop:residual} the new residual is $r' = \delta(\pi' - \pihat') = \delta(I-\Pi)\Delta = (I-\Pi)\,r$.

(ii) The vector $r' = (I-\Pi)r$ is orthogonal to the range of $\Abar Z$, i.e.\
$(\Abar Z)^\top r' = 0$. In the next center-only iteration the center moves by
$\Delta' = r'/\delta$, and the same computation as in~(i) gives the reduced gradient
$-(\Abar Z)^\top\Delta' = 0$: the direction is zero and $y$ does not move. The dual becomes
$\pi' + \Delta'$; since the new center is $\pi'$, the residual is
$\delta\bigl((\pi'+\Delta') - \pi'\bigr) = r'$ again, and the argument repeats. The multipliers follow from \cref{cor:multipliers-minimizer}. A shift of $\pi$ by
$\Delta'$ changes $g_j$ by $-a_j^\top\Delta'$, hence $(\nu_x)_j = g_j - g_{\rep(j)}$ by
$-u_j^\top\Delta'$; and $(\nu_s)_i = \pi_i$ changes by $\Delta'_i$.
\end{proof}

\Cref{prop:center-only} shows that in a center-only phase the dual walks along a fixed
direction with step length $\lVert r'\rVert/\delta$, while the primal iterate stays put. The
walk ends when a multiplier changes sign, so that the working set changes, or when pricing
finds an improving column. If the residual $r'$ is small, many such iterations are needed,
which explains the long tail of center-only iterations observed in \cref{sec:results-phases}.
Each of them is cheap: one direction computation, which returns zero, and one multiplier test.


\section{The Complete Method}
\label{sec:complete}

\Cref{alg:hot-qpas} combines the results of
\cref{chap:qpas,chap:nullspace,chap:updates}. It is called in the first line of the loop of
regularized column generation (\cref{alg:rcg}), and its state, consisting of the iterate, the
extended ordering and the factorization, is kept from one call to the next.

\begin{algobox}{Hot-Started Nullspace Active-Set Method}{alg:hot-qpas}
\State \textbf{Input:} state of the previous call (iterate $y$, ordering, factor of $\Heff$),
       new columns, new center $\pihat$
\State Put the new columns into $\Wx$ with weight $0$ \Comment{\cref{thm:hot-start}}
\State $\pi\gets\pihat + (b - \Abar y)/\delta$
\Repeat
    \If{the last step was not a full step}
        \State Compute $p$ by \cref{alg:nullspace-iteration} (lines 2--4)
    \Else
        \State $p\gets0$ \Comment{\cref{prop:full-step}}
    \EndIf
    \If{$p\approx0$}
        \State Compute the multipliers (\cref{cor:multipliers-minimizer})
        \If{all multipliers $\ge-\varepsilon$} \textbf{return} $y$, $\pi$, $\mu$ \EndIf
        \State Release the most negative index: \code{rowadd} (column) or downdate (surplus)
        \State If the new pivot is too small, resolve the singularity (\cref{alg:singular})
    \Else
        \State Ratio test over $\Sset\cup\Nx\cup\Ns$; $y\gets y + \alpha p$;
               $\pi\gets\pi - \frac{\alpha}{\delta}\Abar p$
        \If{blocked by $q$}
            \State Expand: \code{rowdel} ($q\in\Nx$), representative swap ($q\in\Sset$), or
                   update ($q\in\Ns$)
        \EndIf
    \EndIf
    \State Refactorize if $\tau$ updates have passed since the last refactorization
\Until{returned}
\end{algobox}

\chapter{Singular Reduced Hessians}
\label{chap:singular}

\vspace*{\fill}

The nullspace method requires the reduced Hessian to be positive definite
(\cref{ass:nullspace}~(ii)). The Hessian of the regularized master problem is only positive
semidefinite, and its rank is bounded by the number of linking rows, while the number of
columns keeps growing. It is therefore unavoidable that a reduction of the working set
occasionally produces a nullspace direction along which the objective has no curvature.

\medskip

This chapter shows where such singularities can arise, how they are detected from the
pivots of the factorization, and how they are resolved. The resolution is a
\textbf{zero-curvature step}: along a direction in the kernel of the reduced Hessian, the
objective decreases linearly, and the step is taken until a variable reaches its bound. We
prove that adding this variable to the working set removes the singularity, so that the
interrupted reduction can be completed. The treatment covers released columns and released
surpluses with one argument, and keeps every working-set variable exactly at its bound.

\vspace*{\fill}

\newpage


\section{Where Singularity Can Arise}
\label{sec:where-singular}

By \cref{cor:dimension}, the reduced Hessian is positive definite if and only if
$\Heff = (\Aeff Z_x)^\top(\Aeff Z_x)$ is, i.e.\ if and only if the vectors
$\Aeff z_j$, $j\in\Nx$, are linearly independent. \Cref{prop:expand-pd} shows that expansions
preserve this property. Singularity can therefore only arise in one of the two reductions
of \cref{thm:H-reduce}.

\begin{prop}[Singular reductions] \label{prop:singular-cases}
Let $\Heff$ be positive definite before the reduction.
\begin{enumerate}
    \item[(iv)] After releasing column $r$, the matrix $\Heff'$ of~\eqref{eq:bordered} is
          singular if and only if $\Aeff z_r\in\spn\{\Aeff z_k : k\in\Nx\}$, which is the case if
          and only if the Schur complement $\sigma_r = h_{rr} - h^\top\Heff^{-1}h$ vanishes.
    \item[(v)] After releasing the surplus of row $i$, the matrix $\Heff' = \Heff - v_iv_i^\top$
          is singular if and only if $\omega_i = 1 - v_i^\top\Heff^{-1}v_i$ vanishes.
\end{enumerate}
Moreover $\sigma_r\ge0$ and $0\le\omega_i\le1$.
\end{prop}

\begin{proof}
(iv) $\Heff'$ is the Gram matrix of the vectors $\Aeff z_k$, $k\in\Nx$, and $u = \Aeff z_r$. It
is singular if and only if these vectors are linearly dependent; since the vectors $\Aeff z_k$
alone are independent, this means $u\in\spn\{\Aeff z_k\}$. For the bordered matrix
$\Heff' = \left[\begin{smallmatrix}\Heff & h\\ h^\top & h_{rr}\end{smallmatrix}\right]$, block
elimination gives $\det\Heff' = \det\Heff\cdot(h_{rr} - h^\top\Heff^{-1}h) = \det\Heff\cdot\sigma_r$,
so $\Heff'$ is singular if and only if $\sigma_r = 0$.

(v) By the matrix determinant lemma, $\det(\Heff - v_iv_i^\top) = \det\Heff\cdot(1 - v_i^\top\Heff^{-1}v_i)
= \det\Heff\cdot\omega_i$, so $\Heff'$ is singular if and only if $\omega_i = 0$.

In both cases $\Heff'$ is a Gram matrix, of $[\,\Aeff Z_x\;\;u\,]$ in case~(iv) and of $\Aeff'Z_x$
($\Aeff$ without row $i$) in case~(v). It is therefore positive semidefinite, so
$\det\Heff'\ge0$, and with $\det\Heff>0$ we get $\sigma_r\ge0$ and $\omega_i\ge0$. Finally
$\omega_i\le1$ because $\Heff^{-1}$ is positive definite, so $v_i^\top\Heff^{-1}v_i\ge0$.
\end{proof}

\Cref{prop:singular-cases} also shows that a downdate in case~(v) can never make $\Heff$
indefinite: in exact arithmetic, the worst case is a singular matrix.


\section{Detection}
\label{sec:detection}

Testing $\sigma_r$ or $\omega_i$ before the update would cost an additional solve for every
reduction. Instead, the update is performed, and the pivots $D_{kk}$ of the factorization
$\Heff^{\mathrm{emb}} = LDL^\top$ are inspected afterwards.

\begin{definition}[Pivot invariant] \label{def:pivot-invariant}
For $k\in\Nx$, let $\gamma_k = \lVert a_k - a_{\rep(k)}\rVert^2$ be the \textbf{pivot scale} of
index $k$ (with all rows, not only those of $\Aeff$). The factorization satisfies the
\textbf{pivot invariant} with tolerance $\varepsilon_{\mathrm{piv}}>0$ if
\[
D_{kk}\;>\;\varepsilon_{\mathrm{piv}}\,\gamma_k\qquad\text{for all }k\in\Nx .
\]
\end{definition}

The scale $\gamma_k$ uses all rows on purpose: a column whose coefficients lie entirely in rows
with free surpluses has $(\Heff)_{kk} = 0$, and a scale based on $\Heff$ would hide this
singularity.

The determinant of $\Heff'$ is the product of its pivots. In case~(iv) the update adds one new
pivot, and in case~(v) the pivots change by the factor $\omega_i$ in total, since
$\det\Heff' = \det\Heff\cdot\omega_i$. If all pivots satisfied the invariant before the update, a
vanishing $\sigma_r$ or $\omega_i$ must therefore show up as a pivot that violates it after the
update. This detection is a numerical heuristic rather than a theorem about floating-point
arithmetic, but it follows directly from the determinant identities of
\cref{prop:singular-cases}. When it fires, the update is undone by a refactorization with the
previous working set (\cref{def:refactorization}), which also removes any large entries
created by dividing by a tiny pivot. We use $\varepsilon_{\mathrm{piv}} = 10^{-10}$.


\section{Zero-Curvature Directions}
\label{sec:zero-curvature}

Assume now that the iterate $y$ is a subspace minimizer for the working set $\Wset$, that the
index to be released has a negative multiplier, and that its release would make the reduced
Hessian singular. We construct a direction along which the objective decreases linearly and
the working set (apart from the released index) is respected. Recall from
\cref{cor:multipliers-minimizer} what is known at a subspace minimizer:
\begin{equation} \label{eq:minimizer-facts}
\begin{aligned}
g^\top z_k &= g_k - g_{\rep(k)} = 0 \;\;(k\in\Nx), &\qquad \pi_\ell &= 0 \;\;(\ell\in\Ns),\\
g_r - g_{\rep(r)} &= (\nu_x)_r \;\;(r\in\Wx), & \pi_i &= (\nu_s)_i \;\;(i\in\Ws).
\end{aligned}
\end{equation}
In this section $z_k = e_k - e_{\rep(k)}$ denotes the column part in $\R^n$.

\subsection{Released column}

\begin{theorem}[Zero-curvature direction, column] \label{thm:zc-column}
Let $r\in\Wx$ have the multiplier $(\nu_x)_r<0$, and suppose that the release of $r$ is
singular, i.e.\ $\Aeff z_r = \sum_{k\in\Nx}t_k\,\Aeff z_k$ for some coefficients $t$ (these are
$t = \Heff^{-1}h$ with $h$ from~\eqref{eq:bordered}). Define $d = [d_x; d_s]$ by
\begin{equation} \label{eq:zc-column}
d_x = z_r - \sum_{k\in\Nx}t_k\,z_k,\qquad
(d_s)_\ell = (A_Id_x)_\ell\;\;(\ell\in\Ns),\qquad (d_s)_\ell = 0\;\;(\ell\in\Ws).
\end{equation}
Then
\begin{enumerate}
    \item $Cd_x = 0$, and $d_\ell = 0$ for every $\ell\in\Wset\setminus\{r\}$;
    \item $\Abar d = 0$;
    \item $\nabla f(y)^\top d = (\nu_x)_r < 0$.
\end{enumerate}
\end{theorem}

\begin{proof}
(i) Every $z_k = e_k - e_{\rep(k)}$ has one entry $+1$ and one entry $-1$ in the same block, so
$Cz_k = 0$ and therefore $Cd_x = 0$. The vector $d_x$ is supported on $\{r\}\cup\Nx\cup\Sset$, so it
vanishes on $\Wx\setminus\{r\}$, and $d_s$ vanishes on $\Ws$ by definition.

(ii) The rows of $\Abar d$ are $A_Ed_x$ (equality rows) and $A_Id_x - d_s$ (inequality rows). On
the rows of $\Aeff$, i.e.\ the equality rows and the inequality rows in $\Ws$ (where $d_s = 0$),
$\Abar d$ equals $\Aeff d_x = \Aeff z_r - \sum_kt_k\Aeff z_k$, which is zero by the assumption on
$t$. On an inequality row $\ell\in\Ns$, $(\Abar d)_\ell = (A_Id_x)_\ell - (d_s)_\ell = 0$ by the
definition of $d_s$.

(iii) Splitting $g^\top d = g_x^\top d_x + \sum_\ell\pi_\ell(d_s)_\ell$ (recall $g_s = \pi_I$) and
using~\eqref{eq:minimizer-facts},
\[
g^\top d = \bigl(g_r - g_{\rep(r)}\bigr) - \sum_{k\in\Nx}t_k\underbrace{\bigl(g_k - g_{\rep(k)}\bigr)}_{=0}
+ \sum_{\ell\in\Ns}\underbrace{\pi_\ell}_{=0}(d_s)_\ell = (\nu_x)_r . \qedhere
\]
\end{proof}

\subsection{Released surplus}

\begin{theorem}[Zero-curvature direction, surplus] \label{thm:zc-surplus}
Let the surplus of row $i\in\Ws$ have the multiplier $(\nu_s)_i = \pi_i < 0$, and suppose that
its release is singular, $\omega_i = 1 - v_i^\top\Heff^{-1}v_i = 0$. Let $t = \Heff^{-1}v_i$, so
that $v_i^\top t = 1$. Define $d = [d_x;d_s]$ by
\begin{equation} \label{eq:zc-surplus}
d_x = \sum_{k\in\Nx}t_k\,z_k,\qquad
(d_s)_\ell = (A_Id_x)_\ell\;\;(\ell\in\Ns\cup\{i\}),\qquad (d_s)_\ell = 0\;\;\text{otherwise}.
\end{equation}
Then
\begin{enumerate}
    \item $Cd_x = 0$, $d_x$ vanishes on $\Wx$, $d_s$ vanishes on $\Ws\setminus\{i\}$, and $(d_s)_i = 1$;
    \item $\Abar d = 0$;
    \item $\nabla f(y)^\top d = \pi_i < 0$.
\end{enumerate}
\end{theorem}

\begin{proof}
(i) As in \cref{thm:zc-column}, $Cd_x = 0$ and $d_x$ vanishes on $\Wx$. Further, $d_x = Z_xt$,
so $(d_s)_i = (A_IZ_xt)_i = (A_I)_{i,:}Z_x\,t = v_i^\top t = 1$, using
$v_i = Z_x^\top(A_I)_{i,:}^\top$ from~\eqref{eq:rank-one-up}.

(ii) Let $\Aeff'$ be $\Aeff$ without row $i$; its rows are the equality rows and the inequality
rows in $\Ws\setminus\{i\}$. By \cref{thm:H-reduce}, $(\Aeff'Z_x)^\top(\Aeff'Z_x) = \Heff - v_iv_i^\top$,
hence
\[
\lVert\Aeff'd_x\rVert^2 = t^\top(\Heff - v_iv_i^\top)\,t = t^\top v_i - (v_i^\top t)^2 = 1 - 1 = 0,
\]
using $\Heff t = v_i$ and $v_i^\top t = 1$. So $\Aeff'd_x = 0$, which means that $\Abar d$ vanishes
on the rows of $\Aeff'$ (where $d_s = 0$). On the remaining inequality rows
$\ell\in\Ns\cup\{i\}$, $(\Abar d)_\ell = (A_Id_x)_\ell - (d_s)_\ell = 0$ by construction.

(iii) Using~\eqref{eq:minimizer-facts}, $g_x^\top d_x = \sum_kt_k(g_k - g_{\rep(k)}) = 0$ and
$\pi_\ell = 0$ for $\ell\in\Ns$, so only the component of surplus $i$ remains:
$g^\top d = \pi_i\,(d_s)_i = \pi_i$.
\end{proof}

\begin{remark} \label{rem:zc-kappa}
For a ``$\le$'' row with slack coefficient $\kappa_i = +1$ (\cref{rem:kappa-summary}), the
column part is scaled by $-\kappa_i$, $d_x = -\kappa_i\sum_kt_kz_k$, and the surplus part is
$(d_s)_\ell = -\kappa_\ell(A_Id_x)_\ell$. Then again $(d_s)_i = 1$ and
$g^\top d = -\kappa_i\pi_i = (\nu_s)_i < 0$.
\end{remark}

\begin{cor}[Linear decrease] \label{cor:linear-decrease}
For the directions of \cref{thm:zc-column,thm:zc-surplus},
\[
f(y + \alpha d) = f(y) + \alpha\,\nabla f(y)^\top d,\qquad \pi(y+\alpha d) = \pi(y)
\qquad\text{for all }\alpha .
\]
\end{cor}

\begin{proof}
Expanding the quadratic $f$, $f(y+\alpha d) = f(y) + \alpha\nabla f(y)^\top d +
\frac{\alpha^2}{2\delta}\lVert\Abar d\rVert^2$, and $\Abar d = 0$. The dual estimate
$\pi(y) = \pihat + \frac{1}{\delta}(b - \Abar y)$ depends on $y$ only through $\Abar y$, which does
not change.
\end{proof}

\begin{remark}[Fully masked columns] \label{rem:masked-column}
If all nonzero coefficients of $u = a_r - a_{\rep(r)}$ lie in rows with free surpluses, then
$\Aeff z_r = 0$, so $h = 0$ and $h_{rr} = 0$. \Cref{thm:zc-column} then applies with $t = 0$:
the direction is $z_r$ together with the surplus components that absorb $A_Iz_r$.
\end{remark}


\section{The Zero-Curvature Step Removes the Singularity}
\label{sec:zc-step}

Along $d$ the objective decreases linearly (\cref{cor:linear-decrease}), so the method moves
as far as the bounds allow. In this section $\rho$ denotes the released index: the column $r$ of
\cref{thm:zc-column} or the surplus $i$ of \cref{thm:zc-surplus}.

\begin{lemma}[Finite step] \label{lem:finite-step}
The directions of \cref{thm:zc-column,thm:zc-surplus} have at least one negative component.
Hence the ratio test
\[
\alpha^\star = \min_{\ell:\;d_\ell<0}\;-\frac{y_\ell}{d_\ell}
\]
is finite, and it is attained at a \textbf{blocking index} $q$ in
$\Sset\cup\Nx\cup\Ns$.
\end{lemma}

\begin{proof}
Suppose $d\ge0$. By part~(i) of the theorems, every block sum of $d_x$ is zero ($Cd_x = 0$).
A nonnegative vector whose block sums are zero is zero, so $d_x = 0$. Then
$\Abar d = [\,A_Ed_x;\; A_Id_x - d_s\,] = [\,0;\;-d_s\,]$, and $\Abar d = 0$ (part~(ii)) forces
$d_s = 0$. But $d = 0$ contradicts $\nabla f(y)^\top d < 0$ (part~(iii)). So some component of
$d$ is negative, and the minimum is taken over a nonempty set. The blocking index cannot lie in
$\Wset\setminus\{\rho\}$, where $d$ vanishes, and it is not the released index $\rho$, whose
component of $d$ is positive ($d_r = 1$ or $(d_s)_i = 1$). Hence $q\in\Sset\cup\Nx\cup\Ns$.
\end{proof}

After the step, the blocking index $q$ is added to the working set by an expansion
(\cref{sec:expanding}), and the release is attempted again. The following theorem shows that
this attempt cannot fail for the same reason.

\begin{theorem}[Removal of the dependency] \label{thm:removal}
Let $\rho$ be the released index and $q$ the blocking index of
\cref{lem:finite-step}. Then the reduced Hessian of the working set
$\Wset' = (\Wset\setminus\{\rho\})\cup\{q\}$ is positive definite. In other words, after the
zero-curvature step and the expansion by $q$, the retried release of $\rho$ succeeds.
\end{theorem}

\begin{proof}
For any working set with nullspace basis $Z$, the reduced Hessian $Z^\top\Abar^\top\Abar Z$ is
positive definite if and only if $\Abar Z\eta\ne0$ for all $\eta\ne0$, i.e.\ if and only if
$\Abar$ is injective on the nullspace $\Null(D)$ of the working set. We use this for three
spaces:
\begin{align*}
N_0 &= \Null(D_\Wset) && \text{(before the release)},\\
V &= \Null(D_{\Wset\setminus\{\rho\}}) && \text{(after the release)},\\
V' &= V\cap\{y : y_q = 0\} = \Null(D_{\Wset'}) && \text{(after the release and the expansion by $q$)}.
\end{align*}
\emph{Step 1: $\ker\Abar\cap V = \spn\{d\}$.} Before the release the reduced Hessian was positive
definite, so $\ker\Abar\cap N_0 = \{0\}$. Removing one constraint from the working set adds one
dimension, $N_0\subset V$ with $\dim V = \dim N_0 + 1$. If $\ker\Abar\cap V$ contained two
linearly independent vectors, some nonzero combination of them would satisfy the one extra
constraint ($y_\rho = 0$) and therefore lie in $N_0\cap\ker\Abar = \{0\}$, a contradiction. So
$\dim(\ker\Abar\cap V)\le1$. The direction $d$ lies in $V$ (part~(i) of
\cref{thm:zc-column,thm:zc-surplus}) and in $\ker\Abar$ (part~(ii)), and $d\ne0$, so
$\ker\Abar\cap V = \spn\{d\}$.

\emph{Step 2: $\ker\Abar\cap V' = \{0\}$.} Every vector of $\ker\Abar\cap V'$ is a multiple
$\beta d$ of $d$ by Step~1, and it has a zero component $q$. Since $d_q<0$, this forces
$\beta = 0$. Hence $\Abar$ is injective on $V' = \Null(D_{\Wset'})$, and the reduced Hessian of
$\Wset'$ is positive definite.
\end{proof}

The retried release starts from the same kind of state as the first attempt: a subspace
minimizer with a negative multiplier for $\rho$.

\begin{prop}[The iterate stays a subspace minimizer] \label{prop:stays-minimizer}
After the zero-curvature step and the expansion by $q$, the iterate is a subspace minimizer
of the new working set $(\Wset\setminus\{\rho\})\cup\{q\}$ with $\rho$ still fixed, and the
multiplier of $\rho$ is unchanged.
\end{prop}

\begin{proof}
By \cref{cor:linear-decrease}, $\pi$ and hence the gradient $g = \cbar - \Abar^\top\pi$ do not change
along the step. The expansion by $q$ changes the nullspace basis as in \cref{thm:Z-update}: it
removes a vector, or, if $q$ is a representative, replaces vectors $z_j$ by
$z_j - z_{r'}$ for the replacement $r'$. In all cases the new basis vectors are combinations of
old ones, whose reduced gradient entries $g^\top z$ were zero; so the new reduced gradient is zero
as well. The multiplier of $\rho$ is computed from the unchanged gradient by
\cref{cor:multipliers-minimizer}.
\end{proof}

\begin{remark}[The representative of the released column blocks] \label{rem:rep-blocks}
In \cref{thm:zc-column}, the blocking index may be the representative $\rep(r)$ of the block
of $r$. Then $r$, which now has the positive weight $\alpha^\star$, can simply take over as
the new representative (case~(ii) of \cref{lem:expand-ordering} with replacement $r$). This
completes the release at once.
\end{remark}

\Cref{alg:singular} summarizes the procedure. Each pass either completes the release or adds
one more index to the working set. The number of passes is therefore bounded by the number of
free variables.

\begin{algobox}{Release with Singularity Resolution}{alg:singular}
\State \textbf{Input:} subspace minimizer $y$, index $\rho$ (column $r$ or surplus $i$) with
       negative multiplier
\Loop
    \State Apply the reduction update (\code{rowadd} for a column, downdate for a surplus)
    \If{the pivot invariant holds}
        \State \textbf{return} \Comment{release completed}
    \EndIf
    \State Undo by refactorization with $\rho$ still in the working set
    \State Solve $\Heff t = h$ (column) or $\Heff t = v_i$ (surplus)
    \State Form $d$ by~\eqref{eq:zc-column} or~\eqref{eq:zc-surplus}
    \State Ratio test along $d$; $y\gets y+\alpha^\star d$ \Comment{$\pi$ is unchanged}
    \If{$\rho = r$ and $q = \rep(r)$}
        \State Make $r$ the representative, put $q$ into $\Wx$; \textbf{return}
               \Comment{\cref{rem:rep-blocks}}
    \EndIf
    \State Expand the working set by $q$ \Comment{\cref{thm:removal}: the next pass succeeds}
\EndLoop
\end{algobox}

\begin{remark}[Why not remove the largest coefficient?] \label{rem:argmax}
A simpler rule moves the index $k = \arg\max_{k\in\Nx}|t_k|$ of the largest coefficient of the
linear dependence into the working set and retries. This also removes the dependency, but
it ignores the value of $x_k$. If $x_k>0$, the working-set constraint $x_k = 0$ is violated by
the iterate, and its multiplier no longer certifies optimality: the method can stop with
$x_k>0$ and a positive multiplier. The zero-curvature step avoids this. The iterate stays
feasible, every working-set variable is exactly at its bound, the objective decreases
strictly, and the index added to the working set is the one selected by a ratio test rather
than by a heuristic. The same argument covers released columns and released surpluses.
\end{remark}

\chapter{A Dense Factorization of the Effective Reduced Hessian}
\chaptermark{Dense Factorization}
\label{chap:dense}

\vspace*{\fill}

\Cref{chap:updates} expressed every change of the working set as a row addition, a row
deletion or a rank-one modification of an $LDL^\top$ factorization, and realized these
operations with a sparse factorization of fixed dimension $\nmax$. This is the natural choice
for a matrix that is embedded in an index range of the size of the column capacity. The matrix
that actually has to be factorized, however, is small and dense: $\Heff$ never has more rows
than there are linking rows, and almost all of its entries are nonzero.

\medskip

This chapter develops a \textbf{dense factorization} that exploits these two properties. It
stores only the active indices, in the order in which they entered $\Nx$, so that a new index
is always \textbf{appended} as the last row of the factor. A removed index is deleted from the
factor by a \textbf{rank-one update} of the rows behind it, and the factor is compacted at the
same time. We derive these operations from the classical rank-one modification algorithm and
show four properties that make the dense factor attractive for the active-set method: the new
pivot of a released column is exactly its Schur complement, so the singularity test of
\cref{chap:singular} reduces to a single number; removals and updates never decrease a pivot;
downdates fail exactly when the modified matrix is not positive definite; and the factor never
contains a stored entry that does not belong to the current matrix, so that the periodic
refactorization of \cref{sec:refactorization} becomes unnecessary.

\medskip

The two factorizations are interchangeable: the active-set method only requires the
operations of \cref{sec:ldl-updates}, and both provide them. The chapter closes with a
comparison of their costs; the measurements follow in \cref{sec:results-factorizations}.

\vspace*{\fill}

\newpage


\section{Why a Dense Factor}
\label{sec:why-dense}

Three observations from the previous chapters motivate a different representation of the
factor.

\paragraph{The factorized matrix is small.} By \cref{cor:dimension}, as long as the reduced
Hessian is positive definite,
\begin{equation} \label{eq:dense-size}
|\Nx|\;\le\;\rank(\Aeff)\;\le\;m .
\end{equation}
The release of a dependent column (\cref{prop:singular-cases}) adds one more index before it is
detected and undone, so at most $m+1$ indices are ever active at the same time. The dimension of
$\Heff$ is therefore bounded by the number of linking rows and does not depend on the column
capacity $\nmax$ or on the number of generated columns. The fixed-length representation of
\cref{sec:fixed-length}, in contrast, gives the factor the dimension $\nmax$, which is typically
one to two orders of magnitude larger than $m$.

\paragraph{The factorized matrix is dense.} The entry $(\Heff)_{jj'} = u_j^\top\Meff u_{j'}$
of~\eqref{eq:implicit} is nonzero as soon as the differences $u_j = a_j - a_{\rep(j)}$ and
$u_{j'}$ share a row of $\Aeff$. Two columns of the same block share their representative, so
\[
(\Heff)_{jj'} = a_j^\top\Meff a_{j'} - a_j^\top\Meff a_{S_k} - a_{S_k}^\top\Meff a_{j'}
+ a_{S_k}^\top\Meff a_{S_k}
\qquad(\psi(j) = \psi(j') = k),
\]
which vanishes only by cancellation. Columns of different blocks compete for the same linking
rows, which is the purpose of these rows, and therefore overlap as well. The measurements of
\cref{sec:results-fill} confirm that the factor of $\Heff$ is practically dense, in agreement
with \cref{rem:natural-ordering}. A sparse factorization of a dense matrix stores the same
numbers as a dense one, together with row indices and column pointers, and its operations
follow these index structures instead of streaming through contiguous memory.

\paragraph{Sparse updates leave structural zeros.} As shown in \cref{sec:refactorization}, row
deletions and rank-one modifications of a sparse factor keep the positions of the removed
entries in its sparsity pattern. Every later operation pays for them until the factor is
rebuilt, and the refactorization interval has to balance the two costs.

\medskip

A dense factor of the active block only has none of these drawbacks, provided that it can be
updated as cheaply as the sparse one. The rest of this chapter shows that it can.


\section{Compact Representation}
\label{sec:compact}

Throughout this chapter, $t = |\Nx|$ is the number of active indices, and
$\Nx = \{j_1,\dots,j_t\}$ lists them in the order in which they entered $\Nx$: $j_1$ has been
active the longest and $j_t$ is the most recent. The \textbf{position} of an active index is
$\operatorname{pos}(j_\ell) = \ell$.

\begin{definition}[Compact representation] \label{def:compact}
The \textbf{compact effective reduced Hessian} is the $t\times t$ matrix
\[
\hat H_{\ell\ell'} = (\Heff)_{j_\ell j_{\ell'}}\qquad(\ell,\ell' = 1,\dots,t),
\]
i.e.\ $\Heff$ with its rows and columns arranged in activation order. The \textbf{dense
factor} is its $LDL^\top$ factorization without pivoting,
\[
\hat H = LDL^\top,\qquad L\in\R^{t\times t}\ \text{unit lower triangular},\qquad
D = \diag(d_1,\dots,d_t),
\]
stored as a dense array together with the list $j_1,\dots,j_t$ and the position map
$\operatorname{pos}$.
\end{definition}

The factor exists and is unique whenever $\Heff$ is positive definite, and then all pivots
$d_\ell$ are positive \cite{golub2013matrix}. The only object of size $\nmax$ is the integer
array $\operatorname{pos}$, with the value $0$ for inactive indices; the floating-point data
occupy at most $(m+1)^2$ entries by~\eqref{eq:dense-size}. The compact representation solves
the same systems as the embedded one:

\begin{lemma}[Solves in compact form] \label{lem:compact-solve}
Let $\zeta\in\R^{\nmax}$ vanish outside $\Nx$, and let $\hat\zeta_\ell = \zeta_{j_\ell}$. If
$\hat\eta$ solves $\hat H\hat\eta = \hat\zeta$, then $\eta$ with $\eta_{j_\ell} = \hat\eta_\ell$ and
$\eta_j = 0$ for $j\notin\Nx$ solves $\Heff^{\mathrm{emb}}\eta = \zeta$.
\end{lemma}

\begin{proof}
By \cref{lem:embedding}, the solution of $\Heff^{\mathrm{emb}}\eta = \zeta$ vanishes outside
$\Nx$, and its restriction to $\Nx$ solves $\Heff\eta_x = \zeta_{\Nx}$. The matrix $\hat H$ is $\Heff$
with rows and columns permuted by the same permutation as the entries of $\hat\zeta$ and
$\hat\eta$, so $\hat H\hat\eta = \hat\zeta$ is this system written in activation order.
\end{proof}

A solve therefore consists of a gather of $\zeta$ into activation order, a forward substitution
with $L$, a division by the pivots, a backward substitution with $L^\top$ and a scatter of the
result, at the cost of $2t^2 + O(t)$ arithmetic operations, independent of $\nmax$.

\medskip

The activation order is not a fill-reducing ordering, and it need not be: the matrix is dense,
so every ordering of its rows produces a dense factor. It is chosen because it turns a row
addition into the cheapest possible operation, an append.

\begin{center}
\renewcommand{\arraystretch}{1.25}
\begin{tabular}{lll}
\toprule
change of the working set & change of $\hat H$ & dense operation \\
\midrule
(i) column leaves $\Nx$ & position $\operatorname{pos}(q)$ removed & delete (\cref{sec:dense-delete}) \\
(ii) representative swap & rows of $\mathcal{N}_x^{(k)}$ rewritten & deletes + appends \\
(iii) surplus becomes active & $+\,\hat v_i\hat v_i^\top$ & rank-one update (\cref{sec:dense-rank-one}) \\
(iv) column enters $\Nx$ & new last row and column & append (\cref{sec:dense-append}) \\
(v) surplus becomes free & $-\,\hat v_i\hat v_i^\top$ & rank-one downdate (\cref{sec:dense-rank-one}) \\
\bottomrule
\end{tabular}
\end{center}

Here $\hat v_i$ is the vector $v_i$ of~\eqref{eq:rank-one-up} in activation order. A
representative swap is carried out as in \cref{sec:ldl-updates}: the replacement $S_k'$ is
deleted, and every row of $\mathcal{N}_x^{(k)}\setminus\{S_k'\}$ is deleted and appended again with
its new vector $z_j' = z_j - z_{S_k'}$. These rows thereby move to the end of the activation order.


\section{Appending an Index}
\label{sec:dense-append}

When a column $r$ is released (case~(iv) of \cref{thm:H-reduce}), $\Heff$ is bordered by the
vector $\gamma_r$ and the scalar $\gamma_{rr}$ of~\eqref{eq:bordered}. In the compact
representation the new index takes position $t+1$, so the new matrix is
\[
\hat H' = \begin{bmatrix}\hat H & \hat\gamma\\ \hat\gamma^\top & \gamma_{rr}\end{bmatrix},
\qquad \hat\gamma_\ell = (\gamma_r)_{j_\ell}.
\]

\begin{theorem}[Append] \label{thm:append}
Let $\hat H = LDL^\top$ be positive definite. Then
\[
\hat H' = \begin{bmatrix} L & 0\\ l^\top & 1\end{bmatrix}
\begin{bmatrix} D & 0\\ 0 & d_{t+1}\end{bmatrix}
\begin{bmatrix} L & 0\\ l^\top & 1\end{bmatrix}^\top
\qquad\text{with}\qquad
Lw = \hat\gamma,\quad l = D^{-1}w,\quad d_{t+1} = \gamma_{rr} - w^\top D^{-1}w .
\]
The factor of $\hat H$ is not changed, and the new row costs one forward substitution,
$t^2 + O(t)$ operations.
\end{theorem}

\begin{proof}
Multiplying out the right-hand side gives the blocks
\[
LDL^\top = \hat H,\qquad LDl = Lw = \hat\gamma,\qquad
l^\top Dl + d_{t+1} = w^\top D^{-1}w + d_{t+1} = \gamma_{rr},
\]
which are the blocks of $\hat H'$. The first factor is unit lower triangular and the middle one
diagonal, so this is the $LDL^\top$ factorization of $\hat H'$ whenever $d_{t+1}\ne0$, by
uniqueness. The vector $w$ is obtained from $Lw = \hat\gamma$ by forward substitution, and $l$
and $d_{t+1}$ in $O(t)$ further operations.
\end{proof}

The new pivot has a simple meaning. It is exactly the quantity that decides whether the
release is singular.

\begin{cor}[The new pivot is the Schur complement] \label{cor:append-pivot}
The new pivot of \cref{thm:append} is the Schur complement of~\eqref{eq:bordered},
\[
d_{t+1} = \sigma_r = \gamma_{rr} - \gamma_r^\top\Heff^{-1}\gamma_r
= \bigl\lVert (I - \Pi_{\mathrm{eff}})\,\Aeff z_r\bigr\rVert^2 ,
\]
where $\Pi_{\mathrm{eff}}$ is the orthogonal projector onto the range of $\Aeff Z_x$. In particular
$d_{t+1}\ge0$, and $d_{t+1} = 0$ if and only if the release of $r$ is singular.
\end{cor}

\begin{proof}
From $Lw = \hat\gamma$ and $\hat H = LDL^\top$,
\[
w^\top D^{-1}w = \hat\gamma^\top L^{-\top}D^{-1}L^{-1}\hat\gamma = \hat\gamma^\top\hat H^{-1}\hat\gamma
= \gamma_r^\top\Heff^{-1}\gamma_r ,
\]
since the permutation to activation order does not change this quadratic form. Hence
$d_{t+1} = \sigma_r$. With $U = \Aeff Z_x$ and $u = \Aeff z_r$ we have $\Heff = U^\top U$,
$\gamma_r = U^\top u$ and $\gamma_{rr} = u^\top u$, so
\[
\sigma_r = u^\top u - u^\top U(U^\top U)^{-1}U^\top u = u^\top(I - \Pi_{\mathrm{eff}})\,u
= \lVert(I-\Pi_{\mathrm{eff}})\,u\rVert^2,
\]
because $\Pi_{\mathrm{eff}} = U(U^\top U)^{-1}U^\top$ is symmetric and idempotent. The last statement is
part~(i) of \cref{prop:singular-cases}.
\end{proof}

\begin{remark}[Singularity detection with the dense factor] \label{rem:dense-detection}
\Cref{cor:append-pivot} makes the detection of \cref{sec:detection} exact for released columns.
An append leaves all existing pivots unchanged, and the new pivot, scaled by the pivot scale
$\lVert u_r\rVert^2$ of \cref{def:pivot-invariant}, is
\[
\frac{d_{t+1}}{\lVert u_r\rVert^2} = \frac{\lVert(I-\Pi_{\mathrm{eff}})\,\Aeff z_r\rVert^2}{\lVert u_r\rVert^2}
\;\le\;\sin^2\theta_r ,
\]
where $\theta_r$ is the angle between $\Aeff z_r$ and the range of $\Aeff Z_x$ (with equality if no
row of $u_r$ is masked). The pivot invariant thus tests directly whether the new direction is
numerically independent of the old ones. In the sparse factorization of \cref{prop:rowadd} the
same information is spread over the new pivot and the downdated pivots of the later indices.
A dependent column can moreover be removed again by deleting the last position, which by
\cref{thm:delete} below costs no arithmetic at all. (Our implementation undoes it by a
refactorization for both factorizations, which happens only a handful of times per run.)
\end{remark}


\section{Rank-One Modifications}
\label{sec:dense-rank-one}

The surplus cases~(iii) and~(v) change $\Heff$ by $\pm v_iv_i^\top$, and we will see in
\cref{sec:dense-delete} that the deletion of an index requires a rank-one update as well. All
three are instances of the following problem: given the factor $\hat H = LDL^\top$, a vector
$v\in\R^t$ and a scalar $\varphi$, compute the factor of $\hat H + \varphi\,vv^\top$. We use the
method C1 of Gill, Golub, Murray and Saunders \cite{gill1974methods}, and derive it here in a
form that also yields the properties we need.

\begin{theorem}[Rank-one modification] \label{thm:rank-one}
Let $\hat H = LDL^\top$ with $d_1,\dots,d_t>0$, let $v\in\R^t$ and $\varphi\in\R$, and suppose that
$\hat H + \varphi\,vv^\top$ is positive definite. Define $\varphi_1 = \varphi$, $v^{(0)} = v$, and
for $j = 1,\dots,t$
\begin{equation} \label{eq:c1}
\chi_j = v^{(j-1)}_j,\qquad \bar d_j = d_j + \varphi_j\chi_j^2,\qquad
\beta_j = \frac{\varphi_j\chi_j}{\bar d_j},\qquad \varphi_{j+1} = \frac{d_j}{\bar d_j}\,\varphi_j,\qquad
v^{(j)} = v^{(j-1)} - \chi_j\,Le_j .
\end{equation}
Then $\hat H + \varphi\,vv^\top = \bar L\bar D\bar L^\top$ with $\bar D = \diag(\bar d_1,\dots,\bar d_t)$ and the
unit lower triangular matrix $\bar L$ whose columns are
\[
\bar Le_j = Le_j + \beta_j\,v^{(j)}\qquad(j = 1,\dots,t).
\]
The vector $v^{(j)}$ vanishes in its first $j$ entries, so column $j$ of $\bar L$ differs from
column $j$ of $L$ only below the diagonal.
\end{theorem}

\begin{proof}
\emph{Step 1: reduction to a diagonal matrix.} Let $\tilde v = L^{-1}v$. Then
\[
\hat H + \varphi\,vv^\top = L\,(D + \varphi\,\tilde v\tilde v^\top)\,L^\top .
\]
If $D + \varphi\tilde v\tilde v^\top = \tilde L\bar D\tilde L^\top$, then
$\hat H + \varphi vv^\top = (L\tilde L)\,\bar D\,(L\tilde L)^\top$, and $\bar L = L\tilde L$ is unit lower triangular
as a product of two such matrices. It therefore suffices to factorize the diagonal matrix plus
a rank-one term.

\emph{Step 2: factor of $D + \varphi\tilde v\tilde v^\top$.} We show by induction over $t$ that
\begin{equation} \label{eq:c1-inner}
\bar d_j = d_j + \varphi_j\tilde v_j^2,\qquad \tilde L_{rj} = \beta_j\tilde v_r\;\;(r>j),\qquad
\beta_j = \varphi_j\tilde v_j/\bar d_j,\qquad \varphi_{j+1} = \varphi_jd_j/\bar d_j .
\end{equation}
Split off the first row and column:
\[
D + \varphi\tilde v\tilde v^\top =
\begin{bmatrix} d_1 + \varphi\tilde v_1^2 & \varphi\tilde v_1\tilde v_2^\top\\
\varphi\tilde v_1\tilde v_2 & D_2 + \varphi\tilde v_2\tilde v_2^\top\end{bmatrix},
\]
with $\tilde v = [\tilde v_1;\tilde v_2]$ and $D = \diag(d_1, D_2)$. One step of symmetric Gaussian
elimination gives the pivot $\bar d_1 = d_1 + \varphi\tilde v_1^2$, which is positive because the
matrix is positive definite, the first column $\varphi\tilde v_1\tilde v_2/\bar d_1 = \beta_1\tilde v_2$ of
$\tilde L$ below the diagonal, and the Schur complement
\[
D_2 + \varphi\tilde v_2\tilde v_2^\top - \frac{\varphi^2\tilde v_1^2}{\bar d_1}\tilde v_2\tilde v_2^\top
= D_2 + \varphi\Bigl(1 - \frac{\varphi\tilde v_1^2}{\bar d_1}\Bigr)\tilde v_2\tilde v_2^\top
= D_2 + \underbrace{\varphi\,\frac{d_1}{\bar d_1}}_{=\varphi_2}\,\tilde v_2\tilde v_2^\top .
\]
The Schur complement of a positive definite matrix is positive definite and has the same form
as the original matrix, with $D_2$, $\tilde v_2$ and $\varphi_2$, so the induction hypothesis
gives~\eqref{eq:c1-inner} for the remaining indices.

\emph{Step 3: the columns of $\bar L$.} By Step~2, column $j$ of $\tilde L$ is
$e_j + \beta_j\sum_{r>j}\tilde v_re_r$, so
\[
\bar Le_j = L\tilde Le_j = Le_j + \beta_j\sum_{r>j}\tilde v_r\,Le_r .
\]
Since $L\tilde v = v$, the sum equals $v - \sum_{i\le j}\tilde v_i\,Le_i$. We show that this is
$v^{(j)}$ and that $\chi_j = \tilde v_j$. Indeed, $v^{(j)} = v - \sum_{i\le j}\chi_i\,Le_i$ by its
recursion. If $\chi_i = \tilde v_i$ for $i<j$, then $v^{(j-1)} = \sum_{r\ge j}\tilde v_r\,Le_r$, whose
$j$-th entry is $\tilde v_j$ because $L$ is unit lower triangular; hence $\chi_j = \tilde v_j$. By
induction, $\chi_j = \tilde v_j$ for all $j$, and $\bar Le_j = Le_j + \beta_jv^{(j)}$. The same
representation shows that $v^{(j)} = \sum_{r>j}\tilde v_rLe_r$ vanishes in its first $j$ entries.
Inserting $\chi_j = \tilde v_j$ into~\eqref{eq:c1-inner} gives~\eqref{eq:c1}.
\end{proof}

The recursion~\eqref{eq:c1} computes $\tilde v = L^{-1}v$ by forward substitution and builds the
new factor in the same pass, column by column. \Cref{alg:c1} states it in the form in which it
is implemented: the vector $v^{(j)}$ overwrites $v$, and the new column overwrites the old one.

\begin{algobox}{Rank-One Modification $LDL^\top + \varphi\,vv^\top$ (Method C1)}{alg:c1}
\State \textbf{Input:} $L$, $d$ of $\hat H = LDL^\top$, vector $v$, scalar $\varphi$
\State $j_0\gets$ first index with $v_{j_0}\ne0$ \Comment{columns before $j_0$ are unchanged}
\For{$j = j_0,\dots,t$}
    \State $\chi\gets v_j$; \textbf{if} $\chi = 0$ \textbf{then continue} \Comment{column $j$ and $\varphi$ unchanged}
    \State $\bar d\gets d_j + \varphi\chi^2$
    \If{$\bar d\le0$} \textbf{return} \emph{not positive definite} \EndIf
    \State $\beta\gets\varphi\chi/\bar d$; \;$\varphi\gets\varphi\,d_j/\bar d$; \;$d_j\gets\bar d$
    \For{$r = j+1,\dots,t$}
        \State $v_r\gets v_r - \chi\,L_{rj}$; \;$L_{rj}\gets L_{rj} + \beta\,v_r$
    \EndFor
\EndFor
\end{algobox}

If $\chi_j = 0$, then $\bar d_j = d_j$, $\beta_j = 0$ and $\varphi_{j+1} = \varphi_j$, and the column is
unchanged; this justifies the skip in \cref{alg:c1}. The cost is $2(t - j_0)^2 + O(t)$
operations, where $j_0$ is the position of the first nonzero of $v$. The following properties
of~\eqref{eq:c1} are used below.

\begin{prop}[Pivots of a rank-one modification] \label{prop:c1-pivots}
In the situation of \cref{thm:rank-one}:
\begin{enumerate}
    \item If $\varphi>0$, then $\varphi_j>0$ and $\bar d_j\ge d_j$ for all $j$: an update never
          decreases a pivot.
    \item The pivots change by the factor
          \[
          \prod_{j=1}^{t}\frac{\bar d_j}{d_j} = 1 + \varphi\,v^\top\hat H^{-1}v .
          \]
    \item If $\varphi<0$ and the recursion~\eqref{eq:c1} is run without the assumption that
          $\hat H + \varphi vv^\top$ is positive definite, then all $\bar d_j$ are positive if and only if
          $\hat H + \varphi vv^\top$ is positive definite, i.e.\ if and only if
          $1 + \varphi\,v^\top\hat H^{-1}v>0$.
\end{enumerate}
\end{prop}

\begin{proof}
(i) By induction: if $\varphi_j>0$, then $\bar d_j = d_j + \varphi_j\chi_j^2\ge d_j>0$ and
$\varphi_{j+1} = \varphi_jd_j/\bar d_j>0$.

(ii) The determinant of an $LDL^\top$ factorization is the product of its pivots, so
\[
\prod_j\frac{\bar d_j}{d_j} = \frac{\det(\hat H + \varphi vv^\top)}{\det\hat H}
= 1 + \varphi\,v^\top\hat H^{-1}v
\]
by the matrix determinant lemma \cite{golub2013matrix}.

(iii) The recursion performs symmetric Gaussian elimination without pivoting on
$D + \varphi\tilde v\tilde v^\top$ (Step~2 of the proof of \cref{thm:rank-one}). Elimination without
pivoting produces positive pivots $\bar d_1,\dots,\bar d_t$ if and only if all leading principal
minors are positive, i.e.\ if and only if the matrix is positive definite, and
$D + \varphi\tilde v\tilde v^\top$ is positive definite if and only if
$\hat H + \varphi vv^\top = L(D+\varphi\tilde v\tilde v^\top)L^\top$ is. Since $\hat H$ is positive definite and
$\hat H + \varphi vv^\top$ differs from it by a rank-one term, the latter is positive definite if and
only if its determinant is positive (its smallest eigenvalue is the only one that can become
nonpositive, by interlacing), i.e.\ by~(ii) if and only if $1 + \varphi v^\top\hat H^{-1}v>0$.
\end{proof}

For the two surplus cases, $\varphi = \pm1$ and $v = \hat v_i$. An activated surplus ($\varphi = 1$)
can never fail and only increases pivots. A released surplus ($\varphi = -1$) produces positive
pivots exactly when $\omega_i = 1 - v_i^\top\Heff^{-1}v_i>0$, i.e.\ exactly when its release is not
singular (part~(ii) of \cref{prop:singular-cases}); the product of the pivot ratios is $\omega_i$.
\Cref{alg:c1} stops as soon as a pivot is not positive, and the method then proceeds as in
\cref{alg:singular}.

\begin{remark}[Numerical behaviour] \label{rem:c1-stability}
For $\varphi>0$ no cancellation occurs in $\bar d_j = d_j + \varphi_j\chi_j^2$, and method~C1 is
numerically stable \cite{gill1974methods}. For $\varphi<0$ the subtraction can cancel, and the
error grows when $\omega_i$ is small, i.e.\ when the downdated matrix is nearly singular. This is
precisely the situation that the pivot invariant of \cref{def:pivot-invariant} detects, and it
is then resolved by the zero-curvature step of \cref{chap:singular}. Away from it, the factor
stays accurate over long runs; \cref{sec:results-factorizations} measures its residual.
\end{remark}


\section{Removing an Index}
\label{sec:dense-delete}

When the index $q\in\Nx$ leaves the nullspace set (case~(i) of \cref{thm:H-expand}), its row
and column are removed from $\hat H$. Let $\ell = \operatorname{pos}(q)$ and partition the
positions into the groups $1 = \{1,\dots,\ell-1\}$, $2 = \{\ell\}$ and $3 = \{\ell+1,\dots,t\}$:
\[
L = \begin{bmatrix} L_{11} & & \\ l_{21}^\top & 1 & \\ L_{31} & l_{32} & L_{33}\end{bmatrix},\qquad
D = \diag(D_1, d_\ell, D_3).
\]

\begin{theorem}[Delete] \label{thm:delete}
Removing row and column $\ell$ from $\hat H = LDL^\top$ gives the matrix with the factor
\[
\begin{bmatrix} L_{11} & 0\\ L_{31} & \bar L_{33}\end{bmatrix},\qquad
\diag(D_1, \bar D_3),\qquad\text{where}\quad
\bar L_{33}\bar D_3\bar L_{33}^\top = L_{33}D_3L_{33}^\top + d_\ell\,l_{32}l_{32}^\top .
\]
The trailing factor is a rank-one update with $\varphi = d_\ell>0$, computed by \cref{alg:c1}. The
blocks $L_{11}$, $D_1$ and $L_{31}$ are kept.
\end{theorem}

\begin{proof}
Multiplying out $LDL^\top$ block by block gives, for the blocks that remain after the removal,
\[
\hat H_{11} = L_{11}D_1L_{11}^\top,\qquad \hat H_{31} = L_{31}D_1L_{11}^\top,\qquad
\hat H_{33} = L_{31}D_1L_{31}^\top + d_\ell\,l_{32}l_{32}^\top + L_{33}D_3L_{33}^\top .
\]
The proposed factor has the blocks $L_{11}D_1L_{11}^\top = \hat H_{11}$,
$L_{31}D_1L_{11}^\top = \hat H_{31}$ and
\[
L_{31}D_1L_{31}^\top + \bar L_{33}\bar D_3\bar L_{33}^\top
= L_{31}D_1L_{31}^\top + L_{33}D_3L_{33}^\top + d_\ell\,l_{32}l_{32}^\top = \hat H_{33}.
\]
It is unit lower triangular with a diagonal middle factor, so by uniqueness it is the $LDL^\top$
factorization of the reduced matrix. The trailing matrix $L_{33}D_3L_{33}^\top + d_\ell l_{32}l_{32}^\top$
is positive definite because $d_\ell>0$, so \cref{thm:rank-one} applies with $\varphi = d_\ell$.
\end{proof}

\begin{figure}[ht]
\centering
\begin{tikzpicture}[x=0.36cm, y=0.36cm, font=\small]
    % before: lower triangle of L, t = 10, deleted position ell = 4
    \fill[cblue!25] (0,10) -- (0,0) -- (10,0) -- cycle;
    \fill[cred!35] (3,0) rectangle (4,7);                % column ell below the diagonal (l_32)
    \fill[cred!35] (0,6) -- (0,7) -- (3,7) -- (3,6) -- cycle; % row ell (l_21)
    \fill[cgreen!35] (4,0) -- (4,6) -- (10,0) -- cycle;   % L_33
    \draw[thick] (0,10) -- (0,0) -- (10,0) -- cycle;
    \draw (3,0) -- (3,7); \draw (4,0) -- (4,6); \draw (0,6) -- (4,6); \draw (0,7) -- (3,7);
    \node at (1.2,3) {$L_{31}$};
    \node at (6.2,1.6) {$L_{33}$};
    \node[cred!70!black, above] at (3.5,7) {$\ell$};
    \node[below] at (5,-0.3) {$L$ before};
    % arrow
    \draw[-{Stealth[length=3mm]}, thick] (11.2,4) -- (14.8,4)
        node[midway, above, align=center] {\scriptsize delete $\ell$};
    % after: t = 9
    \begin{scope}[xshift=5.8cm]
    \fill[cblue!25] (0,9) -- (0,0) -- (9,0) -- cycle;
    \fill[cgreen!35] (3,0) -- (3,6) -- (9,0) -- cycle;
    \draw[thick] (0,9) -- (0,0) -- (9,0) -- cycle;
    \draw (3,0) -- (3,6); \draw (0,6) -- (3,6);
    \node at (1.2,3) {$L_{31}$};
    \node at (5.4,1.6) {$\bar L_{33}$};
    \node[below] at (4.5,-0.3) {$L$ after};
    \end{scope}
    % legend
    \begin{scope}[xshift=10.2cm, yshift=0.6cm]
    \fill[cred!35] (0,7) rectangle (0.7,7.7); \node[right] at (0.9,7.35) {\scriptsize removed: $l_{21}$, $d_\ell$, $l_{32}$};
    \fill[cgreen!35] (0,5.6) rectangle (0.7,6.3); \node[right] at (0.9,5.95) {\scriptsize updated by $d_\ell\,l_{32}l_{32}^\top$};
    \fill[cblue!25] (0,4.2) rectangle (0.7,4.9); \node[right] at (0.9,4.55) {\scriptsize kept, rows below $\ell$ shifted up};
    \end{scope}
\end{tikzpicture}
\caption{Deleting position $\ell$ from the dense factor (\cref{thm:delete}). The leading
columns are kept, the trailing triangle receives a rank-one update with the removed column
$l_{32}$ and the positive weight $d_\ell$, and all positions behind $\ell$ move up by one, so
that the factor stays contiguous.}
\label{fig:dense-delete}
\end{figure}

In an implementation, the update of the trailing block and the compaction are one pass:
\cref{alg:c1} reads column $\ell+s$ of the old trailing block and writes the updated column to
position $\ell+s-1$, which has been read completely in the previous step
(\cref{fig:dense-delete}). The rows $\ell+1,\dots,t$ of the leading columns are moved up by one
position, and the activation list and the position map are updated.

\begin{prop}[Cost of a deletion] \label{prop:delete-cost}
Deleting position $\ell$ costs $2(t-\ell)^2 + O(t)$ arithmetic operations for the rank-one update
and $O\bigl(\ell\,(t-\ell)\bigr)$ memory moves for the compaction. Deleting the most recent index
($\ell = t$) costs no arithmetic; deleting the oldest one ($\ell = 1$) costs as much as a full
rank-one update.
\end{prop}

\begin{proof}
The trailing block has $t-\ell$ positions, and \cref{alg:c1} on it costs $2(t-\ell)^2 + O(t-\ell)$
operations. The block $L_{31}$ has $t-\ell$ rows and $\ell-1$ columns, all of which are moved.
\end{proof}

The cost of a deletion therefore depends on how long the removed index has been active. In the
active-set method most removed indices have entered $\Nx$ recently: a released column moves
away from its bound (\cref{lem:descent}), and the columns that block are typically those with
small weights, which have been released last. The measurements of
\cref{sec:results-factorizations} show that a deletion touches about a quarter of the factor on
average.

\begin{cor}[Expansions preserve the pivot invariant] \label{cor:expand-pivots}
An expansion of type~(i) or~(iii) of \cref{lem:expand-ordering} does not decrease any pivot of
an index that stays in $\Nx$. In particular, if the pivot invariant of
\cref{def:pivot-invariant} holds before such an expansion, it holds after it.
\end{cor}

\begin{proof}
In case~(i), the pivots $D_1$ are kept and $\bar D_3$ results from a rank-one update with
$\varphi = d_\ell>0$, so $\bar D_3\ge D_3$ entrywise by part~(i) of \cref{prop:c1-pivots}; each pivot
moves with its index to the position one up. In case~(iii), the update $+\hat v_i\hat v_i^\top$ has
$\varphi = 1>0$. The pivot scales $\lVert u_j\rVert^2$ of the remaining indices do not change, since
their representatives do not change.
\end{proof}

In exact arithmetic the same holds for the sparse factorization, whose row deletion is the same
rank-one update (\cref{sec:ldl-updates}). \Cref{cor:expand-pivots} shows that the pivot
invariant can only be violated by a reduction, which is where \cref{chap:singular} checks it.


\section{Structural Zeros and Refactorization}
\label{sec:dense-no-zeros}

The dense factor stores exactly the $t(t-1)/2$ entries below the diagonal of the $t\times t$
array, the $t$ pivots, and the activation list. This is all that it ever stores:

\begin{prop}[No structural zeros] \label{prop:no-zeros}
After every sequence of appends, deletions and rank-one modifications, the leading
$t\times t$ block of the stored array is the factor $L$ of the current matrix $\hat H$, and no
stored entry belongs to a matrix that is no longer current.
\end{prop}

\begin{proof}
By induction over the operations. \Cref{thm:append} writes the new last row and pivot and keeps
the leading block. \Cref{thm:rank-one} overwrites the columns from $j_0$ on and keeps the
others, without changing the dimension. \Cref{thm:delete} keeps the leading columns, overwrites
the trailing triangle with $\bar L_{33}$, and moves everything behind position $\ell$ up by one,
so that the result again occupies exactly the leading $(t-1)\times(t-1)$ block. Entries outside
the leading block are never read, and the next append overwrites the entries of its row before
using them.
\end{proof}

The cost of every operation therefore depends only on the current $t$, and the factor never has
to be rebuilt to restore its efficiency: there is nothing like the growth of the stored pattern
of \cref{sec:refactorization}. A refactorization is only needed for numerical reasons, namely
when a pivot violates the invariant of \cref{def:pivot-invariant} or a downdate fails
(\cref{alg:singular}). Such a refactorization appends the active indices one by one as in
\cref{def:refactorization}, at the cost of about $t^3/3$ operations
(\cref{prop:refactor-cost}); afterwards their activation order is the increasing index order. An
optional periodic refactorization could be used to bound the accumulation of rounding errors,
but \cref{sec:results-factorizations} shows that it is not needed in practice.

\begin{remark}[Reset] \label{rem:dense-reset}
Resetting the factor to the identity only clears the position map of the active indices, in
$O(t)$. The stored entries need not be zeroed, because every append writes its whole row before
it is read.
\end{remark}


\section{Comparison with the Sparse Factorization}
\label{sec:dense-comparison}

\Cref{tab:dense-vs-sparse} compares the operations of the two factorizations. For the sparse
factorization, the costs are those of a factor without structural zeros, i.e.\ directly after a
refactorization; between refactorizations they grow with the stored pattern.

\begin{table}[ht]
\centering
\small
\renewcommand{\arraystretch}{1.3}
\begin{tabular}{>{\raggedright}p{0.2\textwidth}>{\raggedright}p{0.36\textwidth}>{\raggedright\arraybackslash}p{0.36\textwidth}}
\toprule
operation & sparse factor of $\Heff^{\mathrm{emb}}$ (\cref{chap:updates}) & dense factor of $\hat H$ (this chapter) \\
\midrule
dimension & $\nmax$, identity outside $\Nx$ & $t = |\Nx|\le m$ \\
release of a column, (iv) & row addition at index $r$: solve, new column, downdate of the later indices & append: one forward substitution, $t^2$ \\
removal of a column, (i) & row deletion: rank-one update of the later indices & delete: rank-one update of the later positions, $2(t-\ell)^2$ \\
surplus, (iii) and (v) & rank-one update or downdate along the elimination tree & \cref{alg:c1} from the first nonzero, at most $2t^2$ \\
solve & sparse solve on $\reach_L(\Nx)$ & two triangular solves, $2t^2$ \\
new pivot of a release & $D_{22}'\ge\sigma_r$; later pivots change & $d_{t+1} = \sigma_r$; no other pivot changes \\
stored entries & grow with every update (\cref{sec:refactorization}) & exactly $t(t-1)/2$ (\cref{prop:no-zeros}) \\
periodic refactorization & required, every $\tau$ updates & not required \\
floating-point memory & proportional to the stored pattern & at most $(m+1)^2$ \\
\bottomrule
\end{tabular}
\caption{The two factorizations of the effective reduced Hessian. Costs are arithmetic
operations; $\ell$ is the position of the removed index.}
\label{tab:dense-vs-sparse}
\end{table}

For a dense $\Heff$, both factorizations perform the same arithmetic up to constant factors. The
dense one does it on contiguous memory with simple loops and triangular solves, without index
lookups, without an $\nmax$-dimensional embedding and without structural zeros. On the other
hand, the dense factor costs $(m+1)^2$ floating-point numbers and $O(t^2)$ operations per update
regardless of the sparsity of $\Heff$. A sparse factor remains preferable when $\Heff$ is
genuinely sparse, for instance when the blocks have disjoint sets of linking rows, and $m$ is
large. For the crew scheduling instances of \cref{chap:results}, where $\Heff$ is dense, the
dense factorization is the better choice, as \cref{sec:results-factorizations} shows.

\medskip

Nothing in \cref{chap:qpas,chap:nullspace,chap:singular} depends on the choice. The
active-set method, the hot start of \cref{thm:hot-start} and the singularity treatment of
\cref{alg:singular} only use the operations of the table, so the two factorizations
produce the same sequence of iterates in exact arithmetic, and \cref{alg:hot-qpas} applies
unchanged, with the refactorization in its last line triggered only by the pivot invariant.

\chapter{Computational Results}
\label{chap:results}

\vspace*{\fill}

This chapter evaluates the hot-started nullspace active-set method on generated instances of
the \textbf{airline crew scheduling problem}. Every instance is solved in two variants: with
equality coverage rows only, and with 30\% of the coverage rows relaxed to ``$\ge$''
inequalities. The active-set method is run with both factorizations of the effective reduced
Hessian: the sparse factorization of \cref{chap:updates} and the dense factorization of
\cref{chap:dense}.

\medskip

We first check \textbf{correctness}: every column generation method must reproduce the
optimum of the full linear program, which is still small enough to be solved directly. We
then study the \textbf{convergence behaviour} of regularized column generation, which
consists of two clearly separated kinds of iterations, and the \textbf{computational time},
broken down by solver, by iteration and by operation inside the active-set method, up to
instances with almost a thousand linking rows.

\medskip

The last sections look inside the method. We follow the sizes of the nullspace sets over a run
with inequality rows, measure the \textbf{contents of the sparse factorization} and the value
of periodic refactorization, and finally compare the \textbf{two factorizations} operation by
operation, in memory and in accuracy.

\vspace*{\fill}

\newpage


\section{Test Problem: Airline Crew Scheduling}
\label{sec:test-problem}

\subsection{The model}

In the crew scheduling problem, $K$ pilots have to be assigned schedules that together
operate a set of $m$ flights (\cref{ex:crew}). The schedules of a pilot are the source--sink
paths in a \textbf{time--space network}: a directed acyclic graph $G = (V,E)$ whose nodes
represent states of a pilot at given times, and whose edges represent the allowed transitions
between them. The edge costs of pilot $k$ form the vector $c^{(k)}\in\R^{|E|}$, and every edge
operates a subset of the flights, recorded in the \textbf{task matrix}
$A_{\mathrm{task}}\in\{0,1\}^{m\times|E|}$. With $\xi_k\in\{0,1\}^{|E|}$ the edge vector of the
path of pilot $k$ and $F$ the node--edge incidence matrix of $G$ ($F_{\mathrm{tail}(e),e} = 1$,
$F_{\mathrm{head}(e),e} = -1$), the linear relaxation reads
\begin{subequations} \label{eq:crew}
\begin{align}
\min_{\xi_1,\dots,\xi_K}\quad & \sum_{k=1}^{K} c^{(k)\top} \xi_k\\
\text{s.t.}\quad & \Bigl(\sum_{k=1}^{K}A_{\mathrm{task}}\,\xi_k\Bigr)_i = 1 && (i\in\mathcal{R}_E),
\label{eq:crew-eq}\\
& \Bigl(\sum_{k=1}^{K}A_{\mathrm{task}}\,\xi_k\Bigr)_i \ge 1 && (i\in\mathcal{R}_I),
\label{eq:crew-ineq}\\
& F \xi_k = b^{(k)},\quad 0\le \xi_k\le 1 && (k = 1,\dots,K),
\end{align}
\end{subequations}
where $b^{(k)}$ has the entry $+1$ at the source, $-1$ at the sink and $0$ elsewhere, and
$\mathcal{R}_E$, $\mathcal{R}_I$ are the sets of equality and inequality coverage rows. This is a
multi-block linear program (\cref{def:multiblock}) with identical block polytopes. A
coverage row in $\mathcal{R}_E$ requires a flight to be operated exactly once; a row in
$\mathcal{R}_I$ requires it at least once, which allows a pilot to travel on a flight as a
passenger. The columns of the master problem are complete schedules, and the pricing
problem~\eqref{eq:pricing} of pilot $k$ is a shortest path problem with the edge weights
$c^{(k)} - A_{\mathrm{task}}^\top\pi$. These weights can be negative, but the graph is
acyclic, so one pass of dynamic programming in topological order solves it in
$O(|V|+|E|)$.

\subsection{Instance generator}

The instances are generated randomly but reproducibly.
\begin{itemize}
    \item \textbf{Network.} The nodes $1,\dots,|V|$ are ordered in time. A backbone edge
          connects every node to the next one, which guarantees a source--sink path. A skip
          edge from a node to any node at least two steps later exists independently with a
          fixed probability. An edge that spans $\ell$ steps costs $\ell + U(0,1)$, multiplied
          by $1 + 0.1\,U(0,1)$ for each pilot separately, where $U(0,1)$ is a uniform random
          number.
    \item \textbf{Flights.} Each skip edge operates a random set of up to ten flights whose
          position in time is close to that of the edge. A flight is only attached to edges
          that overlap in time, so that a single schedule operates every flight at most once.
          The task matrix therefore has a banded structure.
    \item \textbf{Feasibility.} Only flights that can be covered exactly once in the linear
          relaxation are kept, so that the equality variant is feasible.
    \item \textbf{Inequality rows.} In the inequality variant, a random 30\% of the coverage
          rows are turned into ``$\ge1$'' rows.
\end{itemize}

\Cref{tab:instances} lists the six instance sizes. The instances I1--I4 are used for the
comparison of all methods; S5 and S6, with up to 960 coverage rows, extend the study of
scaling. The last column is the number of variables of the full edge-based
formulation~\eqref{eq:crew}, which HiGHS solves directly to provide the reference optimum
$z^\star$.

\begin{table}[ht]
\centering
\renewcommand{\arraystretch}{1.2}
\begin{tabular}{lrrrrrr}
\toprule
instance & $|V|$ & $|E|$ & $K$ & $m$ & $m_I$ (inequality variant) & variables of~\eqref{eq:crew} \\
\midrule
\tabinput{data/tab_instances.tex}
\bottomrule
\end{tabular}
\caption{Test instances. $K$ is the number of pilots, $m$ the number of coverage rows, and
$m_I$ the number of ``$\ge$'' rows in the inequality variant.}
\label{tab:instances}
\end{table}


\section{Implementation and Setup}
\label{sec:setup}

All methods are implemented in Julia~1.12 \cite{bezanson2017julia} in one framework: the
column generation driver, the pricing algorithm, the column pool and the statistics are
shared, and only the solver of the restricted master problem differs. We compare six
methods:

\begin{center}
\renewcommand{\arraystretch}{1.2}
\resizebox{\textwidth}{!}{%
\begin{tabular}{ll}
\toprule
method & master problem \\
\midrule
Exact LP (HiGHS) & none: formulation~\eqref{eq:crew} solved directly by the HiGHS simplex method \\
LP CG (HiGHS) & classical linear master problem, HiGHS simplex, warm-started \\
Explicit QP CG (HiGHS) & explicit regularized master problem~\eqref{eq:explicit}, HiGHS QP solver \\
Eliminated QP CG (HiGHS) & eliminated regularized master problem~\eqref{eq:eliminated}, HiGHS QP \\
QPAS CG, sparse factor & eliminated regularized master problem, \cref{alg:hot-qpas} with the factorization of \cref{chap:updates} \\
QPAS CG, dense factor & eliminated regularized master problem, \cref{alg:hot-qpas} with the factorization of \cref{chap:dense} \\
\bottomrule
\end{tabular}
}
\end{center}

HiGHS \cite{huangfu2018parallelizing} (version 1.15, with its simplex method for the LPs and
its active-set QP solver for the QPs) is used through a persistent model, to which new columns
are added incrementally, so that its LP solves are warm-started. The LP master starts with
artificial columns of cost $M = 10^4$ on both sides of every linking row (\cref{sec:column-generation}),
which keep it feasible until enough columns are generated. The eliminated QP master is
run on the equality variant only. The sparse
factorization of the active-set method uses CHOLMOD \cite{chen2008cholmod} from SuiteSparse~7.10
for the factorization and its updates. The dense factorization is our own implementation of
\cref{chap:dense}; its triangular solves call the BLAS routine \code{trsv} of OpenBLAS, and
its rank-one modifications are plain loops over the columns of the stored array.

\medskip

All column generation runs use the same settings: the initial columns are the cheapest path
of every pilot, pricing visits the pilots in random order and adds the first improving
column, the proximal point rule moves the center after every iteration, and $\delta = 1$,
$\varepsilon_{\mathrm{c}} = 10^{-6}$, $\varepsilon_{\mathrm{rc}} = 10^{-6}$. The active-set
method uses the column capacity $\nmax = 20\,000$ and the pivot tolerance
$\varepsilon_{\mathrm{piv}} = 10^{-10}$ with both factorizations. The sparse factorization is
refactorized every $\tau = 150$ updates unless stated otherwise; the dense factorization is
only refactorized when the pivot invariant is violated (\cref{sec:dense-no-zeros}).

\medskip

The HiGHS QP masters are slow, so they are run to completion on I1 only; on the equality
variant of I2 they are stopped after 120 seconds, which is enough to compare their path and
their time per master problem with the active-set method. The LP master is run on I1--I3 with
a limit of 900 seconds, and the direct LP solve of S5 and S6 is stopped after 300 seconds. All
experiments ran single-threaded on an AMD Ryzen~7 7700 processor with 32\,GB of memory under
Windows~11. Times of the active-set method are the faster of two runs, except for runs longer
than one minute on S5 and S6, which were run once; all other times are from a single run.


\section{Correctness}
\label{sec:results-correctness}

\Cref{tab:comparison-eq,tab:comparison-ineq} compare all methods on both variants. Besides the
status, the number of column generation iterations, the number of generated columns and the
times, they report three accuracy measures relative to the reference optimum $z^\star$:
\begin{itemize}
    \item the relative error of the cost, $(c^\top\lambda - z^\star)/|z^\star|$;
    \item the relative gap of the best Lagrangian bound, $(z^\star - \theta)/|z^\star|$, which
          is nonnegative for every valid bound (\cref{prop:lagrangian-bound});
    \item the infeasibility $\lVert b - A\lambda\rVert_\infty$ of the final weights in the linking
          rows (for inequality rows only the violated part).
\end{itemize}

\begin{table}[!htbp]
\centering
\scriptsize
\renewcommand{\arraystretch}{1.15}
\setlength{\tabcolsep}{4pt}
\resizebox{\textwidth}{!}{%
\begin{tabular}{llcrrrrccc}
\toprule
 & method & status & CG iter. & columns & time [s] & master [ms/iter.] & cost error & bound gap & infeasibility \\
\midrule
\tabinput{data/tab_comparison_eq.tex}
\bottomrule
\end{tabular}
}
\caption{Equality variant: all methods. ``time'' denotes a method stopped by its time limit
(120\,s for the HiGHS QP masters on I2). Master times are averages per column generation
iteration.}
\label{tab:comparison-eq}
\end{table}

\begin{table}[!htbp]
\centering
\scriptsize
\renewcommand{\arraystretch}{1.15}
\setlength{\tabcolsep}{4pt}
\resizebox{\textwidth}{!}{%
\begin{tabular}{llcrrrrccc}
\toprule
 & method & status & CG iter. & columns & time [s] & master [ms/iter.] & cost error & bound gap & infeasibility \\
\midrule
\tabinput{data/tab_comparison_ineq.tex}
\bottomrule
\end{tabular}
}
\caption{Inequality variant (30\% ``$\ge$'' rows). The eliminated HiGHS QP master is not
run on this variant, and the explicit HiGHS QP master is run on I1 only.}
\label{tab:comparison-ineq}
\end{table}

\paragraph{All methods reach the optimum.} Every column generation run that finishes within
its time limit stops with an optimal status, and its Lagrangian bound agrees with $z^\star$
to within $2\cdot10^{-7}$ relative or better. This confirms the exact termination criterion
of \cref{thm:exact-termination} in practice, including the sign conditions of the inequality
rows: the bound is evaluated at the projected duals $\pi^+$, since at duals that violate them it
would not be valid and could exceed $z^\star$ (\cref{sec:bound-ineq,rem:bound-sign}).

\paragraph{The accuracy of the stabilized methods.} The stabilized methods stop when
$\lVert\pi-\pihat\rVert_\infty\le\varepsilon_{\mathrm{c}}\max(1,\lVert\pi\rVert_\infty)$, which by
\cref{prop:stop-accuracy} allows an infeasibility of up to $\delta\varepsilon_{\mathrm{c}}\max(1,\lVert\pi\rVert_\infty)$
in the linking rows. The infeasibility column shows exactly this order of magnitude, about
$10^{-5}$ for the larger instances, where the duals are of order $10$ to $100$. The cost
$c^\top\lambda$ of such a slightly infeasible solution can be off by more than the
infeasibility itself: on I2 it is $2\cdot10^{-5}$ below $z^\star$. The Lagrangian bound, which
certifies the optimal value, is unaffected and agrees with $z^\star$ to about $2\cdot10^{-7}$. A
tighter $\varepsilon_{\mathrm{c}}$ removes the cost error at the price of a few more center-only iterations. In a
test on a smaller instance with $\delta = 5$, the cost error dropped from $10^{-6}$ for
$\varepsilon_{\mathrm{c}} = 10^{-6}$ to $10^{-15}$ for $\varepsilon_{\mathrm{c}} = 10^{-9}$, with one additional iteration.

\paragraph{The active-set method and HiGHS solve the same problem.} On instance I1, where
all stabilized methods finish, the explicit and eliminated HiGHS QPs and the active-set method
perform almost the same number of column generation iterations and generate almost the same
number of columns: they follow the same sequence of master problems, as they should. The small
differences come from the accuracy of the master solves, which affects the duals passed to
pricing and hence the order in which columns are found.

\paragraph{The two factorizations solve the same problem.} The two variants of the active-set
method differ only in the linear algebra, so in exact arithmetic they produce the same iterates
(\cref{sec:dense-comparison}). In floating point they agree closely: their final objectives
differ by at most $6\cdot10^{-7}$ relative on all twelve runs (\cref{tab:backends} below),
with the same accuracy against $z^\star$, and on I2 they perform exactly the same 11\,974
iterations. Where the iteration counts differ, rounding differences have broken a tie in a
ratio test or a pricing round differently, after which the runs continue on slightly
different paths. The differences are below 7\% on all runs but one: on the inequality
variant of S6 the dense variant needs 9\,574 iterations against 6\,926, with the same final
objective.


\section[Convergence and the Two Phases]{Convergence and the Two Phases of Regularized Column Generation}
\label{sec:results-phases}

\Cref{fig:convergence} shows the convergence of the master objective on instance I2 for all
column generation methods, together with the change of the duals from one iteration to the
next. Here and in the rest of this section, ``QPAS'' denotes the active-set method with the
dense factorization; the sparse one follows the same path (it performs the same iterations on
I2).

\begin{figure}[ht]
\centering
\begin{tikzpicture}
\begin{groupplot}[group style={group size=2 by 1, horizontal sep=1.6cm},
    width=0.5\textwidth, height=0.42\textwidth, ymode=log,
    xlabel={CG iteration}, tick label style={font=\scriptsize}, label style={font=\small},
    title style={font=\small}, legend style={font=\scriptsize, at={(0.98,0.98)}, anchor=north east},
    every axis plot/.append style={line width=0.7pt}]
\nextgroupplot[title={relative error of the master objective}, ymin=1e-9, ymax=1]
    \addplot[cred] table[x=it, y=gap, col sep=comma] {data/conv_I2_0_lp.csv};
    \addplot[cgreen] table[x=it, y=gap, col sep=comma] {data/conv_I2_0_explicit_qp.csv};
    \addplot[cpurple, densely dashed] table[x=it, y=gap, col sep=comma] {data/conv_I2_0_eliminated_qp.csv};
    \addplot[cblue] table[x=it, y=gap, col sep=comma] {data/conv_I2_0_qpas_dense.csv};
    \legend{LP CG, explicit QP, eliminated QP, QPAS}
\nextgroupplot[title={dual change $\lVert\pi_t - \pi_{t-1}\rVert_\infty$}, ymin=1e-6]
    \addplot[cred, only marks, mark size=0.3pt] table[x=it, y=dpi, col sep=comma] {data/conv_I2_0_lp.csv};
    \addplot[cblue, only marks, mark size=0.3pt] table[x=it, y=dpi, col sep=comma] {data/conv_I2_0_qpas_dense.csv};
    \legend{LP CG, QPAS}
\end{groupplot}
\end{tikzpicture}
\caption{Instance I2 (equality variant). Left: relative distance of the master objective to
$z^\star$. Right: change of the duals between consecutive iterations. The two HiGHS QP masters
stop at their time limit of 120 seconds after about 2\,000 iterations.}
\label{fig:convergence}
\end{figure}

\paragraph{Stabilization.} The duals of the linear master problem jump by up to $2\cdot10^4$
in the first iterations, when the artificial columns of cost $M = 10^4$ are still active, and
by about $2$ in a typical iteration of the second half of the run. The regularized duals change
by about $10^{-4}$ in a typical iteration of the second half. The linear master problem,
however, reaches the optimum in fewer iterations on this instance. Its jumps explore the dual
space quickly, while the proximal term keeps the regularized duals close to the center, so
that they need many small steps. Regularization therefore does not reduce the number of
iterations here. What it buys is iterations whose master problems differ only slightly from
each other, which is exactly what the hot start exploits: an iteration of the active-set
method costs $0.08$\,ms on average with the dense factorization and $0.19$\,ms with the
sparse one, against $29$\,ms for the warm-started LP master (\cref{tab:comparison-eq}).

\paragraph{The QP masters follow the same path.} The curves of the explicit QP, the
eliminated QP and the active-set method are almost indistinguishable over the first 2\,000
iterations, after which the HiGHS QP masters run out of time.

\medskip

The right panel of \cref{fig:convergence} shows a second feature of the regularized methods:
from some point on, the duals change by very small amounts in most iterations. These are the
center-only iterations of \cref{sec:center-update}. \Cref{fig:phases} shows how the two kinds
of iterations are distributed over a run.

\begin{figure}[ht]
\centering
\begin{tikzpicture}
\begin{groupplot}[group style={group size=2 by 2, horizontal sep=1.6cm, vertical sep=0.6cm,
    xlabels at=edge bottom, xticklabels at=edge bottom},
    width=0.5\textwidth, height=0.3\textwidth, xlabel={CG iteration},
    tick label style={font=\scriptsize}, label style={font=\small}, title style={font=\small},
    every axis plot/.append style={line width=0.7pt}]
\nextgroupplot[title={I2, equality variant}, ylabel={columns}]
    \addplot[cblue] table[x=it, y=ncol, col sep=comma] {data/conv_I2_0_qpas_dense.csv};
\nextgroupplot[title={I2, inequality variant}]
    \addplot[cblue] table[x=it, y=ncol, col sep=comma] {data/conv_I2_30_qpas_dense.csv};
\nextgroupplot[ylabel={share adding a column}, ymin=0, ymax=1.05]
    \addplot[const plot mark mid, fill=cblue!35, draw=cblue] table[x=it, y=share, col sep=comma]
        {data/phases_I2_0.csv} \closedcycle;
\nextgroupplot[ymin=0, ymax=1.05]
    \addplot[const plot mark mid, fill=cblue!35, draw=cblue] table[x=it, y=share, col sep=comma]
        {data/phases_I2_30.csv} \closedcycle;
\end{groupplot}
\end{tikzpicture}
\caption{The two phases of regularized column generation (active-set method, instance I2).
Top: number of columns. Bottom: share of the iterations in each window of 100 iterations
that add a column; the remaining iterations are center-only iterations.}
\label{fig:phases}
\end{figure}

\paragraph{Two phases.} A run starts with a \textbf{column phase}, in which (almost) every
iteration adds a column. It is followed by a \textbf{proximal phase}, in which most
iterations are center-only iterations that move the center towards the restricted dual
optimum, interrupted by occasional iterations that find a new column because the duals have
moved. The phases are not strictly separated: columns keep appearing far into the proximal
phase.
\Cref{tab:phases} quantifies this for all instances.

\begin{table}[ht]
\centering
\small
\renewcommand{\arraystretch}{1.15}
\setlength{\tabcolsep}{4pt}
\resizebox{\textwidth}{!}{%
\begin{tabular}{lcrrrrrrrr}
\toprule
 & & & & & first & last & \multicolumn{2}{c}{mean inner iter.} & master time \\
instance & ineq. & CG iter. & with column & center-only & center-only & column & column & center-only & center-only \\
\midrule
\tabinput{data/tab_phases.tex}
\bottomrule
\end{tabular}
}
\caption{Column-adding and center-only iterations of the active-set method (dense
factorization). ``inner iter.'' counts the iterations of the active-set method per master
problem.}
\label{tab:phases}
\end{table}

The two kinds of iterations differ sharply in cost. A column-adding iteration needs on
average tens of active-set iterations, since the new column is released and the working set
adapts to it. A center-only iteration needs at most two on average: as shown in
\cref{prop:center-only}, once the working set has settled, the direction is zero, and the
master problem is solved by a single multiplier test. The length of the proximal phase
varies a lot between instances. On I2 and I3 (equality variant) center-only iterations are
the large majority, on I4 they are a minority. On all instances except the tiny I1 their share
of the master time is much smaller than their share of the iterations.

\Cref{fig:optimality} shows how the optimality measures evolve over the two phases for the
active-set method on I2. In the column phase, the cost error, the error of the master
objective and the infeasibility decrease together. In the proximal phase, the Lagrangian
bound becomes available, since all pilots are priced in every center-only iteration. The
bound approaches $z^\star$ from below at the same rate at which the master objective
approaches it, so the two errors are indistinguishable in the plot. The infeasibility, and
with it the cost error, decrease more slowly and in steps: each step is a change of the
working set that ends one of the dual walks of \cref{prop:center-only}.

\begin{figure}[ht]
\centering
\begin{tikzpicture}
\begin{axis}[width=0.78\textwidth, height=0.42\textwidth, ymode=log, xlabel={CG iteration},
    ymin=1e-9, ymax=10, tick label style={font=\scriptsize}, label style={font=\small},
    legend style={font=\scriptsize, at={(0.98,0.98)}, anchor=north east, legend columns=2},
    every axis plot/.append style={line width=0.7pt}]
    \addplot[cred!45, only marks, mark=*, mark size=0.7pt] table[x=it, y=bgap, col sep=comma] {data/bound_I2_0_qpas_dense.csv};
    \addplot[corange] table[x=it, y=costgap, col sep=comma] {data/conv_I2_0_qpas_dense.csv};
    \addplot[cgreen] table[x=it, y=infeas, col sep=comma] {data/conv_I2_0_qpas_dense.csv};
    \addplot[cblue, densely dashed] table[x=it, y=gap, col sep=comma] {data/conv_I2_0_qpas_dense.csv};
    \legend{bound gap $(z^\star-\theta)/|z^\star|$, cost error $|c^\top\lambda - z^\star|/|z^\star|$,
            infeasibility, master objective error}
\end{axis}
\end{tikzpicture}
\caption{Optimality measures of the active-set method over a run (I2, equality variant). The
Lagrangian bound is only available in iterations that price all pilots, which are the
iterations that find no column.}
\label{fig:optimality}
\end{figure}


\section{Computational Time}
\label{sec:results-time}

\paragraph{Master time per iteration.} \Cref{fig:master-time} shows the time of the master
solve per column generation iteration on I2. The time of the LP master and of the two HiGHS QP
masters grows with the number of columns. The HiGHS QP solvers take 80--95\,ms per iteration
in the second half of their run, since they rebuild their internal factorization for every
master problem. The active-set method stays between $0.003$ and $1.6$\,ms per iteration with
the dense factorization and between $0.006$ and $6.5$\,ms with the sparse one, with medians
of $0.04$ and $0.07$\,ms in the second half of the run. That is about three orders of
magnitude below the HiGHS QP masters and almost three orders of magnitude below the median of
$30$\,ms of the warm-started LP master.

\begin{figure}[ht]
\centering
\begin{tikzpicture}
\begin{axis}[width=0.78\textwidth, height=0.42\textwidth, ymode=log, xlabel={CG iteration},
    ylabel={master time per iteration [ms]}, tick label style={font=\scriptsize},
    label style={font=\small}, legend style={font=\scriptsize, at={(0.98,0.5)}, anchor=east},
    every axis plot/.append style={line width=0.8pt}]
    \addplot[cred, dashdotted] table[x=it, y=mtavg, col sep=comma] {data/conv_I2_0_lp.csv};
    \addplot[cgreen] table[x=it, y=mtavg, col sep=comma] {data/conv_I2_0_explicit_qp.csv};
    \addplot[cpurple, densely dashed] table[x=it, y=mtavg, col sep=comma] {data/conv_I2_0_eliminated_qp.csv};
    \addplot[corange] table[x=it, y=mtavg, col sep=comma] {data/conv_I2_0_qpas.csv};
    \addplot[cblue] table[x=it, y=mtavg, col sep=comma] {data/conv_I2_0_qpas_dense.csv};
    \legend{LP CG, explicit QP, eliminated QP, QPAS (sparse), QPAS (dense)}
\end{axis}
\end{tikzpicture}
\caption{Master time per column generation iteration on I2 (moving average over 50
iterations).}
\label{fig:master-time}
\end{figure}

\paragraph{Where the time goes.} \Cref{fig:breakdown} splits the time of the active-set
method into its operations, for both factorizations, as shares of the total time of the run
with the sparse factorization. With the sparse factorization, the largest share, 32--48\%, is
the direction computation (the right-hand side and the sparse solve with the factor), followed
by the updates of the factorization and the periodic refactorization, together 26--55\%; the
refactorization share grows with the instance size. The dense factorization cuts the
direction computation by a factor of 2.5 to 4, the updates by a factor of 3 to 4.5, and the
refactorizations almost entirely. What remains is dominated by operations that do not touch
the factor: the multiplier test, the only operation whose cost grows with the number of
columns, is now about as expensive as the direction computation on the equality variants and
takes between a fifth and two fifths of the time of the dense runs. Pricing takes less than
7\% of the total time on all instances shown.

\begin{figure}[ht]
\centering
\begin{tikzpicture}
\begin{axis}[xbar stacked, table/col sep=comma, width=0.8\textwidth, height=0.95\textwidth, xmin=0, xmax=100,
    xlabel={share of the total time of the sparse run [\%]}, ytick=data,
    yticklabels from table={data/breakdown.csv}{label},
    y dir=reverse, bar width=6.5pt, tick label style={font=\scriptsize}, label style={font=\small},
    legend style={font=\scriptsize, at={(0.5,-0.08)}, anchor=north, legend columns=3},
    reverse legend=false]
    \addplot[fill=cblue!70, draw=cblue] table[x=direction, y expr=\coordindex, col sep=comma] {data/breakdown.csv};
    \addplot[fill=corange!70, draw=corange] table[x=multipliers, y expr=\coordindex, col sep=comma] {data/breakdown.csv};
    \addplot[fill=cgrey!50, draw=cgrey] table[x=step, y expr=\coordindex, col sep=comma] {data/breakdown.csv};
    \addplot[fill=cgreen!60, draw=cgreen] table[x=updates, y expr=\coordindex, col sep=comma] {data/breakdown.csv};
    \addplot[fill=cred!60, draw=cred] table[x=refactor, y expr=\coordindex, col sep=comma] {data/breakdown.csv};
    \addplot[fill=cpurple!50, draw=cpurple] table[x=pricing, y expr=\coordindex, col sep=comma] {data/breakdown.csv};
    \legend{direction, multipliers, step, updates, refactorization, pricing}
\end{axis}
\end{tikzpicture}
\caption{Time breakdown of the active-set method and of pricing with the sparse and the dense
factorization, as shares of the total time of the sparse run on the same instance. ``updates''
are the expansions and reductions of the working set without the refactorizations.}
\label{fig:breakdown}
\end{figure}

\Cref{tab:counters} lists the events inside the active-set method for every run. Over a full
run, the working set changes tens to hundreds of thousands of times, and every change is a
single update of the factorization. Representative swaps are rare, because representatives
carry the largest weights. Singular reduced Hessians occur, but only a handful of times per
run, and each is resolved by one zero-curvature step. The final dimension $|\Nx|$ of $\Heff$
stays below the number of coverage rows, as guaranteed by \cref{cor:dimension}. The sparse
factorization is rebuilt hundreds to thousands of times per run, almost all of them periodic
refactorizations; the dense one is rebuilt at most 16 times, each time because a singular
release had to be undone or a pivot had decayed.

\begin{table}[ht]
\centering
\scriptsize
\renewcommand{\arraystretch}{1.15}
\setlength{\tabcolsep}{3.5pt}
\resizebox{\textwidth}{!}{%
\begin{tabular}{lcrrrrrcrrrrr}
\toprule
 & & AS iter. & steps (full) & expans. & swaps & reduct. & surplus exp./red. & sing. & updates & refact. & $|\Nx|$ & $|\Ns|$ \\
\midrule
\tabinput{data/tab_qpas_counters.tex}
\bottomrule
\end{tabular}
}
\caption{Events inside the active-set method (dense factorization) over a whole column
generation run: active-set iterations, primal steps (of which full steps), expansions,
representative swaps, reductions, expansions/reductions of surpluses, singular steps,
factorization updates, refactorizations (sparse / dense factorization), and the final sizes of
$\Nx$ and $\Ns$.}
\label{tab:counters}
\end{table}

\paragraph{Scaling.} \Cref{fig:scaling} shows the total time of the active-set method with
both factorizations and of the direct LP solve as the instances grow, up to S6 with 960
coverage rows, and \cref{tab:scaling} gives the details of the two largest instances. From I2
to S6 the time of the dense variant grows almost exactly with the third power of $m$. Of this,
the number of active-set iterations accounts for a factor of about $m^2$ (from 62\,000 on I2 to
950\,000 on S6, \cref{tab:counters}) and the time per iteration for a factor of about $m$. The
latter is less than the $O(|\Nx|^2)$ of the factor operations because, at these sizes, the
parts of an iteration that grow only linearly, such as the right-hand side, the step and the
multiplier test, still take a large share of its time. The sparse variant grows faster, about
with $m^{3.5}$, because its refactorizations take an increasing share of the time, and the
direct LP solve grows faster still. With the dense
factorization, the complete column generation run is faster than the direct solution of the
full LP from I2 on in the equality variant and on all instances in the inequality variant, by
a factor of three to four on I4 and S5. On S6, and on the inequality variant of S5, the direct
LP solve did not finish within 300 seconds, while the active-set method with the dense
factorization solved both variants of S6 in 84--95 seconds. The inequality variant is
consistently faster than the equality variant, since it needs fewer active-set iterations
(\cref{sec:results-ineq}).

\begin{figure}[ht]
\centering
\begin{tikzpicture}
\begin{axis}[width=0.72\textwidth, height=0.5\textwidth, xmode=log, ymode=log,
    xlabel={coverage rows $m$}, ylabel={total time [s]}, tick label style={font=\scriptsize},
    label style={font=\small}, legend style={font=\scriptsize, at={(0.02,0.98)}, anchor=north west},
    every axis plot/.append style={line width=0.8pt}, log basis x=10, unbounded coords=jump]
    \addplot[corange, mark=triangle*, mark size=1.8pt] table[x=m, y=qpas, col sep=comma] {data/scaling_eq.csv};
    \addplot[corange, dashed, mark=triangle, mark size=1.8pt] table[x=m, y=qpas, col sep=comma] {data/scaling_ineq.csv};
    \addplot[cblue, mark=*, mark size=1.5pt] table[x=m, y=dense, col sep=comma] {data/scaling_eq.csv};
    \addplot[cblue, dashed, mark=o, mark size=1.5pt] table[x=m, y=dense, col sep=comma] {data/scaling_ineq.csv};
    \addplot[black!75, mark=square*, mark size=1.5pt] table[x=m, y=exact, col sep=comma] {data/scaling_eq.csv};
    \addplot[black!75, dashed, mark=square, mark size=1.5pt] table[x=m, y=exact, col sep=comma] {data/scaling_ineq.csv};
    \addplot[cgrey, dotted, domain=40:1000, samples=2] {1.1e-7*x^3};
    \legend{QPAS sparse (equality), QPAS sparse (inequality), QPAS dense (equality),
            QPAS dense (inequality), exact LP (equality), exact LP (inequality), $\propto m^3$}
\end{axis}
\end{tikzpicture}
\caption{Total time against the number of coverage rows for instances I1--I4, S5 and S6. The
exact LP of S6 and of the inequality variant of S5 did not finish within 300 seconds.}
\label{fig:scaling}
\end{figure}

\begin{table}[ht]
\centering
\scriptsize
\renewcommand{\arraystretch}{1.15}
\setlength{\tabcolsep}{3.5pt}
\resizebox{\textwidth}{!}{%
\begin{tabular}{lccccrrrrrrcc}
\toprule
 & ineq. & exact LP [s] & factor & status & CG iter. & columns & time [s] & master [ms/iter.] & refact. & $|\Nx|$ & bound gap & infeasibility \\
\midrule
\tabinput{data/tab_scaling.tex}
\bottomrule
\end{tabular}
}
\caption{The large instances S5 and S6. The bound gap is only available where the direct LP
solve finished within its limit of 300 seconds.}
\label{tab:scaling}
\end{table}


\section{Inequality Rows}
\label{sec:results-ineq}

With 30\% inequality rows, column generation needs fewer active-set iterations on all
instances, and fewer column generation iterations on I1--I4, considerably fewer on I1--I3
(\cref{tab:comparison-eq,tab:comparison-ineq,tab:counters}): over-covering a flight is allowed, so the
optimal schedules are found with fewer columns, and the proximal phase is shorter.

\medskip

\Cref{fig:sets} follows the sizes of the nullspace sets over a run with inequality rows.
Within the first iterations, about 35 surpluses are released: the initial schedules
over-cover their rows, so the duals of these rows become negative. These rows drop out of
$\Aeff$, and by \cref{cor:dimension} they also reduce the room for nullspace columns,
$|\Nx|\le m - |\Ns|$. Over the rest of the run $|\Ns|$ only changes slowly, between 37 and 51,
ending at 48, while individual surpluses are released and fixed again thousands of times
(\cref{tab:counters}). The number of nullspace columns grows during the column phase and ends
at $|\Nx| = 306$, just below its bound $m - |\Ns| = 308$: at the end of the run, the effective
reduced Hessian is almost as large as it can be.

\begin{figure}[ht]
\centering
\begin{tikzpicture}
\begin{axis}[width=0.78\textwidth, height=0.42\textwidth, xlabel={CG iteration},
    ylabel={set size}, tick label style={font=\scriptsize}, label style={font=\small},
    legend style={font=\scriptsize, at={(0.98,0.5)}, anchor=east},
    every axis plot/.append style={line width=0.8pt}, ymin=0]
    \addplot[cgrey, dashed] table[x=it, y=m, col sep=comma] {data/sets_I3_30.csv};
    \addplot[cred] table[x=it, y=bound, col sep=comma] {data/sets_I3_30.csv};
    \addplot[cblue] table[x=it, y=N, col sep=comma] {data/sets_I3_30.csv};
    \addplot[corange] table[x=it, y=Ns, col sep=comma] {data/sets_I3_30.csv};
    \legend{$m$, $m - |\Ns|$ (bound on $|\Nx|$), $|\Nx|$, $|\Ns|$}
\end{axis}
\end{tikzpicture}
\caption{Nullspace columns and free surpluses over a run (I3, inequality variant, dense
factorization), sampled every 20 iterations.}
\label{fig:sets}
\end{figure}

\medskip

Surplus operations are frequent (\cref{tab:counters}): with inequality rows, 14--19\% of all
expansions and reductions of the working set on I2--S6 are expansions and reductions of
surpluses, each of which is a rank-one update or downdate of the factorization. The average
time per master problem of the inequality variant is of the same order as that of the equality
variant (\cref{tab:comparison-eq,tab:comparison-ineq}). The remaining differences come from the
mix of iterations. On I2, for instance, the inequality variant has relatively few of the cheap
center-only iterations (\cref{tab:phases}), which raises its average. Singular reduced
Hessians, of columns or surpluses, occurred at most six times per run on I1--I4. The singular
surplus release of \cref{thm:zc-surplus} is exercised by the test suite of the implementation
on constructed instances.


\section{Fill-In and Refactorization of the Sparse Factor}
\label{sec:results-fill}

The factorization of $\Heff$ is updated tens of thousands of times per run
(\cref{tab:counters}). \Cref{fig:fill} shows what this does to the contents of the sparse
factor, using the quantities of \cref{def:stored-effective}, without refactorization, with
refactorization every 150 updates, and, for comparison, for the dense factor of
\cref{chap:dense}.

\begin{figure}[ht]
\centering
\begin{tikzpicture}
\begin{groupplot}[group style={group size=3 by 1, horizontal sep=0.5cm, yticklabels at=edge left},
    width=0.38\textwidth, height=0.4\textwidth, xlabel={CG iteration}, ymode=log,
    ymin=100, ymax=2e6, tick label style={font=\scriptsize}, label style={font=\small},
    title style={font=\small}, legend style={font=\scriptsize, at={(0.98,0.02)}, anchor=south east},
    every axis plot/.append style={line width=0.7pt}]
\nextgroupplot[title={sparse, never refactorized}, ylabel={off-diagonal entries of $L$}]
    \addplot[cred] table[x=it, y=stored, col sep=comma] {data/fill_I3_tau0.csv};
    \addplot[cblue] table[x=it, y=effective, col sep=comma] {data/fill_I3_tau0.csv};
    \addplot[cgrey, dashed] table[x=it, y=dense, col sep=comma] {data/fill_I3_tau0.csv};
    \legend{stored, nonzero, dense equiv.}
\nextgroupplot[title={sparse, $\tau = 150$}]
    \addplot[cred] table[x=it, y=stored, col sep=comma] {data/fill_I3_tau150.csv};
    \addplot[cblue] table[x=it, y=effective, col sep=comma] {data/fill_I3_tau150.csv};
    \addplot[cgrey, dashed] table[x=it, y=dense, col sep=comma] {data/fill_I3_tau150.csv};
\nextgroupplot[title={dense factor}]
    \addplot[cred] table[x=it, y=stored, col sep=comma] {data/fill_I3_dense_tau0.csv};
    \addplot[cblue] table[x=it, y=effective, col sep=comma] {data/fill_I3_dense_tau0.csv};
    \addplot[cgrey, dashed] table[x=it, y=dense, col sep=comma] {data/fill_I3_dense_tau0.csv};
\end{groupplot}
\end{tikzpicture}
\caption{Contents of the factor $L$ of $\Heff$ over a run (I3, equality variant), sampled
every 25 iterations: stored off-diagonal entries, the nonzero ones among them, and the
off-diagonal entries of a dense factor of the same dimension. In the right panel the three
curves coincide.}
\label{fig:fill}
\end{figure}

\paragraph{The true factor is dense.} In all panels the nonzero entries coincide with the
dense equivalent: the factor of $\Heff$ has essentially no exploitable sparsity. This confirms
\cref{rem:natural-ordering} and the starting point of \cref{chap:dense}: a fill-reducing
ordering cannot help, and the cost of a solve grows quadratically with $|\Nx|$.

\paragraph{Structural zeros accumulate.} Without the periodic refactorization, the stored
pattern of the sparse factor only grows. On I3 it reaches almost a million entries in the
column phase, about thirty times the size of the actual factor, and almost all of them are
structural zeros. Every solve and every update pays for them, and they connect active indices
with inactive ones, which enlarges the reach of the sparse solves. The drop around iteration
1\,400 is a refactorization triggered by a degraded pivot (\cref{sec:detection}). It resets the
pattern, after which it grows again and stays at 1.2 to 2.6 times the size of the factor. With
refactorization every 150 updates, the stored pattern stays within a small margin of the
effective one for most of the run. The dense factor stores exactly the dense equivalent at
every moment (\cref{prop:no-zeros}).

\paragraph{The refactorization interval.} \Cref{tab:refactor} and \cref{fig:refactor} vary the
refactorization interval of the sparse factorization. Without refactorization, the method is
three to five times slower, and the sparse factor needs six to eleven times more memory.
Refactorizing too often is also expensive, since every rebuild costs $O(|\Nx|^3)$
(\cref{prop:refactor-cost}): with an interval of 25, I3 needs almost 3\,000 rebuilds and is
70\% slower than with the default. Between these extremes the optimum is flat: any interval
between 100 and 600 is within about 10\% of the best time, and the default of 150 is within
3\% of the best on both instances. For the dense factorization, which has no structural zeros,
a periodic refactorization only costs time: with $\tau = 150$ it is 5--10\% slower than without.
Both dense variants are faster than the best sparse one by a factor of 2.1 to 2.6.

\begin{table}[ht]
\centering
\small
\renewcommand{\arraystretch}{1.15}
\resizebox{\textwidth}{!}{%
\begin{tabular}{lllrrrrrrr}
\toprule
 & factor & interval $\tau$ & time [s] & refact. & max.\ stored & max.\ zeros & mean zero share & reach$/|\Nx|$ & peak memory [MB] \\
\midrule
\tabinput{data/tab_refactor.tex}
\bottomrule
\end{tabular}
}
\caption{Effect of the refactorization interval (equality variants). ``max.\ zeros'' is the
largest number of stored entries that are zero; for the dense factor these are entries of
$L$ that happen to vanish, not structural zeros. ``reach$/|\Nx|$'' is the average number of
rows touched by a solve per nullspace column. ``never'' still refactorizes when a pivot
degrades. These runs form a separate series (times are the faster of two runs), so the times
differ slightly from \cref{tab:comparison-eq}.}
\label{tab:refactor}
\end{table}

\begin{figure}[ht]
\centering
\begin{tikzpicture}
\begin{axis}[width=0.68\textwidth, height=0.42\textwidth, xmode=log, xlabel={refactorization interval $\tau$},
    ylabel={total time [s]}, xtick={25,50,100,150,300,600,2000},
    xticklabels={25,50,100,150,300,600,never}, tick label style={font=\scriptsize},
    label style={font=\small}, legend style={font=\scriptsize, at={(0.02,0.98)}, anchor=north west},
    ymin=0, unbounded coords=jump, every axis plot/.append style={line width=0.8pt}]
    \addplot[cblue, mark=*, mark size=1.5pt] table[x=tau, y=time, col sep=comma] {data/refactor_I2.csv};
    \addplot[corange, mark=square*, mark size=1.5pt] table[x=tau, y=time, col sep=comma] {data/refactor_I3.csv};
    \addplot[cblue, dashed, domain=20:2500] table[x=tau, y=time, col sep=comma] {data/refactor_I2_dense.csv};
    \addplot[corange, dashed] table[x=tau, y=time, col sep=comma] {data/refactor_I3_dense.csv};
    \legend{I2 sparse, I3 sparse, I2 dense, I3 dense}
\end{axis}
\end{tikzpicture}
\caption{Total time of the active-set method against the refactorization interval of the
sparse factorization. The dashed lines show the dense factorization without periodic
refactorization.}
\label{fig:refactor}
\end{figure}


\section{The Sparse and the Dense Factorization}
\label{sec:results-factorizations}

\Cref{tab:backends} compares the two factorizations on all instances and variants. Both
variants of the active-set method reach the same objective and the same accuracy
(\cref{sec:results-correctness}), and their iteration counts are close, so the time
difference is caused by the linear algebra.

\begin{table}[ht]
\centering
\scriptsize
\renewcommand{\arraystretch}{1.15}
\setlength{\tabcolsep}{4pt}
\resizebox{\textwidth}{!}{%
\begin{tabular}{lcrrrrrrrrrrc}
\toprule
 & & & CG iter. & \multicolumn{3}{c}{master time [s]} & \multicolumn{2}{c}{total time [s]} & refact. & \multicolumn{2}{c}{factor memory [MB]} & objective \\
\cmidrule(lr){5-7}\cmidrule(lr){8-9}\cmidrule(lr){11-12}
instance & ineq. & $m$ & sparse / dense & sparse & dense & speedup & sparse & dense & sparse / dense & sparse & dense & rel.\ difference \\
\midrule
\tabinput{data/tab_backends.tex}
\bottomrule
\end{tabular}
}
\caption{The two factorizations of the effective reduced Hessian. Master times are summed over
the run, the speedup is their ratio, and the factor memory is the peak of the memory held by
the factorization.}
\label{tab:backends}
\end{table}

\paragraph{Speed.} The dense factorization makes the master problems 2.3 to 4.2 times
faster, and the speedup grows with the number of linking rows: 2.4 to 2.8 at $m = 225$ and $356$,
3.3 at $m = 480$, 3.6 at $m = 668$ and 4.2 at $m = 960$ (equality variants). Two effects
contribute. First, every operation on the factor is cheaper, as shown below. Second, the sparse
factorization has to be rebuilt periodically, and since a rebuild costs $O(|\Nx|^3)$ while the
interval stays fixed, these rebuilds take a growing share of its time
(\cref{fig:breakdown}): almost 40\% on S6, with more than 4\,000 rebuilds. The dense
factorization was rebuilt at most 16 times per run.

\paragraph{Cost per operation.} \Cref{tab:ops} lists the average time of the individual
factor operations as measured inside the runs. On the instances of the study, a solve and a
row addition are about five times faster with the dense factorization, and a row deletion five
to seven times. Assembling
a new row of $\Heff$ from the column pool, which is the same code for both factorizations,
now costs more than adding it to the factor.

\begin{table}[ht]
\centering
\small
\renewcommand{\arraystretch}{1.15}
\begin{tabular}{lrrrrrrrr}
\toprule
 & \multicolumn{2}{c}{solve} & \multicolumn{2}{c}{row addition} & \multicolumn{2}{c}{row deletion} & \multicolumn{2}{c}{assembly of a row} \\
\cmidrule(lr){2-3}\cmidrule(lr){4-5}\cmidrule(lr){6-7}\cmidrule(lr){8-9}
instance & sparse & dense & sparse & dense & sparse & dense & sparse & dense \\
\midrule
\tabinput{data/tab_ops.tex}
\bottomrule
\end{tabular}
\caption{Average time per factor operation in $\mu$s over a column generation run (equality
variants). Row additions include those of the refactorizations.}
\label{tab:ops}
\end{table}

To separate the cost of the operations from the behaviour of the active-set method,
\cref{fig:micro} and \cref{tab:micro} measure them in isolation. The matrix is the Gram matrix
of $k+1$ random sparse columns with entries in $\{-1,0,1\}$, which is dense like $\Heff$, both
factorizations hold its first $k$ indices, and every operation is repeated and the median time
is reported.

\begin{figure}[ht]
\centering
\begin{tikzpicture}
\begin{groupplot}[group style={group size=2 by 1, horizontal sep=1.4cm},
    width=0.5\textwidth, height=0.42\textwidth, xmode=log, ymode=log, xlabel={active size $k$},
    tick label style={font=\scriptsize}, label style={font=\small}, title style={font=\small},
    legend style={font=\scriptsize, at={(0.02,0.98)}, anchor=north west},
    every axis plot/.append style={line width=0.8pt}]
\nextgroupplot[title={solve and row addition}, ylabel={time per operation [$\mu$s]}]
    \addplot[corange, mark=*, mark size=1.2pt] table[x=k, y=sparse, col sep=comma] {data/micro_solve.csv};
    \addplot[cblue, mark=*, mark size=1.2pt] table[x=k, y=dense, col sep=comma] {data/micro_solve.csv};
    \addplot[corange, dashed, mark=square*, mark size=1.2pt] table[x=k, y=sparse, col sep=comma] {data/micro_append.csv};
    \addplot[cblue, dashed, mark=square*, mark size=1.2pt] table[x=k, y=dense, col sep=comma] {data/micro_append.csv};
    \legend{solve (sparse), solve (dense), append (sparse), append (dense)}
\nextgroupplot[title={deletion and rank-one update}]
    \addplot[corange, mark=*, mark size=1.2pt] table[x=k, y=sparse, col sep=comma] {data/micro_delete_first.csv};
    \addplot[cblue, mark=*, mark size=1.2pt] table[x=k, y=dense, col sep=comma] {data/micro_delete_first.csv};
    \addplot[corange, dashed, mark=square*, mark size=1.2pt] table[x=k, y=sparse, col sep=comma] {data/micro_update.csv};
    \addplot[cblue, dashed, mark=square*, mark size=1.2pt] table[x=k, y=dense, col sep=comma] {data/micro_update.csv};
    \legend{delete oldest (sparse), delete oldest (dense), update (sparse), update (dense)}
\end{groupplot}
\end{tikzpicture}
\caption{Median time of the factor operations against the number $k$ of active indices, for a
dense Gram matrix. The sparse factor has the dimension $\nmax = 20\,000$.}
\label{fig:micro}
\end{figure}

\begin{table}[ht]
\centering
\small
\renewcommand{\arraystretch}{1.15}
\resizebox{\textwidth}{!}{%
\begin{tabular}{lrrrrrrrr}
\toprule
 & \multicolumn{2}{c}{$k = 100$} & \multicolumn{2}{c}{$k = 350$} & \multicolumn{2}{c}{$k = 1000$} & \multicolumn{2}{c}{$k = 2000$} \\
\cmidrule(lr){2-3}\cmidrule(lr){4-5}\cmidrule(lr){6-7}\cmidrule(lr){8-9}
operation & sparse & dense & sparse & dense & sparse & dense & sparse & dense \\
\midrule
\tabinput{data/tab_micro.tex}
\bottomrule
\end{tabular}
}
\caption{Median time per operation in $\mu$s, and the memory of the factor, against the number
$k$ of active indices.}
\label{tab:micro}
\end{table}

The dense factorization is faster for every operation and every size. The largest gains are
for the operations that the sparse factorization performs through index structures: solves
and appends are 2 to 10 times faster, refactorizations 6 to 12 times. Deletions and rank-one
modifications, which both factorizations perform by the same column-oriented rank-one
recursion, are 1.2 to 2.6 times faster. For the sparse factorization the time per operation
grows faster than $k^2$ beyond $k\approx1000$, the dense one stays close to $k^2$. The memory
of the dense factor is seven to nine times smaller, since it neither stores indices nor embeds
the factor in $\nmax$ dimensions.

\paragraph{Deletions.} By \cref{prop:delete-cost}, the cost of a deletion depends on the
position of the removed index. Over a run, the trailing part $(t-\ell)/t$ that a deletion
updates was on average 0.22 to 0.36 of the factor on all instances: as anticipated in
\cref{sec:dense-delete}, most removed indices had been active only for a short time. The
dense factor never needed more than $m+1$ positions, as guaranteed by~\eqref{eq:dense-size} and
the one index of an undone singular release (\cref{sec:why-dense}).

\paragraph{Accuracy.} \Cref{fig:residual} shows the residual
$\max|(LDL^\top)_{\Nx\Nx} - \Heff|/\max(1,\max|\Heff|)$ of the factor over a run, with $\Heff$
recomputed from the columns at every sample. For both factorizations it stays at the level of
the rounding error, below $3\cdot10^{-15}$ on all runs of I1--I4, and it does not grow over
the run. The dense factor is rebuilt only a handful of times per run and yet does not drift,
even over almost 150\,000 updates on I4: the rank-one modifications of \cref{alg:c1} are
stable in practice, as \cref{rem:c1-stability} suggests.

\begin{figure}[ht]
\centering
\begin{tikzpicture}
\begin{groupplot}[group style={group size=2 by 1, horizontal sep=1.4cm, yticklabels at=edge left},
    width=0.5\textwidth, height=0.38\textwidth, ymode=log, ymin=1e-17, ymax=1e-13,
    xlabel={CG iteration}, tick label style={font=\scriptsize}, label style={font=\small},
    title style={font=\small}, legend style={font=\scriptsize, at={(0.98,0.98)}, anchor=north east},
    every axis plot/.append style={line width=0.7pt}]
\nextgroupplot[title={I3, equality variant}, ylabel={relative residual}]
    \addplot[corange] table[x=it, y=residual, col sep=comma] {data/resid_I3_0_sparse.csv};
    \addplot[cblue] table[x=it, y=residual, col sep=comma] {data/resid_I3_0_dense.csv};
    \legend{sparse, dense}
\nextgroupplot[title={I4, inequality variant}]
    \addplot[corange] table[x=it, y=residual, col sep=comma] {data/resid_I4_30_sparse.csv};
    \addplot[cblue] table[x=it, y=residual, col sep=comma] {data/resid_I4_30_dense.csv};
\end{groupplot}
\end{tikzpicture}
\caption{Relative residual of the factorization over a run, sampled every 20 iterations.}
\label{fig:residual}
\end{figure}


\section{Discussion}
\label{sec:discussion}

The experiments support the main claims of this work.
\begin{itemize}
    \item \textbf{Correctness.} The active-set method solves the same sequence of master
          problems as a general-purpose QP solver, reproduces the optimum of the full linear
          program, and certifies it by the Lagrangian bound, for equality and inequality rows
          and with both factorizations.
    \item \textbf{Speed.} With the dense factorization, a master problem of the active-set
          method costs between $0.003$ and $13$\,ms on average on our instances, and $0.08$\,ms
          on I2. That is almost three orders of magnitude less than a HiGHS QP master and two
          to three orders of magnitude less than a warm-started LP master. This makes regularized column
          generation much faster than classical column generation, although it often needs
          more iterations: on I3, 3.0\,s against 663\,s. The complete run is faster than the
          direct solution of the full LP from I2 on, and on the largest instances the direct
          solve does not finish within the time in which the active-set method is done.
    \item \textbf{Inequality rows} add no new kind of linear algebra: the surplus operations
          are rank-one modifications of the same factorization, free surpluses lower the bound
          $m - |\Ns|$ on the size of $\Heff$, the time per master problem stays of the same
          order, and the runs are shorter.
    \item \textbf{The factorization.} The factor of $\Heff$ is dense. A sparse factorization
          needs periodic refactorization to control its structural zeros, and a fixed interval
          of a few hundred updates works well. A dense factorization in activation order
          needs none, is two to four times faster with a speedup that grows with $m$, needs an
          order of magnitude less memory, and is just as accurate.
\end{itemize}

The experiments also show the limits of the approach. Since $\Heff$ is dense, the cost of the
method grows with the cube of the number of linking rows, and the \textbf{proximal phase} can
be long, since the duals approach the restricted optimum in many small steps
(\cref{prop:center-only}). With the dense factorization the factor is no longer the dominant
cost: the multiplier test, which touches every column of the working set, now costs as much as
the direction computation. Natural next steps are therefore an incremental multiplier test that
avoids rescanning all columns, and an adaptive stabilization parameter in the proximal phase.
A change of $\delta$ does not affect the factorization (\cref{rem:factor-independent}) and is
therefore free for this method.

\chapter{Conclusion}
\label{chap:conclusion}

This work developed a hot-started nullspace active-set method for the quadratic master
problems of regularized column generation, for multi-block linear programs whose linking
constraints may be equalities or inequalities of either direction.

\section{Summary}

The starting point was a \textbf{unified formulation} (\cref{chap:formulation}): surplus
variables bring every linking row into equality form, and the sign conditions of the
inequality rows become ordinary dual constraints of the surplus columns. Proximal
regularization of the dual then leads to a convex quadratic master problem
(\cref{chap:regularization}). Because all linking rows are equalities, its auxiliary variable
can be eliminated for all rows, and the result has only the convexity constraints and
nonnegativity bounds. Two identities underlie everything that follows: the duals of the
linking rows are an explicit affine function of the primal iterate, and the gradient of the
quadratic objective consists of the reduced costs of the columns and the duals of the
inequality rows.

\medskip

For the active-set method (\cref{chap:qpas}), the constraint matrix of a working set admits a
\textbf{reordering into five sets}: representatives, working-set columns, nullspace columns,
active surpluses and free surpluses. The surpluses violate the column generation property of
the convexity matrix, but this only adds a group with identity structure (active surpluses)
and a group of empty columns (free surpluses). The
reordering yields an explicit nullspace basis of vectors $e_j - e_{\rep(j)}$ and unit vectors,
which is never stored (\cref{chap:nullspace}). The reduced KKT system \textbf{decouples}: the
surplus part of the direction is explicit, and only the \textbf{effective reduced Hessian}
$\Heff = (\Aeff Z_x)^\top(\Aeff Z_x)$ has to be factorized. Here $\Aeff$ holds the equality rows
and the inequality rows with active surplus. The column part of the method is thus the method
for equality rows, applied to a linking matrix whose rows change with the surplus working set.
The Lagrange multipliers are the reduced costs and the dual signs, so the optimality test of
the master problem and the pricing test are the same test.

\medskip

Every change of the working set is a row addition, a row deletion or a rank-one modification
of an $LDL^\top$ factorization of $\Heff$ (\cref{chap:updates}). A new column or a new
stability center leaves the factorization unchanged, which gives a hot start of every master
problem. Singular reduced Hessians are resolved by an exact zero-curvature step, which removes
the dependency and keeps every working-set variable except the released one at its bound, for
released columns and released surpluses alike (\cref{chap:singular}).

\medskip

The factorization was realized in two ways. A sparse factorization of fixed dimension
$\nmax$, updated with CHOLMOD, accumulates structural zeros, which periodic refactorization in
increasing index order controls (\cref{chap:updates}). Since $\Heff$ is small and dense, it can
instead be factorized densely in the order in which its indices became active
(\cref{chap:dense}). A released column is then appended, a removed index is deleted by a
rank-one update of the positions behind it, and the surplus operations are rank-one updates
and downdates by the method of Gill, Golub, Murray and Saunders. This factor never contains a
stale entry and needs no periodic refactorization; its new pivot is exactly the Schur
complement of a released column, its downdates fail exactly when the downdated matrix is
singular, and removals never decrease a pivot.

\medskip

On generated crew scheduling instances (\cref{chap:results}), the method reproduces the optimum
of the full linear program and certifies it by the Lagrangian bound, with both factorizations.
With the dense factorization, its master problems cost between $0.003$ and $13$\,ms on
average, almost three orders of magnitude less than a general-purpose QP solver, and fast
enough to make regularized column generation much faster than classical column generation with
a warm-started LP master: 3.0\,s instead of 663\,s on the largest instance on which both were
run. The dense factorization is two to four times faster than the sparse one, with a speedup
that grows with the number of linking rows, needs an order of magnitude less memory, and is
just as accurate. The complete column generation run is faster than solving the full linear
program directly from the second instance on, by a factor of three to four on the larger
instances, and on the largest ones the direct solve did not finish in the time in which the
active-set method was done. Inequality rows need no new kind of linear algebra, keep the time
per master problem of the same order, and lead to shorter runs.


\section{Outlook}

The experiments point to four directions for further work.
\begin{itemize}
    \item \textbf{A cheaper multiplier test.} With the dense factorization, the linear algebra
          is no longer the dominant cost. The multiplier test, which computes the reduced cost of
          every working-set column at every subspace minimizer, now costs as much as the
          direction computation. Since a step changes $\pi$ only by a multiple of $\Abar p$, the
          reduced costs could be updated incrementally, or the test could be restricted to a
          candidate list as in partial pricing.
    \item \textbf{Adaptive stabilization in the proximal phase.} The long tail of center-only
          iterations is a property of the proximal point method with a fixed $\delta$
          (\cref{prop:center-only}). Since $\delta$ does not enter the factorization, it can be
          decreased in this phase at no cost to the linear algebra.
    \item \textbf{Range rows.} Linking rows with two finite bounds need surpluses with upper
          bounds, which fit into the same framework as an additional working-set state.
    \item \textbf{Inexact master solves.} Since the multipliers are the reduced costs, pricing
          could be called at any subspace minimizer instead of only at the optimum of the
          master problem. This would turn column generation into a single active-set method
          over all columns.
\end{itemize}


\backmatter
\begin{thebibliography}{99}
\addcontentsline{toc}{chapter}{Bibliography}

\bibitem{bezanson2017julia}
J.~Bezanson, A.~Edelman, S.~Karpinski, and V.~B.~Shah.
\newblock Julia: A fresh approach to numerical computing.
\newblock \emph{SIAM Review}, 59(1):65--98, 2017.

\bibitem{briant2008comparison}
O.~Briant, C.~Lemar\'echal, Ph.~Meurdesoif, S.~Michel, N.~Perrot, and F.~Vanderbeck.
\newblock Comparison of bundle and classical column generation.
\newblock \emph{Mathematical Programming}, 113(2):299--344, 2008.

\bibitem{chen2008cholmod}
Y.~Chen, T.~A.~Davis, W.~W.~Hager, and S.~Rajamanickam.
\newblock Algorithm 887: CHOLMOD, supernodal sparse Cholesky factorization and
  update/downdate.
\newblock \emph{ACM Transactions on Mathematical Software}, 35(3):22:1--22:14, 2008.

\bibitem{dantzig1960decomposition}
G.~B.~Dantzig and P.~Wolfe.
\newblock Decomposition principle for linear programs.
\newblock \emph{Operations Research}, 8(1):101--111, 1960.

\bibitem{davis1999modifying}
T.~A.~Davis and W.~W.~Hager.
\newblock Modifying a sparse Cholesky factorization.
\newblock \emph{SIAM Journal on Matrix Analysis and Applications}, 20(3):606--627, 1999.

\bibitem{davis2005row}
T.~A.~Davis and W.~W.~Hager.
\newblock Row modifications of a sparse Cholesky factorization.
\newblock \emph{SIAM Journal on Matrix Analysis and Applications}, 26(3):621--639, 2005.

\bibitem{desaulniers2005column}
G.~Desaulniers, J.~Desrosiers, and M.~M.~Solomon, editors.
\newblock \emph{Column Generation}.
\newblock Springer, New York, 2005.

\bibitem{dumerle1999stabilized}
O.~du~Merle, D.~Villeneuve, J.~Desrosiers, and P.~Hansen.
\newblock Stabilized column generation.
\newblock \emph{Discrete Mathematics}, 194(1--3):229--237, 1999.

\bibitem{fletcher1971quadratic}
R.~Fletcher.
\newblock A general quadratic programming algorithm.
\newblock \emph{Journal of the Institute of Mathematics and its Applications}, 7(1):76--91, 1971.

\bibitem{gill1974methods}
P.~E.~Gill, G.~H.~Golub, W.~Murray, and M.~A.~Saunders.
\newblock Methods for modifying matrix factorizations.
\newblock \emph{Mathematics of Computation}, 28(126):505--535, 1974.

\bibitem{gill1981practical}
P.~E.~Gill, W.~Murray, and M.~H.~Wright.
\newblock \emph{Practical Optimization}.
\newblock Academic Press, London, 1981.

\bibitem{golub2013matrix}
G.~H.~Golub and C.~F.~Van~Loan.
\newblock \emph{Matrix Computations}.
\newblock Johns Hopkins University Press, Baltimore, 4th edition, 2013.

\bibitem{huangfu2018parallelizing}
Q.~Huangfu and J.~A.~J.~Hall.
\newblock Parallelizing the dual revised simplex method.
\newblock \emph{Mathematical Programming Computation}, 10(1):119--142, 2018.

\bibitem{kelley1960cutting}
J.~E.~Kelley.
\newblock The cutting-plane method for solving convex programs.
\newblock \emph{Journal of the Society for Industrial and Applied Mathematics},
  8(4):703--712, 1960.

\bibitem{kiwiel1990proximity}
K.~C.~Kiwiel.
\newblock Proximity control in bundle methods for convex nondifferentiable minimization.
\newblock \emph{Mathematical Programming}, 46(1--3):105--122, 1990.

\bibitem{lubbecke2005selected}
M.~E.~L\"ubbecke and J.~Desrosiers.
\newblock Selected topics in column generation.
\newblock \emph{Operations Research}, 53(6):1007--1023, 2005.

\bibitem{nocedal2006numerical}
J.~Nocedal and S.~J.~Wright.
\newblock \emph{Numerical Optimization}.
\newblock Springer, New York, 2nd edition, 2006.

\bibitem{rockafellar1976monotone}
R.~T.~Rockafellar.
\newblock Monotone operators and the proximal point algorithm.
\newblock \emph{SIAM Journal on Control and Optimization}, 14(5):877--898, 1976.


\bibitem{schrijver1986theory}
A.~Schrijver.
\newblock \emph{Theory of Linear and Integer Programming}.
\newblock Wiley, Chichester, 1986.

\end{thebibliography}


\end{document}
